\documentclass[10pt]{article}
\usepackage[utf8]{inputenc}
\usepackage[T1]{fontenc}
\usepackage{graphicx}
\usepackage[export]{adjustbox}
\graphicspath{ {./images/} }
\usepackage{hyperref}
\hypersetup{colorlinks=true, linkcolor=blue, filecolor=magenta, urlcolor=cyan,}
\urlstyle{same}
\usepackage{amsmath}
\usepackage{amsfonts}
\usepackage{amssymb}
\usepackage[version=4]{mhchem}
\usepackage{stmaryrd}
\usepackage{bbold}
\usepackage{mathrsfs}
\usepackage{caption}
\usepackage{multirow}

%New command to display footnote whose markers will always be hidden
\let\svthefootnote\thefootnote
\newcommand\blfootnotetext[1]{%
  \let\thefootnote\relax\footnote{#1}%
  \addtocounter{footnote}{-1}%
  \let\thefootnote\svthefootnote%
}
%Overriding the \footnotetext command to hide the marker if its value is `0`
\let\svfootnotetext\footnotetext
\renewcommand\footnotetext[2][?]{%
  \if\relax#1\relax%
    \ifnum\value{footnote}=0\blfootnotetext{#2}\else\svfootnotetext{#2}\fi%
  \else%
    \if?#1\ifnum\value{footnote}=0\blfootnotetext{#2}\else\svfootnotetext{#2}\fi%
    \else\svfootnotetext[#1]{#2}\fi%
  \fi
}
\DeclareUnicodeCharacter{22C5}{\ifmmode\cdot\else{$\cdot$}\fi}
\DeclareUnicodeCharacter{0394}{\ifmmode\Delta\else{$\Delta$}\fi}
\DeclareUnicodeCharacter{22EE}{\ifmmode\vdots\else{$\vdots$}\fi}
\title{Chapter 21: Treatment Effects}

\author{}
\date{}

\begin{document}
\maketitle

\includegraphics[max width=\textwidth, alt={}]{2e17339c-7635-4ef8-aded-2191d8d853ff-0934_108_734_119_284}
\end{center}}
\subsection*{21.1 Introduction}
We now explicitly cover the problem of estimating an average treatment effect (ATE), sometimes called an average causal effect. An ATE is a special case of an average partial effect-it is an APE for a binary explanatory variable-and therefore many of the econometric models and methods that we have used in previous chapters can be applied or adapted to the problem of estimating ATEs.

Estimating ATEs has become important in the program evaluation literature, such as the evaluation of job-training programs or school voucher programs. Many of the early applications of the methods described in this chapter were to medical interventions, and some of the language (such as "treatment group" and "control group") is a carryover from those early applications. But the methods have proven to be useful in situations where experiments are clearly impractical.

The organizing principle of a modern approach to program evaluation is the counterfactual framework pioneered by Rubin (1974)-in fact, the framework has been dubbed the Rubin causal model (RCM)-and since adopted by many authors in statistics, econometrics, and many other fields, including Rosenbaum and Rubin (1983), Heckman (1992, 1997), Imbens and Angrist (1994), Angrist, Imbens, and Rubin (1996), Manski (1996), Heckman, Ichimura, and Todd (1997), and Angrist (1998). Research on estimating treatment or causal effects using Rubin's framework continues unabated. This chapter is intended to provide an introduction and a fairly detailed treatment of the most commonly used methods. Recent surveys include Heckman, Lalonde, and Smith (2000), Imbens (2004), Heckman and Vytlacil (2007a, 2007b), and Imbens and Wooldridge (2009), to name a handful.

Counterfactual thinking is not reserved for estimating average treatment effects using the RCM framework. Recall that in Chapter 9 we discussed how sensible applications of simultaneous equations models entail being able to think about each equation in isolation from other equations in the system. (We called this the autonomy requirement.) For example, a demand function is defined for each possible price, even though when we collect data the prices we observe are (usually assumed to be) equilibrium prices determined by the intersection of supply and demand. Therefore, the reasoning underlying the RCM should be familiar even if the particulars are not.

This chapter mainly focuses on binary treatments, although Section 21.6.2 briefly describes approaches that are available when the treatment takes on more than two values. It is also possible to consider continuous "treatments," which would make coverage of simultaneous equations models and explicit counterfactuals possible. This chapter does not attempt such a unification. Pearl (2000) provides a counterfactual setting that encompasses SEMs.

Most approaches to estimating average treatment effects fall into one of three approaches. The first exploits ignorability or unconfoundedness of the treatment conditional on a set of observed covariates. As we will see in Section 21.3, this approach is analogous to the proxy variable solution to the omitted variables problem that we discussed in Chapter 4. In fact, one approach to estimating treatment effects is to use linear regression with many controls: in effect, the treatment is exogenous once we control for enough observed factors. As we will see, an important benefit of ignorability of treatment is that no functional form or distributional assumptions are needed to identify the population parameters of interest (even though, as a practical matter, we may make parametric assumptions).

A second approach allows selection into treatment to depend on unobserved (and observed) factors. Traditionally, we would say that the treatment is "endogenous." In this case, we rely on the availability of instrumental variables (IVs) in order to identify and estimate average treatment effects. Sometimes standard IV estimators identify the effects of interest, but in other cases we rely on control function methods. Depending on the quantity we hope to estimate, we generally need to impose restrictions on functional forms or distributions or both. However, we will also discuss the work by Imbens and Angrist (1994) that provides a useful interpretation of IV estimation under very weak assumptions. We discuss IV approaches in Section 21.4.

This chapter also provides an introduction to regression discontinuity designs, where treatment-or the probability of treatment-is a discontinuous function of an observed forcing variable. If underlying regression functions are assumed to be smooth in the forcing variable, the discontinuity of treatment can be used to identify a local treatment effect. Section 21.5 considers both sharp and fuzzy designs.

There is much ongoing research on estimating average treatment effects. Some of these active areas are touched on in Section 21.6. For example, Sections 21.6.2 and 21.6.3 consider multivalued treatments and multiple treatments, respectively.

Section 21.6.4 gives a brief discussion of estimating treatment effects with panel data, showing how standard unobserved effects models can be obtained when one assumes unconfoundedness of the history of treatments conditional on time-constant unobserved heterogeneity (and observed covariates). An alternative approach, which assumes unconfoundedness conditional on the observed past history, is also covered.

\subsection*{21.2 A Counterfactual Setting and the Self-Selection Problem}
The modern literature on treatment effects begins with a counterfactual, where each individual (or other agent) has an outcome with and without treatment (where\\
"treatment" is interpreted very broadly). This section draws heavily on Heckman (1992, 1997), Imbens and Angrist (1994), and Angrist, Imbens, and Rubin (1996) (hereafter AIR). Let $y_{1}$ denote the outcome with treatment and $y_{0}$ the outcome without treatment. Because an individual cannot be in both states, we cannot observe both $y_{0}$ and $y_{1}$; in effect, the problem we face is one of missing data. In fact, we will show how to apply the inverse probability weighting methods from Section 19.8 to the problem of estimating average treatment effects.

It is important to see that we have made no assumptions about the distributions of $y_{0}$ and $y_{1}$. In many cases these may be roughly continuously distributed (such as salary), but often $y_{0}$ and $y_{1}$ are binary outcomes (such as a welfare participation indicator), or even corner solution outcomes (such as married women's labor supply). However, some of the assumptions we make will be less plausible for discontinuous random variables, something we discuss after introducing the assumptions.

The following discussion assumes that we have an independent, identically distributed sample from the population. This assumption rules out cases where the treatment of one unit affects another's outcome (possibly through general equilibrium effects, as in Heckman, Lochner, and Taber, 1998). The assumption that treatment of unit $i$ affects only the outcome of unit $i$ is called the stable unit treatment value assumption (SUTVA) in the treatment literature (see, for example, AIR). We are making a stronger assumption because random sampling implies SUTVA.

Let the variable $w$ be a binary treatment indicator, where $w=1$ denotes treatment and $w=0$ otherwise. The triple $\left(y_{0}, y_{1}, w\right)$ represents a random vector from the underlying population of interest. For a random draw $i$ from the population, we write $\left(y_{i 0}, y_{i 1}, w_{i}\right)$. However, as we have throughout, we state assumptions in terms of the population.

To measure the effect of treatment, we are interested in the difference in the outcomes with and without treatment, $y_{1}-y_{0}$. Because this is a random variable (that is, it is individual specific), we must be clear about what feature of its distribution we want to estimate. Several possibilities have been suggested in the literature. In Rosenbaum and Rubin (1983), the quantity of interest is the average treatment effect (ATE),\\
$\tau_{\text {ate }} \equiv \mathrm{E}\left(y_{1}-y_{0}\right)$,\\
which is the expected effect of treatment on a randomly drawn person from the population. Some have criticized this measure as not being especially relevant for policy purposes: because it averages across the entire population, it includes in the average units who would never be eligible for treatment. Heckman (1997) gives the example of a job training program, where we would not want to include millionaires\\
in computing the average effect of a job training program. This criticism is somewhat misleading, as we can-and would-exclude people from the population who would never be eligible. For example, in evaluating a job training program, we might restrict attention to people whose pretraining income is below a certain threshold; wealthy people would be excluded precisely because we have no interest in how job training affects the wealthy. In evaluating the benefits of a program such as Head Start, we could restrict the population to those who are actually eligible for the program or are likely to be eligible in the future. In evaluating the effectiveness of enterprise zones, we could restrict our analysis to block groups whose unemployment rates are above a certain threshold or whose per capita incomes are below a certain level.

A second quantity of interest, and one that has received much recent attention, is the average treatment effect on the treated, which we denote $\tau_{a t t}$ :\\
$\tau_{a t t} \equiv \mathrm{E}\left(y_{1}-y_{0} \mid w=1\right)$.

That is, $\tau_{\text {att }}$ is the mean effect for those who actually participated in the program. As we will see, in some special cases equations (21.1) and (21.2) are equivalent, but generally they differ.

Imbens and Angrist (1994) define another treatment effect, which they call a local average treatment effect (LATE). LATE has the advantage of being estimable using instrumental variables under very weak conditions. It has two potential drawbacks: (1) it measures the effect of treatment on a generally unidentifiable subpopulation; and (2) the definition of LATE depends on the particular instrumental variable that we have available. We will discuss LATE in the simplest setting in Section 21.4.3.

We can expand the definition of both treatment effects by conditioning on covariates. If $x$ is an observed covariate, the ATE conditional on $x$ is simply $\mathrm{E}\left(y_{1}-y_{0} \mid x\right)$; similarly, equation (21.2) becomes $\mathrm{E}\left(y_{1}-y_{0} \mid x, w=1\right)$. By choosing $x$ appropriately, we can define ATEs for various subsets of the population. For example, $x$ can be pretraining income or a binary variable indicating poverty status, race, or gender. Recent work by Heckman and Vytlacil (2006) and Heckman, Urzua, and Vytlacil (2006) unifies the various kinds of average treatment effects by defining the marginal treatment effect. In this chapter, we focus on estimating $\tau_{\text {ate }}, \tau_{\text {att }}$, and these effects on various subpopulations.

As noted previously, the difficulty in estimating equation (21.1) or (21.2) is that we observe only $y_{0}$ or $y_{1}$, not both, for each person. More precisely, along with $w$, the observed outcome is


\begin{equation*}
y=(1-w) y_{0}+w y_{1}=y_{0}+w\left(y_{1}-y_{0}\right) \tag{21.3}
\end{equation*}


Therefore, the question is, How can we estimate $\tau_{\text {ate }}$ or $\tau_{\text {att }}$ with a random sample on $y$ and $w$ (and usually some observed covariates)?

First, suppose that the treatment indicator $w$ is statistically independent of $\left(y_{0}, y_{1}\right)$, as would occur when treatment is randomized across agents. One implication of independence between treatment status and the potential outcomes is that $\tau_{\text {ate }}$ and $\tau_{\text {att }}$ are identical: $\mathrm{E}\left(y_{1}-y_{0} \mid w=1\right)=\mathrm{E}\left(y_{1}-y_{0}\right)$. Furthermore, estimation of $\tau_{\text {ate }}$ is simple. Using equation (21.3), we have\\
$\mathrm{E}(y \mid w=1)=\mathrm{E}\left(y_{1} \mid w=1\right)=\mathrm{E}\left(y_{1}\right)$,\\
where the last equality follows because $y_{1}$ and $w$ are independent. Similarly,\\
$\mathrm{E}(y \mid w=0)=\mathrm{E}\left(y_{0} \mid w=0\right)=\mathrm{E}\left(y_{0}\right)$.\\
It follows that\\
$\tau_{\text {ate }}=\tau_{\text {att }}=\mathrm{E}(y \mid w=1)-\mathrm{E}(y \mid w=0)$.

The right-hand side is easily estimated by a difference in sample means: the sample average of $y$ for the treated units minus the sample average of $y$ for the untreated units. Thus, randomized treatment guarantees that the difference-in-means estimator from basic statistics is unbiased, consistent, and asymptotically normal. In fact, these properties are preserved under the weaker assumption of mean independence: $\mathrm{E}\left(y_{0} \mid w\right)=\mathrm{E}\left(y_{0}\right)$ and $\mathrm{E}\left(y_{1} \mid w\right)=\mathrm{E}\left(y_{1}\right)$.

Randomization of treatment is often infeasible in program evaluation (although randomization of eligibility sometimes is feasible; more on this topic later). In most cases, individuals at least partly determine whether they receive treatment, and their decisions may be related to the benefits of or gain from treatment, $y_{1}-y_{0}$. In other words, there is self-selection into treatment.

It turns out that $\tau_{\text {att }}$ can be consistently estimated as a difference in means under the weaker assumption that $w$ is independent of $y_{0}$, without placing any restriction on the relationship between $w$ and $y_{1}$. To see this point, note that we can always write


\begin{align*}
\mathrm{E}(y \mid w=1)-\mathrm{E}(y \mid w=0) & =\mathrm{E}\left(y_{0} \mid w=1\right)-\mathrm{E}\left(y_{0} \mid w=0\right)+\mathrm{E}\left(y_{1}-y_{0} \mid w=1\right) \\
& =\left[\mathrm{E}\left(y_{0} \mid w=1\right)-\mathrm{E}\left(y_{0} \mid w=0\right)\right]+\tau_{a t t} \tag{21.5}
\end{align*}


If $y_{0}$ is mean independent of $w$, that is,\\
$\mathrm{E}\left(y_{0} \mid w\right)=\mathrm{E}\left(y_{0}\right)$,\\
then the first term in equation (21.5) disappears, and so the difference in means estimator is an unbiased estimator of $\tau_{\text {att }}$. Unfortunately, condition (21.6) is still a pretty\\
strong assumption. For example, suppose that people are randomly made eligible for a voluntary job training program. Condition (21.6) effectively implies that the participation decision is unrelated to what people would earn in the absence of the program. Nevertheless, the standard difference-in-means estimator is consistent for $\tau_{\text {att }}$ in some scenarios where it is inconsistent for $\tau_{\text {ate }}$.

A useful expression relating $\tau_{\text {att }}$ and $\tau_{\text {ate }}$ is obtained by writing $y_{0}=\mu_{0}+v_{0}$ and $y_{1}=\mu_{1}+v_{1}$, where $\mu_{g}=\mathrm{E}\left(y_{g}\right), g=0,1$. Then\\
$y_{1}-y_{0}=\left(\mu_{1}-\mu_{0}\right)+\left(v_{1}-v_{0}\right)=\tau_{\text {ate }}+\left(v_{1}-v_{0}\right)$.\\
Taking the expectation of this equation conditional on $w=1$ gives\\
$\tau_{\text {att }}=\tau_{\text {ate }}+\mathrm{E}\left(v_{1}-v_{0} \mid w=1\right)$.\\
We can think of $v_{1}-v_{0}$ as the person-specific gain from participation-that is, the deviation from the population mean-and so $\tau_{\text {att }}$ differs from $\tau_{\text {ate }}$ by the expected person-specific gain for those who participated. If $y_{1}-y_{0}$ is not mean independent of $w, \tau_{\text {att }}$ and $\tau_{\text {ate }}$ generally differ.

Fortunately, we can estimate $\tau_{\text {ate }}$ and $\tau_{\text {att }}$ under assumptions less restrictive than independence between $\left(y_{0}, y_{1}\right)$ and $w$. In most cases, we can collect data on individual characteristics and relevant pretreatment outcomes-sometimes a substantial amount of data. If, in an appropriate sense, treatment depends on the observables and not on the unobservables determining $\left(y_{0}, y_{1}\right)$, then we can estimate average treatment effects quite generally, as we show in the next section.

\subsection*{21.3 Methods Assuming Ignorability (or Unconfoundedness) of Treatment}
We adopt the framework of the previous section, and, in addition, we let $\mathbf{x}$ denote a vector of observed covariates. Therefore, the population is described by $\left(y_{0}, y_{1}, w, \mathbf{x}\right)$, and we observe $y, w$, and $\mathbf{x}$, where $y$ is given by equation (21.3). When $w$ and $\left(y_{0}, y_{1}\right)$ are allowed to be correlated, we need an assumption in order to identify treatment effects. Rosenbaum and Rubin (1983) introduced the following assumption, which they called ignorability of treatment (given observed covariates $\mathbf{x}$ ):

Assumption ATE. 1 (Ignorability): Conditional on $\mathbf{x}, w$ and $\left(y_{0}, y_{1}\right)$ are independent.

Assumption ATE. 1 has also been called unconfoundedness or simply conditional independence. For many purposes, it suffices to assume ignorability in a conditional mean independence sense:

ASSUMPTION ATE.1' (Ignorability in Mean): (a) $\mathrm{E}\left(y_{0} \mid \mathbf{x}, w\right)=\mathrm{E}\left(y_{0} \mid \mathbf{x}\right)$; and (b) $\mathrm{E}\left(y_{1} \mid \mathbf{x}, w\right)=\mathrm{E}\left(y_{1} \mid \mathbf{x}\right)$.

Naturally, Assumption ATE. 1 implies Assumption ATE.1'. In practice, Assumption ATE.1' might not afford much generality, although it does allow $\operatorname{Var}\left(y_{0} \mid \mathbf{x}, w\right)$ and $\operatorname{Var}\left(y_{1} \mid \mathbf{x}, w\right)$ to depend on $w$. The idea underlying Assumption ATE.1' is this: if we can observe enough information (contained in $\mathbf{x}$ ) that determines treatment, then $\left(y_{0}, y_{1}\right)$ might be mean independent of $w$, conditional on $\mathbf{x}$. Loosely, even though $\left(y_{0}, y_{1}\right)$ and $w$ might be correlated, they are uncorrelated once we partial out $\mathbf{x}$.

Assumption ATE. 1 certainly holds if $w$ is a deterministic function of $\mathbf{x}$, which has prompted some authors in econometrics to call assumptions like ATE. 1 selection on observables; see, for example, Barnow, Cain, and Goldberger (1980), Goldberger (1981), Heckman and Robb (1985), and Moffitt (1996). (We discussed a similar assumption in Section 19.9.3 in the context of missing data and attrition.) The name is fine as a label, but we must realize that Assumption ATE. 1 does allow $w$ to depend on unobservables, albeit in a restricted fashion. If $w=g(\mathbf{x}, a)$, where $a$ is an unobservable random variable independent of $\left(\mathbf{x}, y_{0}, y_{1}\right)$, then Assumption ATE. 1 holds. But $a$ cannot be arbitrarily correlated with $y_{0}$ and $y_{1}$.

To proceed, it is helpful to define the two counterfactual conditional means,


\begin{equation*}
\mu_{0}(\mathbf{x})=\mathrm{E}\left(y_{0} \mid \mathbf{x}\right), \quad \mu_{1}(\mathbf{x})=\mathrm{E}\left(y_{1} \mid \mathbf{x}\right) \tag{21.7}
\end{equation*}


In general, these functions are unknown. But ignorability-in particular, Assumption ATE.1'-along with another assumption that we will discuss shortly, are sufficient to identify $\mu_{g}(\cdot), g=0,1$. We will show this point in the next section. First, it is important to know that, under Assumption ATE.1', the average treatment effect conditional on $\mathbf{x}$ and the average treatment effect on the treated conditional on $\mathbf{x}$ are identical. More precisely, define


\begin{align*}
\tau_{a t e}(\mathbf{x}) & =\mathrm{E}\left(y_{1}-y_{0} \mid \mathbf{x}\right)=\mu_{1}(\mathbf{x})-\mu_{0}(\mathbf{x})  \tag{21.8}\\
\tau_{a t t}(\mathbf{x}) & =\mathrm{E}\left(y_{1}-y_{0} \mid \mathbf{x}, w=1\right) \tag{21.9}
\end{align*}


Then, by Assumption ATE.1', $\mathrm{E}\left(y_{g} \mid \mathbf{x}, w\right)=\mathrm{E}\left(y_{g} \mid \mathbf{x}\right), w=0,1$, and so $\tau_{\text {ate }}(\mathbf{x})= \tau_{\text {att }}(\mathbf{x})$. In general, though-even in cases where $\tau_{\text {att }}(\mathbf{x})$ is identified- $\tau_{\text {ate }}(\mathbf{x})$ and $\tau_{\text {att }}(\mathbf{x})$ can be different.

Intuitively, the ignorability assumption seems to have a better chance of holding when the set of control variables, $\mathbf{x}$, is richer. But one must be careful not to include variables in $\mathbf{x}$ that can themselves be affected by treatment. For example, suppose $w$ is a job training indicator, and $y$ is future labor earnings. We would not want to include in $\mathbf{x}$ a measure of, say, education obtained between the time of assignment\\
and the time labor earnings are measured. By doing so, we in effect hold additional education fixed when it varies in reaction to assignment $w$. Generally, including factors in $\mathbf{x}$ that are affected by $w$ causes ignorability to fail.

Mathematically, Wooldridge (2005d) has shown that ignorability fails when $\mathbf{x}$ is influenced by $w$ in the following setting. Suppose $w$ is actually randomized with respect to $\left(y_{0}, y_{1}\right)$-in which case the standard difference-of-means estimator is unbiased and consistent for $\tau_{\text {ate }}$. But at least one $\mathbf{x}$ is related to assignment in the sense that $\mathrm{D}(\mathbf{x} \mid w=1) \neq \mathrm{D}(\mathbf{x} \mid w=0)$. Using iterated expectations, it can be shown that $\mathrm{E}\left(y_{g} \mid \mathbf{x}, w\right)=\mathrm{E}\left(y_{g} \mid \mathbf{x}\right), g=0,1$ if and only if $\mathrm{E}\left(y_{g} \mid \mathbf{x}\right)=\mathrm{E}\left(y_{g}\right), g=0,1$. That is, ignorability holds only if the covariates do not help to predict the counterfactual outcomes. If $\mathbf{x}$ includes education attained between assignment and measurement of $y, \mathrm{E}\left(y_{g} \mid \mathbf{x}\right)$ almost certainly depends on $\mathbf{x}$, and so unconfoundedness fails.

Good candidates for inclusion in $\mathbf{x}$ are variables measured prior to treatment assignment, including past outcomes on $y$. As shown in Wooldridge (2009c), variables that satisfy instrumental variables assumptions-they are independent of unobservables that affect $\left(y_{0}, y_{1}\right)$ but help predict $w$-should be excluded because their inclusion increases bias in standard regression adjustment estimators unless ignorability holds without the instrument-like variables.

Unfortunately, ignorability is fundamentally untestable because we only observe $(y, w, \mathbf{x})$. In some cases, it can be tested indirectly-see, for example, the discussion in Imbens and Wooldridge (2009). An alternative is to perform a sensitivity analysis similar to studying omitted variables. Imbens and Wooldridge (2009) survey some possibilities.

Assuming that ignorability holds, what is the additional assumption we need to identify the unconditional average treatment effect, $\tau_{a t e}$ ? From the law of iterated expectations,


\begin{equation*}
\tau_{\text {ate }}=\mathrm{E}\left[\tau_{\text {ate }}(\mathbf{x})\right]=\mathrm{E}\left[\mu_{1}(\mathbf{x})-\mu_{0}(\mathbf{x})\right], \tag{21.10}
\end{equation*}


where the expectations are over the distribution of $\mathbf{x}$. As we will see in the next couple of subsections, estimating $\tau_{\text {ate }}$ will require being able to observe both control and treated units for every outcome on $\mathbf{x}$ (a weaker assumption, to be stated precisely, suffices for $\left.\tau_{a t t}\right)$. This assumption is typically called the overlap assumption.

Assumption ATE. 2 (Overlap): For all $\mathbf{x} \in \mathscr{X}$, where $\mathscr{X}$ is the support of the covariates,


\begin{equation*}
0<\mathrm{P}(w=1 \mid \mathbf{x})<1 \tag{21.11}
\end{equation*}


Overlap means that, for any setting of the covariates in the assumed population, there is a chance of seeing units in both the control and treatment groups. If, for\\
example, $\mathrm{P}\left(w=1 \mid \mathbf{x}=\mathbf{x}_{0}\right)=0$, then units having covariate values $\mathbf{x}_{0}$ will never be in the treated group. Generally, we will not be able to estimate an average treatment effect over the population that includes units with $\mathbf{x}=\mathbf{x}_{0}$.

The probability of treatment, as a function of $\mathbf{x}$, plays a very important role in estimating average treatment effects. It is usually called the propensity score, and we denote it\\
$p(\mathbf{x})=\mathrm{P}(w=1 \mid \mathbf{x}), \quad \mathbf{x} \in \mathscr{X}$.

The overlap assumption rules out the possibility that the propensity score is ever zero or one.

Rosenbaum and Rubin (1983) call igorability plus overlap strong ignorability, an assumption that is critical in all approaches to estimating $\tau_{\text {ate }}$. For completeness, we state weaker versions of the ignorability and overlap assumptions that suffice for identifying $\tau_{\text {att }}$ :\\
assumption ATT.1' (Ignorability in Mean): $\mathrm{E}\left(y_{0} \mid \mathbf{x}, w\right)=\mathrm{E}\left(y_{0} \mid \mathbf{x}\right)$.\\
assumption ATT. 2 (Overlap): For all $\mathbf{x} \in \mathscr{X}, \mathrm{P}(w=1 \mid \mathbf{x})<1$.

\subsection*{21.3.1 Identification}
Given the previous ignorability and overlap assumptions, we can establish quite generally that $\tau_{\text {ate }}$ (and, under the weaker assumptions) $\tau_{\text {att }}$ are identified. We do so in two ways, each of which motivates subsequent estimation methods.

Our first approach is based directly on the conditional mean $\mathrm{E}(y \mid \mathbf{x}, w)$. Recall that we can write $y=y_{0}+w\left(y_{1}-y_{0}\right)$. Under the mean version of ignorability,


\begin{align*}
\mathrm{E}(y \mid \mathbf{x}, w) & =\mathrm{E}\left(y_{0} \mid \mathbf{x}, w\right)+w\left[\mathrm{E}\left(y_{1} \mid \mathbf{x}, w\right)-\mathrm{E}\left(y_{0} \mid \mathbf{x}, w\right)\right] \\
& =\mathrm{E}\left(y_{0} \mid \mathbf{x}\right)+w\left[\mathrm{E}\left(y_{1} \mid \mathbf{x}\right)-\mathrm{E}\left(y_{0} \mid \mathbf{x}\right)\right] \\
& \equiv \mu_{0}(\mathbf{x})+w\left[\mu_{1}(\mathbf{x})-\mu_{0}(\mathbf{x})\right] \tag{21.13}
\end{align*}


where getting to equation (21.13) uses Assumption ATE.1'. We have shown\\
$\mathrm{E}(y \mid \mathbf{x}, w=0)=\mu_{0}(\mathbf{x}), \quad \mathrm{E}(y \mid \mathbf{x}, w=1)=\mu_{1}(\mathbf{x})$.

Because we observe $(y, \mathbf{x}, w)$, we can, under the overlap assumption, estimate $m_{0}(\mathbf{x}) \equiv \mathrm{E}(y \mid \mathbf{x}, w=0)$ and $m_{1}(\mathbf{x}) \equiv \mathrm{E}(y \mid \mathbf{x}, w=1)$ quite generally-and this claim has nothing to do with whether ignorability holds. If Assumption ATE.1' holds, then\\
$\tau_{\text {ate }}(\mathbf{x})=m_{1}(\mathbf{x})-m_{0}(\mathbf{x})$.

In other words, when we add overlap, the functions $m_{g}(\cdot), g=0,1$ that we can estimate correspond to the quantities $\mu_{g}(\cdot)$ that we need to estimate in order to estimate $\tau_{\text {ate }}(\mathbf{x})$. If we identify $m_{g}(\cdot)$ for all $\mathbf{x} \in \mathscr{X}$-and this is where overlap comes in-then we can obtain $\tau_{\text {ate }}$ as\\
$\tau_{\text {ate }}=\mathrm{E}\left[m_{1}(\mathbf{x})-m_{0}(\mathbf{x})\right]$,\\
where again the expected value is over the distribution of $\mathbf{x}$. In practice, with a random sample, we use sample averaging. Details are given in the next subsection.

If we are able to identify $\tau(\mathbf{x})$ at all $\mathbf{x} \in \mathscr{X}$, then we can also identify the average treatment effect on any subset of the population defined by $\mathbf{x}$. For example, if we define\\
$\tau_{\text {ate }, \mathscr{R}}=\mathrm{E}\left(y_{1}-y_{0} \mid \mathbf{x} \in \mathscr{R}\right)$,\\
then we have\\
$\tau_{\text {ate }, \mathscr{R}}=\mathrm{E}\left[\tau_{\text {ate }}(\mathbf{x}) \mid \mathbf{x} \in \mathscr{R}\right]$,\\
and so we can average $\tau_{\text {ate }}(\mathbf{x})$ over the subpopulation with $\mathbf{x} \in \mathscr{R}$.\\
To see that the weaker ignorability and overlap assumptions, Assumptions ATT.1' and ATT.2, suffice for identifying $\tau_{a t t}$, note that we can use $y=y_{0}+w\left(y_{1}-y_{0}\right)$ to always write


\begin{align*}
& \mathrm{E}(y \mid \mathbf{x}, w=1)-\mathrm{E}(y \mid \mathbf{x}, w=0) \\
& \quad=\mathrm{E}\left(y_{0} \mid \mathbf{x}, w=1\right)-\mathrm{E}\left(y_{0} \mid \mathbf{x}, w=0\right)+\mathrm{E}\left(y_{1}-y_{0} \mid \mathbf{x}, w=1\right) \\
& \quad=\left[\mathrm{E}\left(y_{0} \mid \mathbf{x}, w=1\right)-\mathrm{E}\left(y_{0} \mid \mathbf{x}, w=0\right)\right]+\tau_{a t t}(\mathbf{x}) \tag{21.18}
\end{align*}


If Assumption ATT.1' holds, then the term in $[\cdot]$ in equation (21.18) is zero, and the difference in estimable means, $m_{1}(\mathbf{x})-m_{0}(\mathbf{x})$, actually identifies $\tau_{a t t}(\mathbf{x})$ :\\
$\tau_{a t t}(\mathbf{x})=m_{1}(\mathbf{x})-m_{0}(\mathbf{x})$.

Assumption ATT.1' allows the gain from treatment, $y_{1}-y_{0}$, to be arbitrarily correlated with treatment, even after conditioning on $\mathbf{x}$. It requires ignorability (in mean) only with respect to $y_{0}$, the outcome in the absence of treatment. Now $\tau_{a t t}$ is obtained as\\
$\tau_{\text {att }}=\mathrm{E}\left[m_{1}(\mathbf{x})-m_{0}(\mathbf{x}) \mid w=1\right]$.

We can use this expression to see how the weaker overlap condition, Assumption ATT.2, suffices. By definition, $m_{1}(\cdot)$ is for the treated subpopulation, $w=1$, and so\\
we only need to estimate this regression function for values of $\mathbf{x}$ corresponding to units in the treated group. In other words, we do not need a positive probability of treatment for all $\mathbf{x}$. But we still must estimate $m_{0}(\cdot)$ at values of $\mathbf{x}$ corresponding to the treated subpopulation. If there are values of $\mathbf{x}$ where treatment is certain, we cannot hope to generally estimate $\mathrm{E}(y \mid \mathbf{x}, w=0)$ for such values because we will not observe control units with the same values of the covariates. The weaker overlap assumption, $\mathrm{P}(w=1 \mid \mathbf{x})<1$, rules out this possibility.

We summarize the previous discussion with a proposition.\\
proposition 21.1: Under Assumption ATE.1', equation (21.15) holds. If we add the overlap Assumption ATE.2, we can obtain $\tau_{\text {ate }}$ as in equation (21.16). If we assume only Assumption ATT.1', we have $\tau_{\text {att }}(\mathbf{x})=m_{1}(\mathbf{x})-m_{0}(\mathbf{x})$; under Assumption ATT.2, $\tau_{\text {att }}$ is identified as in equation (21.20).

A second way to establish identification is to use inverse propensity score weighting. First maintain Assumption ATE.1'. Noting that $w y=w y_{1}$, we have, by iterated expectations,

$$
\begin{aligned}
\mathrm{E}\left[\left.\frac{w y}{p(\mathbf{x})} \right\rvert\, \mathbf{x}\right] & =\mathrm{E}\left[\left.\frac{w y_{1}}{p(\mathbf{x})} \right\rvert\, \mathbf{x}\right]=\mathrm{E}\left\{\left.\mathrm{E}\left[\left.\frac{w y_{1}}{p(\mathbf{x})} \right\rvert\, \mathbf{x}, w\right] \right\rvert\, \mathbf{x}\right\}=\mathrm{E}\left\{\left.\frac{w \mathrm{E}\left(y_{1} \mid \mathbf{x}, w\right)}{p(\mathbf{x})} \right\rvert\, \mathbf{x}\right\} \\
& =\mathrm{E}\left\{\left.\frac{w \mathrm{E}\left(y_{1} \mid \mathbf{x}\right)}{p(\mathbf{x})} \right\rvert\, \mathbf{x}\right\}=\mathrm{E}\left\{\left.\frac{w}{p(\mathbf{x})} \right\rvert\, \mathbf{x}\right\} \mu_{1}(\mathbf{x})=\mu_{1}(\mathbf{x})
\end{aligned}
$$

because $\mathrm{E}(w \mid \mathbf{x})=p(\mathbf{x})$. A similar argument shows that\\
$\mathrm{E}\left[\left.\frac{(1-w) y}{[1-p(\mathbf{x})]} \right\rvert\, \mathbf{x}\right]=\mu_{0}(\mathbf{x})$.\\
Combining these two results and using simple algebra gives\\
$\mathrm{E}\left\{\left.\frac{[w-p(\mathbf{x})] y}{p(\mathbf{x})[1-p(\mathbf{x})]} \right\rvert\, \mathbf{x}\right\}=\mu_{1}(\mathbf{x})-\mu_{0}(\mathbf{x})=\tau_{\text {ate }}(\mathbf{x})$.

Of course, this expression only makes sense for $\mathbf{x}$ such that $0<p(\mathbf{x})<1$. If we maintain Assumption ATE.2, and assume that the expectation exists, we can use iterated expectations to write\\
$\tau_{\text {ate }}=\mathrm{E}\left\{\frac{[w-p(\mathbf{x})] y}{p(\mathbf{x})[1-p(\mathbf{x})]}\right\}$,\\
which, because $w, y$, and $\mathbf{x}$ are all observed, establishes identification of $\tau_{\text {ate }}$ using the propensity score rather than regression functions.

The argument for $\tau_{\text {att }}$ is a little different. Write

$$
\begin{aligned}
{[w-p(\mathbf{x})] y } & =[w-p(\mathbf{x})]\left[y_{0}+w\left(y_{1}-y_{0}\right)\right] \\
& =[w-p(\mathbf{x})] y_{0}+w[w-p(\mathbf{x})]\left(y_{1}-y_{0}\right) \\
& =[w-p(\mathbf{x})] y_{0}+w[1-p(\mathbf{x})]\left(y_{1}-y_{0}\right),
\end{aligned}
$$

where the last equality follows because $w^{2}=w$. Therefore,


\begin{equation*}
\frac{[w-p(\mathbf{x})] y}{[1-p(\mathbf{x})]}=\frac{[w-p(\mathbf{x})] y_{0}}{[1-p(\mathbf{x})]}+w\left(y_{1}-y_{0}\right) . \tag{21.23}
\end{equation*}


Consider the numerator of the first term on the right-hand side of equation (21.23):

$$
\begin{aligned}
\mathrm{E}\left\{[w-p(\mathbf{x})] y_{0} \mid \mathbf{x}\right\} & =\mathrm{E}\left(\mathrm{E}\left\{[w-p(\mathbf{x})] y_{0} \mid \mathbf{x}, w\right\} \mid \mathbf{x}\right)=\mathrm{E}\left\{[w-p(\mathbf{x})] \mathrm{E}\left(y_{0} \mid \mathbf{x}, w\right) \mid \mathbf{x}\right\} \\
& =\mathrm{E}\left\{[w-p(\mathbf{x})] \mathrm{E}\left(y_{0} \mid \mathbf{x}\right) \mid \mathbf{x}\right\}=\mathrm{E}\{[w-p(\mathbf{x})] \mid \mathbf{x}\} \mu_{0}(\mathbf{x})=0 .
\end{aligned}
$$

Therefore,


\begin{equation*}
\mathrm{E}\left\{\left.\frac{[w-p(\mathbf{x})] y}{[1-p(\mathbf{x})]} \right\rvert\, \mathbf{x}\right\}=\mathrm{E}\left[w\left(y_{1}-y_{0}\right) \mid \mathbf{x}\right] \tag{21.24}
\end{equation*}


and so the unconditional expectations are the same, too. But


\begin{align*}
\mathrm{E}\left[w\left(y_{1}-y_{0}\right)\right] & =\mathrm{P}(w=0) \mathrm{E}\left[w\left(y_{1}-y_{0}\right) \mid w=0\right]+\mathrm{P}(w=1) \mathrm{E}\left[w\left(y_{1}-y_{0}\right) \mid w=1\right] \\
& =0+\mathrm{P}(w=1) \mathrm{E}\left[w\left(y_{1}-y_{0}\right) \mid w=1\right] \\
& \equiv \rho \tau_{a t t} \tag{21.25}
\end{align*}


where $\rho \equiv \mathrm{P}(w=1)$ is the unconditional probability of treatment. Putting the pieces together gives


\begin{equation*}
\tau_{\text {att }}=\mathrm{E}\left\{\frac{[w-p(\mathbf{x})] y}{\rho[1-p(\mathbf{x})]}\right\} ; \tag{21.26}
\end{equation*}


notice that this expression only requires the weaker assumption $p(\mathbf{x})<1$ for all $\mathbf{x}$. We summarize with a proposition.\\
proposition 21.2: Under Assumptions ATE.1' and ATE.2, $\tau_{\text {ate }}$ can be expressed as in equation (21.22). Under the weaker assumptions Assumption ATT.1' and ATT.2, $\tau_{\text {att }}$ can be expressed as in equation (21.26).

Wooldridge (1999c) derived expression (21.22) in the context of random coefficient models, while equation (21.26) is essentially due to Dehejia and Wahba (1999), who\\
used the stronger Assumption ATE.1. As we will see, these identification results can be directly turned into estimating equations for $\tau_{\text {ate }}$ and $\tau_{\text {att }}$. We now turn to estimation in the next several subsections.

\subsection*{21.3.2 Regression Adjustment}
The identification strategy based on the regression functions $\mathrm{E}(y \mid \mathbf{x}, w=0)$ and $\mathrm{E}(y \mid \mathbf{x}, w=1)$ leads directly to straightforward estimation approaches. Because we have a random sample on $(y, w, \mathbf{x})$ from the relevant population, $m_{1}(\mathbf{x}) \equiv \mathrm{E}(y \mid \mathbf{x}, w=1)$ and $m_{0}(\mathbf{x}) \equiv \mathrm{E}(y \mid \mathbf{x}, w=0)$ are nonparametrically identified. That is, these are conditional expectations that depend entirely on observables, and so they can be consistently estimated quite generally. (See Härdle and Linton, 1994, or Li and Racine, 2007, for assumptions and methods.) For the purposes of identification, we can just assume $m_{1}(\mathbf{x})$ and $m_{0}(\mathbf{x})$ are known, and the fact that they are known means that $\tau(\mathbf{x})$ is identified. If $\hat{m}_{1}(\mathbf{x})$ and $\hat{m}_{0}(\mathbf{x})$ are consistent estimators (in an appropriate sense), using the random sample of size $N$, a consistent estimator of $A T E$ under fairly weak assumptions is


\begin{equation*}
\hat{\tau}_{\text {ate, reg }}=N^{-1} \sum_{i=1}^{N}\left[\hat{m}_{1}\left(\mathbf{x}_{i}\right)-\hat{m}_{0}\left(\mathbf{x}_{i}\right)\right] \tag{21.27}
\end{equation*}


while a consistent estimator of $A T E_{1}$ is


\begin{equation*}
\hat{\tau}_{\text {att, reg }}=\left(\sum_{i=1}^{N} w_{i}\right)^{-1}\left\{\sum_{i=1}^{N} w_{i}\left[\hat{m}_{1}\left(\mathbf{x}_{i}\right)-\hat{m}_{0}\left(\mathbf{x}_{i}\right)\right]\right\} \tag{21.28}
\end{equation*}


The estimators in equations (21.27) and (21.28) are called regression adjustment estimators of $\tau_{\text {ate }}$ and $\tau_{\text {att }}$, respectively. Notice that $\hat{\tau}_{\text {att, reg }}$ simply averages the differences in predicted values, $\hat{m}_{1}\left(\mathbf{x}_{i}\right)-\hat{m}_{0}\left(\mathbf{x}_{i}\right)$, over the subsample of treated units, $w_{i}=1$.

The key implementation issue in computing $\hat{\tau}_{\text {ate, reg }}$ and $\hat{\tau}_{\text {att, reg }}$ is how to obtain $\hat{m}_{0}(\cdot)$ and $\hat{m}_{1}(\cdot)$. To be as flexible as possible, one could use nonparametric estimators, such as kernel estimators or series estimators. Kernel estimators use "local averaging" or "local smoothing" to estimate a function at a particular point. (See Li and Racine, 2007, for a comprehensive treatment of kernel regression with both continuous and discrete covariates.) In addition to being flexible, such local methods have the benefit of forcing us to confront problems with overlap in the covariate distribution. Consider equation (21.27). The function $\hat{m}_{0}(\cdot)$ is obtained using only those with $w_{i}=0$, and $\hat{m}_{1}(\cdot)$ is obtained using only those with $w_{i}=1$. Thus, in obtaining the summand $\hat{m}_{1}\left(\mathbf{x}_{i}\right)-\hat{m}_{0}\left(\mathbf{x}_{i}\right)$ for, say, someone in the control group, we\\
must obtain $\hat{m}_{1}\left(\mathbf{x}_{i}\right)$, where $\hat{m}_{1}(\cdot)$ was obtained only using those in the treatment group. If $\mathbf{x}_{i}$ is very different from the covariate values of units in the treated sample, it is pretty hopeless to estimate $\hat{m}_{1}(\mathbf{x})$ at $\mathbf{x}=\mathbf{x}_{i}$ using local averaging methods. Similarly, if $i$ denotes a treated unit, we need to evaluate $\hat{m}_{0}(\cdot)$, which is obtained using control units, at covariate values for a treated unit. Generally, the overlap assumption means that in a large enough data set, we will see both control and treated units for a given set of covariate values (or, at least, for "close" covariate values). If we do not see these units, we cannot hope to estimate $\tau_{\text {ate }}$ for the original population.

Equation (21.28) makes it clear that the weaker overlap assumption, Assumption ATT, suffices for obtaining a satisfactory estimate, $\hat{\tau}_{a t t}$. Because $\hat{m}_{1}(\cdot)$ is obtained using treated units, and equation (21.28) sums only over treated units, we allow the possibility that some units in the control group will have covariates very different from the range of covariate values in the treated subset. But we still need to obtain $\hat{m}_{0}\left(\mathbf{x}_{i}\right)$ when $i$ is a treated unit. Therefore, for every treated unit, we hopefully have some units in the control group with similar values for the covariates.

Series estimation involves global approximation using flexible parametric models, where the flexibility of the model increases with the sample size. For a recent treatment, see Li and Racine (2007). Hahn (1998) shows that nonparametric regression adjustment using series estimators can achieve the asymptotic efficiency bound for estimating $\tau_{\text {ate }}$.

The problem with lack of overlap can be seen in a simpler setting. Suppose there is only one binary covariate, $x$, and Assumption ATE.1' holds; for concreteness, $x$ could be an indicator for whether pretraining earnings are below a certain threshold. Suppose that everyone in the relevant population with $x=1$ participates in the program. Then, while we can estimate $\mathrm{E}(y \mid x=1, w=1)$ with a random sample from the population, we cannot estimate $\mathrm{E}(y \mid x=1, w=0)$ because we have no data on the subpopulation with $x=1$ and $w=0$. Intuitively, we only observe the counterfactual $y_{1}$ when $x=1$; we never observe $y_{0}$ for any members of the population with $x=1$. Therefore, $\tau_{\text {ate }}(x)$ is not identified at $x=1$.

If some people with $x=0$ participate while others do not, we can estimate $\mathrm{E}(y \mid x=0, w=1)-\mathrm{E}(y \mid x=0, w=0)$ using a simple difference in averages over the group with $x=0$, and so $\tau_{\text {ate }}(x)$ is identified at $x=0$. But if we cannot estimate $\tau_{\text {ate }}(1)$, we cannot estimate the unconditional ATE because $\tau_{\text {ate }}=\mathrm{P}(x=0)$. $\tau_{\text {ate }}(0)+\mathrm{P}(x=1) \cdot \tau_{\text {ate }}(1)$. In effect, we can only estimate the ATE over the subpopulation with $x=0$, which means that we must redefine the population of interest. This limitation is unfortunate: presumably we would be very interested in the program's effects on the group that always participates.

A similar conclusion holds if the group with $x=0$ never participates in the program. Then $\tau_{\text {ate }}(0)$ is not estimable because $\mathrm{E}(y \mid x=0, w=1)$ is not estimable. If some people with $x=1$ participated while others did not, $\tau_{\text {ate }}(1)$ would be identified, and then we would view the population of interest as the subgroup with $x=1$. There is one important difference between this situation and the one where the $x=1$ group always receives treatment: it may be legitimate to exclude from the population people who have no chance of treatment based on observed covariates. This observation is related to the issue we discussed in Section 21.2 concerning the relevant population for defining $\tau_{\text {ate }}$. If, for example, people with very high preprogram earnings $(x=0)$ have no chance of participating in a job training program, then we would not want to average together $\tau_{\text {ate }}(0)$ and $\tau_{\text {ate }}(1) ; \tau_{\text {ate }}(1)$ by itself is much more interesting.

Although the previous example is extreme, its consequences can arise in more plausible settings. Suppose that $\mathbf{x}$ is a vector of binary indicators for pretraining income intervals. For most of the intervals, the probability of participating is strictly between zero and one. If the participation probability is zero at the highest income level, we simply exclude the high-income group from the relevant population. (If we specify $\tau_{\text {att }}$ as the object of interest, we also will exclude high-income people.) Unfortunately, if participation is certain at low-income levels, we must exclude those low-income groups as well.

One can compute simple statistics to judge whether overlap is a problem. As described in Imbens and Rubin (forthcoming), one can compute the normalized differences, which take the form


\begin{equation*}
\frac{\left(\bar{x}_{1 j}-\bar{x}_{0 j}\right)}{\left(s_{1 j}^{2}+s_{0 j}^{2}\right)^{1 / 2}} \tag{21.29}
\end{equation*}


where $\bar{x}_{g j}$ is the sample average of covariate $j$ for group $g=0,1$ and $s_{g j}$ is the sample standard deviation. Imbens and Rubin suggest that normalized differences above 0.25 are cause for concern. (Notice that the normalized differences are not the $t$ statistic for testing the difference in means. It is the difference in means standardized by a measure of dispersion that is important; the sample size should not play a direct role.) Unfortunately, even if the normalized differences are all small, they only focus on one feature of the marginal distributions. Overlap can still fail in more complicated ways. In the next subsection, we will discuss how to use the estimated propensity score to evaluate overlap. If one or more normalized differences are large, one might need to redefine the population of interest.

It is useful to give a complete treatment of regression adjustment in the parametric case. (It is likely that the same analysis holds when the parametric model is allowed\\
to become more flexible as the sample size grows, at least under standard restrictions.) Let $m_{0}\left(\mathbf{x}, \boldsymbol{\delta}_{0}\right)$ and $m_{1}\left(\mathbf{x}, \boldsymbol{\delta}_{1}\right)$ be parametric functions, where $m_{0}\left(\mathbf{x}, \boldsymbol{\delta}_{0}\right)$ is estimated using the $w_{i}=0$ observations and $m_{1}\left(\mathbf{x}, \boldsymbol{\delta}_{1}\right)$ is estimated using the $w_{i}=1$ observations. We have now covered several models and estimation methods that can be used. In the leading case, we can make the functions linear in parameters and use OLS, but we might want to exploit the nature of $y$; later on, we will discuss this topic further.

Given $\sqrt{N}$-consistent and asymptotically normal estimators $\hat{\boldsymbol{\delta}}_{0}$ and $\hat{\boldsymbol{\delta}}_{1}$, we have the following parametric regression adjustment estimate of $\tau_{\text {ate }}$ :


\begin{equation*}
\hat{\tau}_{a t e, r e g}=N^{-1} \sum_{i=1}^{N}\left[m_{1}\left(\mathbf{x}_{i}, \hat{\boldsymbol{\delta}}_{1}\right)-m_{0}\left(\mathbf{x}_{i}, \hat{\boldsymbol{\delta}}_{0}\right)\right], \tag{21.30}
\end{equation*}


which will be $\sqrt{N}$-consistent and asymptotically normal for $\tau_{\text {ate }}$. Using Problem 12.17, it can be shown that

$$
\begin{aligned}
\operatorname{Avar} \sqrt{N}\left(\hat{\tau}_{\text {ate, reg }}-\tau_{\text {ate }}\right)= & \mathrm{E}\left\{\left[m_{1}\left(\mathbf{x}_{i}, \boldsymbol{\delta}_{1}\right)-m_{0}\left(\mathbf{x}_{i}, \boldsymbol{\delta}_{0}\right)-\tau_{\text {ate }}\right]^{2}\right\} \\
& +\mathrm{E}\left[\nabla_{\boldsymbol{\delta}_{0}} m_{0}\left(\mathbf{x}_{i}, \boldsymbol{\delta}_{0}\right)\right] \mathbf{V}_{0} \mathrm{E}\left[\nabla_{\boldsymbol{\delta}_{0}} m_{0}\left(\mathbf{x}_{i}, \boldsymbol{\delta}_{0}\right)\right]^{\prime} \\
& +\mathrm{E}\left[\nabla_{\boldsymbol{\delta}_{1}} m_{1}\left(\mathbf{x}_{i}, \boldsymbol{\delta}_{1}\right)\right] \mathbf{V}_{1} \mathrm{E}\left[\nabla_{\boldsymbol{\delta}_{1}} m_{1}\left(\mathbf{x}_{i}, \boldsymbol{\delta}_{1}\right)\right]^{\prime}
\end{aligned}
$$

where $\mathbf{V}_{0}$ is the asymptotic variance of $\sqrt{N}\left(\hat{\boldsymbol{\delta}}_{0}-\boldsymbol{\delta}_{0}\right)$ and similarly for $\mathbf{V}_{1}$. This formula makes it clear that it is better to use more efficient estimators for $\boldsymbol{\delta}_{0}$ and $\boldsymbol{\delta}_{1}$. A variance estimator is simple:


\begin{align*}
N \cdot \operatorname{Avar}\left(\widehat{\hat{\tau}_{a t e, \text { reg }}}\right)= & N^{-1} \sum_{i=1}^{N}\left[m_{1}\left(\mathbf{x}_{i}, \hat{\boldsymbol{\delta}}_{1}\right)-m_{0}\left(\mathbf{x}_{i}, \hat{\boldsymbol{\delta}}_{0}\right)-\hat{\tau}_{\text {ate,reg }}\right]^{2} \\
& +\left[N^{-1} \sum_{i=1}^{N} \nabla_{\boldsymbol{\delta}_{0}} m_{0}\left(\mathbf{x}_{i}, \hat{\boldsymbol{\delta}}_{0}\right)\right] \hat{\mathbf{V}}_{0}\left[N^{-1} \sum_{i=1}^{N} \nabla_{\boldsymbol{\delta}_{0}} m_{0}\left(\mathbf{x}_{i}, \hat{\boldsymbol{\delta}}_{0}\right)\right]^{\prime} \\
& +\left[N^{-1} \sum_{i=1}^{N} \nabla_{\boldsymbol{\delta}_{1}} m_{1}\left(\mathbf{x}_{i}, \hat{\boldsymbol{\delta}}_{1}\right)\right] \hat{\mathbf{V}}_{1}\left[N^{-1} \sum_{i=1}^{N} \nabla_{\boldsymbol{\delta}_{1}} m_{1}\left(\mathbf{x}_{i}, \hat{\boldsymbol{\delta}}_{1}\right)\right]^{\prime} \tag{21.31}
\end{align*}


An alternative to an analytical expression is to use the bootstrap, which is straightforward in this context. One simply includes the estimation of both mean functions, the calculation of the difference $m_{1}\left(\mathbf{x}_{i}, \hat{\boldsymbol{\delta}}_{1}\right)-m_{0}\left(\mathbf{x}_{i}, \hat{\boldsymbol{\delta}}_{0}\right)$, and then the averaging of these to obtain $\hat{\tau}_{\text {ate, reg }}$ in the same resampling scheme. All three sources of estimation error-in $\hat{\boldsymbol{\delta}}_{0}$, in $\hat{\boldsymbol{\delta}}_{1}$, and in replacing the expected value over the distribution of $\mathbf{x}_{i}$\\
with the sample average-will be properly accounted for. In fact, as described in Li and Racine (2007), the bootstrap is generally valid for nonparametric procedures, too.

Given the estimated regression functions, it is easy to estimate ATEs over a subset of the population, say $\tau_{\text {ate }, \mathscr{R}}=\mathrm{E}\left(y_{1}-y_{0} \mid \mathbf{x} \in \mathscr{R}\right)$, as\\
$\hat{\tau}_{\text {ate }, \mathscr{R}, \text { reg }}=N_{\mathscr{R}}^{-1} \sum_{i=1}^{N} 1\left[\mathbf{x}_{i} \in \mathscr{R}\right] \cdot\left[m_{1}\left(\mathbf{x}_{i}, \hat{\boldsymbol{\delta}}_{1}\right)-m_{0}\left(\mathbf{x}_{i}, \hat{\boldsymbol{\delta}}_{0}\right)\right]$.\\
As mentioned earlier, linear regression is still the most popular method of regression adjustment. Suppose that $m_{0}\left(\mathbf{x}, \boldsymbol{\delta}_{0}\right)=\alpha_{0}+\mathbf{x} \boldsymbol{\beta}_{0}$ and $m_{1}\left(\mathbf{x}, \boldsymbol{\delta}_{1}\right)=\alpha_{1}+\mathbf{x} \boldsymbol{\beta}_{1}$. Then $\left(\hat{\alpha}_{0}, \hat{\boldsymbol{\beta}}_{0}\right)$ are from the regression $y_{i}$ on $1, \mathbf{x}_{i}$ with $w_{i}=0$ and similarly for $\left(\hat{\alpha}_{1}, \hat{\boldsymbol{\beta}}_{1}\right)$. Then


\begin{equation*}
\hat{\tau}_{\text {ate, reg }}(\mathbf{x})=\left(\hat{\alpha}_{1}-\hat{\alpha}_{0}\right)+\mathbf{x}\left(\hat{\boldsymbol{\beta}}_{1}-\hat{\boldsymbol{\beta}}_{0}\right) \tag{21.32}
\end{equation*}



\begin{equation*}
\hat{\tau}_{\text {ate, reg }}=\left(\hat{\alpha}_{1}-\hat{\alpha}_{0}\right)+\overline{\mathbf{x}}\left(\hat{\boldsymbol{\beta}}_{1}-\hat{\boldsymbol{\beta}}_{0}\right) \tag{21.33}
\end{equation*}


$\hat{\tau}_{a t \tau, r e g}=\left(\hat{\alpha}_{1}-\hat{\alpha}_{0}\right)+\overline{\mathbf{x}}_{1}\left(\hat{\boldsymbol{\beta}}_{1}-\hat{\boldsymbol{\beta}}_{0}\right)$,\\
where $\overline{\mathbf{x}}$ is the sample average over the entire sample and $\overline{\mathbf{x}}_{1}=N_{1}^{-1} \sum_{i=1}^{N} w_{i} \mathbf{x}_{i}$ is the average over the treated subsample. The estimate of $\tau_{\text {ate }, \mathscr{R}}$ is simply $\left(\hat{\alpha}_{1}-\hat{\alpha}_{0}\right)+ \overline{\mathbf{x}}_{\mathscr{R}}\left(\hat{\boldsymbol{\beta}}_{1}-\hat{\boldsymbol{\beta}}_{0}\right)$, where $\overline{\mathbf{x}}_{\mathscr{R}}$ is the sample average over the restricted sample.

Replacing the vector $\mathbf{x}_{i}$ with any functions $\mathbf{h}\left(\mathbf{x}_{i}\right)$ is trivial. The only minor point is that we should demean the regressors $\mathbf{h}_{i}=\mathbf{h}\left(\mathbf{x}_{i}\right)$ using $\overline{\mathbf{h}}=N^{-1} \sum_{i=1}^{N} \mathbf{h}_{i}$ (or over a subset); it makes no sense to use $\mathbf{h}(\overline{\mathbf{x}})$.

If we ignore the sampling variance in the sample averages, we can obtain a standard error for $\hat{\tau}_{\text {ate }, \mathscr{R}, \text { reg }}$ by using a pooled regression. Using the entire sample, run the regression\\
$y_{i}$ on $1, w_{i}, \mathbf{x}_{i}, w_{i}\left(\mathbf{x}_{i}-\overline{\mathbf{x}}_{\mathscr{R}}\right) ;$\\
the coefficient on $w_{i}$ is $\hat{\tau}_{a t e, \mathscr{R}, \text { reg }}$, and we can use the heteroskedasticity-robust standard to form a $t$ statistic or confidence intervals. Technically, we should adjust for the estimation error in $\overline{\mathbf{x}}_{\mathscr{R}}$, but the adjustment may have a small effect (see Problem 6.10). Or, we can use the bootstrap to account for this uncertainty, too. Of course, we get $\hat{\tau}_{\text {ate,reg }}$ by using $\overline{\mathbf{x}}_{\mathscr{R}}=\overline{\mathbf{x}}$. If we are mainly interested in the ATE over a subpopulation, it may be better-by way of increasing overlap-to estimate the regression functions on the restricted sample $\mathbf{x}_{i} \in \mathscr{R}$.

If the range of $y$ is substantively restricted, we may wish to exploit that in estimation. For example, if $y$ is a binary variable, or a fractional response, we can use a logit or probit function and the Bernoulli quasi MLE. If these functions are $G\left(\alpha_{g}+\mathbf{x} \boldsymbol{\beta}_{g}\right)$,\\
$g=0,1$ with $0<G(\cdot)<1$, let $\left(\hat{\alpha}_{g}, \hat{\boldsymbol{\beta}}_{g}\right)$ be the Bernoulli QMLEs on the two subsamples. Then


\begin{equation*}
\hat{\tau}_{\text {ate, reg }}=N^{-1} \sum_{i=1}^{N}\left[G\left(\hat{\alpha}_{1}+\mathbf{x}_{i} \hat{\boldsymbol{\beta}}_{1}\right)-G\left(\hat{\alpha}_{0}+\mathbf{x}_{i} \hat{\boldsymbol{\beta}}_{0}\right)\right] \tag{21.36}
\end{equation*}


is the estimate of $\tau_{\text {ate }}$. As in the general case, $G\left(\hat{\alpha}_{1}+\mathbf{x}_{i} \hat{\boldsymbol{\beta}}_{1}\right)-G\left(\hat{\alpha}_{0}+\mathbf{x}_{i} \hat{\boldsymbol{\beta}}_{0}\right)$ is the difference in average response for unit $i$, given covariates $\mathbf{x}_{i}$, in the two treatment states (regardless of unit $i$ 's actual treatment status). If $y \geq 0$-such as a count variable, but not restricted to that case-an exponential regression function is sensible, in which case\\
$\hat{\tau}_{\text {ate, reg }}=N^{-1} \sum_{i=1}^{N}\left[\exp \left(\hat{\alpha}_{1}+\mathbf{x}_{i} \hat{\boldsymbol{\beta}}_{1}\right)-\exp \left(\hat{\alpha}_{0}+\mathbf{x}_{i} \hat{\boldsymbol{\beta}}_{0}\right)\right]$,\\
where the parameter estimates can be obtained from a quasi-MLE procedure, such as Poisson or gamma regression, or nonlinear least squares.

Equations (21.33), (21.36), and (21.37) make it clear that when parametric models are used for the means, it is possible to estimate $\tau_{\text {ate }}$ while essentially ignoring the overlap assumption. When using parametric models we are at least implicitly assuming that, say, $m_{1}\left(\cdot, \boldsymbol{\delta}_{1}\right)$ holds for all $\mathbf{x} \in \mathscr{X}$, even though we only use the treated subsample to obtain $\hat{\boldsymbol{\delta}}_{1}$. But as in any regression context, we should be careful in extrapolating the estimated mean functions to values of $\mathbf{x}$ far from those used to obtain $\hat{\boldsymbol{\delta}}_{1}$. The resulting estimate of $\tau_{\text {ate }}$ can be sensitive to the exact specification. For example, in the linear case it can be shown (for example, Imbens and Wooldridge, 2009) that


\begin{equation*}
\hat{\tau}_{\text {ate, reg }}=\left(\bar{y}_{1}-\bar{y}_{0}\right)-\left(\overline{\mathbf{x}}_{1}-\overline{\mathbf{x}}_{0}\right)\left(f_{0} \hat{\boldsymbol{\beta}}_{1}-f_{1} \hat{\boldsymbol{\beta}}_{0}\right) \tag{21.38}
\end{equation*}


where $f_{0}=N_{0} /\left(N_{0}+N_{1}\right)$ is the fraction of control observations and $f_{1}$ is the fraction of treated observations. If the difference in means $\overline{\mathbf{x}}_{1}-\overline{\mathbf{x}}_{0}$ is large, changes in the slope estimates can have a large effect on $\hat{\tau}_{\text {ate, reg }}$. Therefore, as a general rule, one should not rely on parametric specifications as a way of overcoming poor overlap in the covariate distributions.

\subsection*{21.3.3 Propensity Score Methods}
The expressions derived in Section 21.3.1 to establish identification of $\tau_{\text {ate }}$ and $\tau_{\text {att }}$ based on the propensity score lead directly to estimators. Practically, though, we need to estimate the propensity score function, $p(\cdot)$. For the moment, let $\hat{p}(\mathbf{x})$ denote such an estimator for any $\mathbf{x} \in \mathscr{X}$ obtained using the random sample $\left\{\left(w_{i}, \mathbf{x}_{i}\right): i=\right.$\\
$1, \ldots, N\}$. Then equation (21.22) suggests the estimator\\
$\hat{\tau}_{\text {ate }, p s w}=N^{-1} \sum_{i=1}^{N}\left[\frac{w_{i} y_{i}}{\hat{p}\left(\mathbf{x}_{i}\right)}-\frac{\left(1-w_{i}\right) y_{i}}{1-\hat{p}\left(\mathbf{x}_{i}\right)}\right]=N^{-1} \sum_{i=1}^{N} \frac{\left[w_{i}-\hat{p}\left(\mathbf{x}_{i}\right)\right] y_{i}}{\hat{p}\left(\mathbf{x}_{i}\right)\left[1-\hat{p}\left(\mathbf{x}_{i}\right)\right]}$,\\
while equation (21.26) suggests\\
$\hat{\tau}_{a t t, p s w}=N^{-1} \sum_{i=1}^{N} \frac{\left[w_{i}-\hat{p}\left(\mathbf{x}_{i}\right)\right] y_{i}}{\hat{\rho}\left[1-\hat{p}\left(\mathbf{x}_{i}\right)\right]}$,\\
where $\hat{\rho}=N_{1} / N$ is the fraction of treated units in the sample. As is evident from the formulas, each estimate is simple to compute given the fitted probabilities of treatment, $\hat{p}\left(\mathbf{x}_{i}\right)$. Interestingly, the estimator in equation (21.39) is the same as an estimator due to Horvitz and Thompson (1952) for handling nonrandom sampling. Not surprisingly, these estimators are generally consistent under Assumptions ATE.1' and ATE. 2 or Assumptions ATT.1' and ATT.2, respectively, and suitable regularity conditions. If $\hat{p}(\cdot)$ is obtained via parametric methods, consistency generally follows from Lemma 12.1.

As for estimating the propensity score, Rosenbaum and Rubin (1983) suggest using a flexible logit model, where various functions of $\mathbf{x}$-for example, levels, squares, and interactions-are included. For discrete components of $\mathbf{x}$, one might define a set of dummy variables indicating the different possible values, and then interact these with continuous covariates and other dummy variables defined similarly. Clearly the sample size is important in deciding on how flexible the model can be. If we use a flexible logit or probit, or any index function $G(\cdot)$ with $0<G(z)<1$ for all $z \in \mathbb{R}$, then there is no danger of $\hat{p}(\mathbf{x})=0$ or $\hat{p}(\mathbf{x})=1$, but ruling out zero and one from the fitted probabilities might simply mask the failure of overlap in the population-just as when we use parametric versions of the regression functions and then extrapolate beyond values of $\mathbf{x}$ used to estimate the two functions.

Hirano, Imbens, and Ridder (2003) (HIR for short) study a nonparametric version of the Rosenbaum and Rubin (1983) approach. In particular, they study $\hat{\tau}_{a t e, p s w}$ when the propensity score is estimated using a flexible logit model, but they explicitly allow the number of functions of the covariates in the logit function to increase as a function of the sample size. Under regularity conditions and an assumption controlling the number of terms in the logit estimation, HIR show that their nonparametric version of the Horvitz-Thompson estimator achieves the semiparametric efficiency bound due to Hahn (1998). Remarkably, the HIR estimator is more efficient-often nontrivially so-than the "estimator" that uses the known propensity scores, $p\left(\mathbf{x}_{i}\right)$, in place of the nonparametric fitted probabilities, $\hat{p}\left(\mathbf{x}_{i}\right)$. Moreover, the HIR estimator is\\
strictly more efficient (except in special cases) than the estimator that uses a maximum likelihood estimator for $\hat{p}(\cdot)$. This fact is related to a result by Robins and Rotnitzky (1995): even if a logit model with a given number of terms is correctly specified for $p(\cdot)$, one generally reduces the asymptotic variance of $\hat{\tau}_{\text {ate,psw }}$ by continuing to add terms even though these have no affect on $\mathrm{P}(w=1 \mid \mathbf{x})$. The HIR estimator essentially takes the Robins and Rotnitzky (1995) result to the limit. (Shortly, we will see how the efficiency gains work in a parametric setting.)

An alternative to global smoothing methods such as the HIR series estimator is to use local smoothing, such as kernel regression. Li, Racine, and Wooldridge (2009) have recently proposed using kernel smoothing that allows both continuous and discrete covariates.

It is informative to show how to obtain the asymptotic variances of $\sqrt{N}\left(\hat{\tau}_{\text {ate,psw }}-\tau_{\text {ate }}\right)$ and $\sqrt{N}\left(\hat{\tau}_{\text {att,psw }}-\tau_{\text {att }}\right)$ in the parametric cases where we assume a correctly specified model for $p(\mathbf{x})$ and use the Bernoulli MLE. For $\hat{\tau}_{a t e, p s w}$, we can directly apply the "surprising" efficiency result in Section 13.10.2. The ignorability-of-treatment assumption (conditional on $\mathbf{x}$, of course) is easily shown to imply the key conditional independence assumption in equation (13.66) (with notation properly adjusted). Because $\hat{\tau}_{a t e, p s w}$ is a sample average, we can work directly off of its "firstorder condition." Let $p(\mathbf{x}, \boldsymbol{\gamma})$ be the correctly specified propensity score model, and define the score from the first-stage propensity score estimation as\\
$\mathbf{d}_{i}=\mathbf{d}\left(w_{i}, \mathbf{x}_{i}, \boldsymbol{\gamma}\right)=\frac{\nabla_{\gamma} p\left(\mathbf{x}_{i}, \boldsymbol{\gamma}\right)^{\prime}\left[w_{i}-p\left(\mathbf{x}_{i}, \boldsymbol{\gamma}\right)\right]}{p\left(\mathbf{x}_{i}, \boldsymbol{\gamma}\right)\left[1-p\left(\mathbf{x}_{i}, \boldsymbol{\gamma}\right)\right]}$,\\
where $y$ is used here to also denote the true population value. Further, define\\
$k_{i}=\frac{\left[w_{i}-p\left(\mathbf{x}_{i}, \boldsymbol{\gamma}\right)\right] y_{i}}{p\left(\mathbf{x}_{i}, \boldsymbol{\gamma}\right)\left[1-p\left(\mathbf{x}_{i}, \boldsymbol{\gamma}\right)\right]}$,\\
which are the summands in $\hat{\tau}_{a t e, p s w}$ with the true propensity score inserted. Then the asymptotic variance of $\sqrt{N}\left(\hat{\tau}_{\text {ate }, p s w}-\tau_{\text {ate }}\right)$ is simply $\operatorname{Var}\left(e_{i}\right)$, where $e_{i}$ is the population residual from the regression $k_{i}$ on $\mathbf{d}_{i}^{\prime}$. We can easily estimate this variance by using the sample version. Let


\begin{equation*}
\hat{\mathbf{d}}_{i}=\mathbf{d}\left(w_{i}, \mathbf{x}_{i}, \hat{\gamma}\right)=\frac{\nabla_{\gamma} p\left(\mathbf{x}_{i}, \hat{\gamma}\right)^{\prime}\left[w_{i}-p\left(\mathbf{x}_{i}, \hat{\gamma}\right)\right]}{p\left(\mathbf{x}_{i}, \hat{\gamma}\right)\left[1-p\left(\mathbf{x}_{i}, \hat{\gamma}\right)\right]} \tag{21.42}
\end{equation*}


be the estimated score, and let


\begin{equation*}
\hat{k}_{i}=\frac{\left[w_{i}-\hat{p}\left(\mathbf{x}_{i}\right)\right] y_{i}}{\hat{p}\left(\mathbf{x}_{i}\right)\left[1-\hat{p}\left(\mathbf{x}_{i}\right)\right]} \tag{21.43}
\end{equation*}


Then, obtain the OLS residuals, $\hat{e}_{i}$, from the regression\\
$\hat{k}_{i} \quad$ on $\quad 1, \hat{\mathbf{d}}_{i}^{\prime}, \quad i=1, \ldots, N$,\\
where the inclusion of unity ensures that we remove the sample average of $\hat{k}_{i}$ (which is $\hat{\tau}_{\text {ate,psw }}$ ) in estimating the variance. Given these residuals, the asymptotic standard error of $\hat{\tau}_{a t e, p s w}$ is\\
$\left[N^{-1} \sum_{i=1}^{N} \hat{e}_{i}^{2}\right]^{1 / 2} / \sqrt{N}$.

The adjustment is particularly simple in the case of logit estimation of the propensity score because then\\
$\hat{\mathbf{d}}_{i}^{\prime}=\mathbf{h}_{i}\left(w_{i}-\hat{p}_{i}\right)$,\\
where $\mathbf{h}_{i} \equiv \mathbf{h}\left(\mathbf{x}_{i}\right)$ is the $1 \times R$ vector of covariates (including unity) and $\hat{p}_{i}=\Lambda\left(\mathbf{h}_{i} \hat{\gamma}\right)= \exp \left(\mathbf{h}_{i} \hat{\gamma}\right) /\left[1+\exp \left(\mathbf{h}_{i} \hat{\gamma}\right)\right]$. We can also see that, as we add elements to $\mathbf{h}_{i}$, the population residuals (and the sample conterparts) from the regression $k_{i}$ on $\mathbf{h}_{i}\left(w_{i}-p_{i}\right)$ will shrink provided the extra elements in $\mathbf{h}_{i}\left(w_{i}-p_{i}\right)$ are partially correlated with $k_{i}$. This is generally the case even though the extra functions in $\mathbf{h}_{i}$ have zero population coefficients in $\mathrm{P}\left(w_{i}=1 \mid \mathbf{x}_{i}\right)$. This was the result noted by Robins and Rotnitzky (1995).

If we ignore estimation of the propensity score, we effectively treat $\left\{\hat{k}_{i}: i=\right. 1, \ldots, N\}$ as being drawn from a random sample, and we use its sample average to estimate the population mean, $\tau_{\text {ate }}$. The naive standard error that we obtain is\\
$\left[N^{-1} \sum_{i=1}^{N}\left(\hat{k}_{i}-\hat{\tau}_{a t e, p s w}\right)^{2}\right]^{1 / 2} / \sqrt{N}$,\\
and this is at least as large as expression (21.45), and sometimes much larger. In the population, the comparision is $\operatorname{Var}\left(k_{i}\right)$ versus $\operatorname{Var}\left[k_{i}-\mathrm{L}\left(k_{i} \mid \mathbf{d}_{i}^{\prime}\right)\right]$, where $\mathrm{L}\left(k_{i} \mid \mathbf{d}_{i}^{\prime}\right)$ is the linear projection. Netting out $\mathrm{L}\left(k_{i} \mid \mathbf{d}_{i}^{\prime}\right)$ produces a smaller variance unless $k_{i}$ and $\mathbf{d}_{i}$ are uncorrelated.

We can also find an appropriate standard error for $\hat{\tau}_{\text {att,psw }}$. Write\\
$\hat{\tau}_{a t t, p s w}=\hat{\rho}^{-1} N^{-1} \sum_{i=1}^{N} \hat{q}_{i}$,\\
where $\hat{q}_{i}=\left[w_{i}-p\left(\mathbf{x}_{i}, \hat{\gamma}\right)\right] y_{i} /\left[1-p\left(\mathbf{x}_{i}, \hat{\gamma}\right)\right]$. Now

$$
\begin{aligned}
\sqrt{N}\left(\hat{\tau}_{a t t, p s w}-\tau_{a t t}\right) & =\hat{\rho}^{-1} N^{-1 / 2} \sum_{i=1}^{N}\left(\hat{q}_{i}-\hat{\rho} \tau_{a t t}\right)=\hat{\rho}^{-1}\left[N^{-1 / 2} \sum_{i=1}^{N} \hat{q}_{i}-\tau_{a t t} N^{-1 / 2} \sum_{i=1}^{N} w_{i}\right] \\
& =\hat{\rho}^{-1}\left[N^{-1 / 2} \sum_{i=1}^{N} r_{i}-\tau_{a t t} N^{-1 / 2} \sum_{i=1}^{N} w_{i}\right]+o_{p}(1)
\end{aligned}
$$

where the $r_{i}$ are the population residuals from regressing $q_{i}$ on $\mathbf{d}_{i}^{\prime}$ (using the same argument as for $\left.\hat{\tau}_{a t e, p s w}\right)$. Therefore,

$$
\sqrt{N}\left(\hat{\tau}_{a t t, p s w}-\tau_{a t t}\right)=\rho^{-1} N^{-1 / 2} \sum_{i=1}^{N}\left(r_{i}-\tau_{a t t} w_{i}\right)+o_{p}(1)
$$

and so estimation of the asymptotic variance is straightforward. The asymptotic standard error of $\hat{\tau}_{a t t, p s w}$ is


\begin{equation*}
\hat{\rho}^{-1}\left[N^{-1} \sum_{i=1}^{N}\left(\hat{r}_{i}-\hat{\tau}_{a t t, p s w} w_{i}\right)^{2}\right]^{1 / 2} \sqrt{N}, \tag{21.48}
\end{equation*}


where $\hat{r}_{i}$ are the residuals from the regression $\hat{q}_{i}$ on $\hat{\mathbf{d}}_{i}^{\prime}$.\\
A different use of the estimated propensity scores is in regression adjustment. A simple, still somewhat popular estimate is obtained from the OLS regression


\begin{equation*}
y_{i} \quad \text { on } \quad 1, w_{i}, \hat{p}\left(\mathbf{x}_{i}\right), \quad i=1, \ldots, N ; \tag{21.49}
\end{equation*}


the coefficient on $w_{i}$, say, $\hat{\tau}_{\text {ate,psreg }}$, is the estimate of $\tau_{\text {ate }}$. The idea is that the estimated propensity score should be sufficient in controlling for correlation between the treatment, $w_{i}$, and the covariates, $\mathbf{x}_{i}$. As it turns out, there is a simple characterization of when $\hat{\tau}_{\text {ate,pseg }}$ consistently estimates $\tau_{\text {ate }}$. The following is a special case of Wooldridge (1999c, Proposition 3.2).

PROPOSITION 21.3: In addition to Assumptions ATE1.1' and ATE.2, assume that $\tau_{\text {ate }}(\mathbf{x})=\mathrm{E}\left(y_{1}-y_{0} \mid \mathbf{x}\right)$ is uncorrelated with $\operatorname{Var}(w \mid \mathbf{x})=p(\mathbf{x})[1-p(\mathbf{x})]$. Then, under standard regularity conditions that ensure that $\hat{\gamma}$ (in a parametric estimation problem) is consistent and $\sqrt{N}$-asymptotically normal, $\hat{\tau}_{\text {ate, psreg }}$ is consistent for $\tau_{\text {ate }}$ and $\sqrt{N}$ asymptotically normal.

The assumption that $\tau_{\text {ate }}(\mathbf{x})=\mu_{1}(\mathbf{x})-\mu_{0}(\mathbf{x})$ is uncorrelated with $\operatorname{Var}(w \mid \mathbf{x})$ may appear unlikely, as both are functions of $\mathbf{x}$. However, remember that correlation measures linear dependence. It would not be surprising for $\tau_{\text {ate }}(\mathbf{x})$ to be monotonic in many elements of $\mathbf{x}$, while $p(\mathbf{x})[1-p(\mathbf{x})]$ is a quadratic in the propensity score. If so, the correlation between these two functions of $\mathbf{x}$ might be small. (This comment is\\
analogous to the fact that if $z$ has a symmetric distribution around zero, $z$ and $z^{2}$ are uncorrelated even though $z^{2}$ is an exact function of $z$.) In the case where the treatment effect is constant, so that $\tau_{\text {ate }}(\mathbf{x})=\tau_{\text {ate }}$, zero correlation holds, and $\hat{\tau}_{\text {ate, psreg }}$ consistently estimates $\tau_{\text {ate }}$ without further assumptions.

The result in Section 13.10.2 can be used to obtain a proper standard error for $\hat{\tau}_{\text {ate,psreg }}$-whether or not different slopes are allowed. But we also know that ignoring the propensity-score estimation results in asymptotically conservative inference. Because regression on the propensity score is computationally simple, bootstrapping the propensity score and regression estimation is attractive as a way to obtain proper standard errors and confidence intervals. By resampling the entire vector ( $y_{i}, w_{i}, \mathbf{x}_{i}$ ), even the estimation error in $\hat{\mu}_{p}$ is easily accounted for.

Rosenbaum and Rubin (1983) take a different approach to deriving regression on the propensity score and matching methods discussed in the next subsection. A key result is that under Assumption ATE.1, treatment is ignorable conditional only on the propensity score, $p(\mathbf{x})$. For completeness, we provide a simple proof.\\
proposition 21.4: Under Assumption ATE.1, $w$ is independent of $\left(y_{0}, y_{1}\right)$ conditional on $p(\mathbf{x})$. Therefore, $\mathrm{E}[y \mid w=0, p(\mathbf{x})]=\mathrm{E}\left[y_{0} \mid p(\mathbf{x})\right]$ and $\mathrm{E}[y \mid w=1, p(\mathbf{x})]= \mathrm{E}\left[y_{1} \mid p(\mathbf{x})\right]$.

Proof: Because $w$ is binary, it suffices to show that $P\left[w=1 \mid y_{0}, y_{1}, p(\mathbf{x})\right]= P[w=1 \mid p(\mathbf{x})]$ or $\mathrm{E}\left[w \mid y_{0}, y_{1}, p(\mathbf{x})\right]=\mathrm{E}[w \mid p(\mathbf{x})]$. But, under Assumption ATE.1, $\mathrm{E}\left[w \mid y_{0}, y_{1}, \mathbf{x}\right]=\mathrm{E}[w \mid \mathbf{x}]=p(\mathbf{x})$, where the second equality follows because $w$ is binary. By iterated expectations,\\
$\mathrm{E}\left[w \mid y_{0}, y_{1}, p(\mathbf{x})\right]=\mathrm{E}\left[\mathrm{E}\left(w \mid y_{0}, y_{1}, \mathbf{x}\right) \mid y_{0}, y_{1}, p(\mathbf{x})\right]=\mathrm{E}\left[p(\mathbf{x}) \mid y_{0}, y_{1}, p(\mathbf{x})\right]=p(\mathbf{x})$,\\
which shows that ignorability of treatment holds conditional on $p(\mathbf{x})$.\\
Next, write $y=(1-w) y_{0}+w y_{1}$, as usual. Then

$$
\begin{aligned}
\mathrm{E}[y \mid w, p(\mathbf{x})] & =(1-w) \mathrm{E}\left[y_{0} \mid w, p(\mathbf{x})\right]+w \mathrm{E}\left[y_{1} \mid w, p(\mathbf{x})\right] \\
& =(1-w) \mathrm{E}\left[y_{0} \mid p(\mathbf{x})\right]+w \mathrm{E}\left[y_{1} \mid p(\mathbf{x})\right]
\end{aligned}
$$

Inserting $w=0$ and $w=1$, respectively, gives the results for the conditional means.\\
Proposition 21.4 has many uses. For one, it implies that, since $\mathrm{E}\left[y_{g} \mid p(\mathbf{x})\right]$ is identified, $g=0,1$, we can identify $\tau_{\text {ate }}=\mathrm{E}\left\{\mathrm{E}\left[y_{1} \mid p(\mathbf{x})\right]-\mathrm{E}\left[y_{0} \mid p(\mathbf{x})\right]\right\}$. In particular, let


\begin{equation*}
r_{0}(p) \equiv \mathrm{E}[y \mid w=0, p(\mathbf{x})=p], \quad r_{1}(p) \equiv \mathrm{E}[y \mid w=1, p(\mathbf{x})=p] \tag{21.50}
\end{equation*}


be the two regression functions that can be identified given the data. Given consistent estimators $\hat{r}_{0}\left(\hat{p}_{i}\right), \hat{r}_{1}\left(\hat{p}_{i}\right)$, where $\hat{p}_{i} \equiv \hat{p}\left(\mathbf{x}_{i}\right)$ are the estimated propensity scores, we get as a general estimate\\
$\hat{\tau}_{\text {ate }, \text { psreg }}=N^{-1} \sum_{i=1}^{N}\left[\hat{r}_{1}\left(\hat{p}_{i}\right)-\hat{r}_{0}\left(\hat{p}_{i}\right)\right]$.

Generally, this estimator is consistent and $\sqrt{N}$-asymptotically normal provided we have the models for $r_{g}(\cdot), g=0,1$ and the propensity score correctly specified (or we use suitable nonparametric methods). Given the nature of $\hat{\tau}_{\text {ate,psreg }}$, bootstrapping is an attractive method for inference.

Heckman, Ichimura, and Todd (1998) propose local smoothers to estimate $r_{0}(\cdot)$ and $r_{1}(\cdot)$; this proposal is attractive because these are functions of a scalar. But one must also estimate the propensity score, and this is often a high-dimensional problem. Hahn (1998) proposed series estimation of $r_{0}(\cdot)$ and $r_{1}(\cdot)$ along with series estimation of $p(\cdot)$. A simpler approach is just to use parametric functions and parametric asymptotics.

If $r_{0}(\cdot)$ and $r_{1}(\cdot)$ are linear, say, $r_{g}(p)=\eta_{g}+\pi_{g} p, g=0,1$, then estimation is straightforward. Two separate regressions can be run, and then the differences in predicted values, $\hat{r}_{1}\left(\hat{p}_{i}\right)-\hat{r}_{0}\left(\hat{p}_{i}\right)$, are averaged to get $\hat{\tau}_{\text {ate,psreg }}$. Equivalently, run the regression


\begin{equation*}
y_{i} \quad \text { on } \quad 1, w_{i}, \hat{p}\left(\mathbf{x}_{i}\right), w_{i} \cdot\left[\hat{p}\left(\mathbf{x}_{i}\right)-\hat{\rho}\right], \quad i=1, \ldots, N, \tag{21.52}
\end{equation*}


where $\hat{\rho}$ is a consistent estimate of $\rho=\mathrm{E}\left[p\left(\mathbf{x}_{i}\right)\right]=\mathrm{P}\left(w_{i}=1\right)$. If $\hat{p}\left(\mathbf{x}_{i}\right)$ is from a logit that includes an intercept, the two natural estimates, $\bar{w}$ and $N^{-1} \sum_{i=1}^{N} \hat{p}\left(\mathbf{x}_{i}\right)$, are identical. The coefficient on $w_{i}$ is $\hat{\tau}_{\text {ate,psreg }}$. If we ignore estimation of $\rho$, the usual heteroskedasticity-robust standard error will be conservative, but bootstrapping can be obtained to get the proper standard error.

The linear models for $\mathrm{E}\left[y_{g} \mid p(\mathbf{x})\right]$ might be too restrictive because $0<p(\mathbf{x})<1$. If $y$ has substantial variation, different functions might be needed. One can always use polynomials in $\hat{p}_{i}$-as formally studied in a series setting by Hahn (1998). A log-odds transformation might improve the fit, where $\widehat{\operatorname{lodd}} s_{i} \equiv \log \left[\hat{p}_{i} /\left(1-\hat{p}_{i}\right)\right]$ and functions of it (such as polynomials) are used as regressors. It is easy to use two separate regressions and then average the differences in predicted values.

Given that $r_{1}(\cdot)$ is estimated using the treated subsample and $r_{0}(\cdot)$ using the control sample, we can see the nature of the overlap assumption. In effect, for each value of the propensity score in the interval $(0,1)$, we should observe both treated and untreated units. We can easily use the estimated propensity scores to determine whether lack of overlap is a problem. For example, we can compute the normalized\\
difference of the propensity scores (or its log-odds ratio), as in equation (21.29), using the propensity score in place of the covariates. It is also informative to plot histograms of the $\hat{p}_{i}$ for the control and treatment groups. We hope to rule out situations where the histograms show bins with data on the control units but no or very few data for the corresponding treatment units, and vice versa.

Regression on the propensity score seems attractive because, compared with regression adjustment using $\mathbf{x}$, it reduces the problem of controlling for a possibly large set of covariates, $\mathbf{x}$, to controlling for a single function of them, $p(\mathbf{x})$. But the parsimony afforded by propensity score regressions is somewhat illusory because we should use a flexible model for the propensity score. Even if we settle on logit or probit, we would typically include flexible functions of $\mathbf{x}$, just as if we were doing regression adjustment. One might argue that, because the treatment $w_{i}$ is binary, we have better leads on suitable models for $\mathrm{P}(w=1 \mid \mathbf{x})$. Nevertheless, as we discussed in Section 21.3.2, we can account for the nature of $y$ when applying regression adjustment. If we use nonparametric approaches, estimating either $p(\mathbf{x})$ or the regression functions $m_{g}(\mathbf{x})=\mathrm{E}(y \mid w=g, \mathbf{x})$ generally requires high-dimensional nonparametric estimation. (Incidentally, if we use a linear probability model for $p(\mathbf{x})$, regressing $y_{i}$ either on the covariates themselves or the propensity score gives the same estimate of $\tau_{\text {ate }}$.)

It is easy to see why regression on the propensity score is generally inefficient. Suppose that assignment is random so that $w$ is independent of $\left(y_{0}, y_{1}, \mathbf{x}\right)$, but $\mathbf{x}$ helps to predict the counterfactual outcomes. Further, assume that we actually know the propensity score, so that regression on the propensity score is $y_{i}$ on $1, w_{i}, p\left(\mathbf{x}_{i}\right)$. Because of random assignment, we do not need to include $p\left(\mathbf{x}_{i}\right)$ in order to obtain a consistent estimator of $\tau_{a t e}$ : the simple difference in sample averages is consistent. Adding $p\left(\mathbf{x}_{i}\right)$ can actually improve efficiency over the difference-in-means estimator if $p(\mathbf{x})$ appears in the linear projections $\mathrm{L}\left[y_{g} \mid 1, p(\mathbf{x})\right]$ : adding $p(\mathbf{x})$ will shrink the error variance. But $p(\mathbf{x})$ is hardly the best function of $\mathbf{x}$ to add to the regression. In the case of a constant treatment effect, the best function of $\mathbf{x}$ to add is $\mu_{0}(\mathbf{x})=\mathrm{E}\left(y_{0} \mid \mathbf{x}\right)$ because this approach leads to the smallest possible mean squared error among functions of $\mathbf{x}$. In effect, the error variance is made as small as possible. Of course, we do not know this mean function, but we can approximate it much better using flexible functions of $\mathbf{x}$ than by using a linear function of just the propensity score.

We now show how regression methods and inverse probability weighted (IPW) methods can be used to estimate treatment effects of job training on labor earnings. The underlying data are from Lalonde (1986), although the particular data set used here is from Dehejia and Wahba (1999).

Example 21.1 (Causal Effects of Job Training on Earnings): We will use two data sets for this example, JTRAIN2.RAW and JTRAIN3.RAW. The first data set is from a job training experiment back in the 1970s for men with poor labor market histories. Being in the treatment group, indicated by train $=1$, was randomly assigned. Out of 445 men, 185 were assigned to the job training program. The response variable is re78, real labor market earnings in 1978. (This variable is zero for a nontrivial fraction of the population.) Training began up to two years prior to 1978. The controls that we use here are age (in years), years of education, dummy variables for being black, being Hispanic, and marital status, and real earnings for the two years prior to the start of the program, 1974 and 1975. For simplicity, these all appear in level form without interactions or other functional forms.

JTRAIN3.RAW contains the same information as JTRAIN2.RAW, but the control group in JTRAIN3.RAW was obtained as a random sample from the Panel Study of Income Dynamics. Thus, JTRAIN3.RAW is essentially a nonexperimental version of JTRAIN2.RAW, which allows us to determine how well the nonexperimental methods for estimating treatment effects compare with experimental estimates (Lalonde's original motivation).

We use four estimation methods: a simple comparision of means, regression adjustment pooling across the control and treated groups, regression adjustment using separate regression functions, and IPW estimation using the propensity score. In each case, the controls are (age, educ, black, hisp, married, re74, re75). Note that for the difference-in-means and pooled regression adjustment there is no difference between $\hat{\tau}_{\text {ate }}$ and $\hat{\tau}_{\text {att }}$.

Before presenting the estimates, we note that there is a severe problem with overlap in JTRAIN3.RAW. For example, for the variable re75, the absolute value of the normalized difference in equation (21.29) is about 1.25, well above the Imbens and Rubin rule of 0.25 . Other covariates have normalized differences above unity, too. Therefore, unless we think the ATE is constant, we can expect problems in trying to estimate $\tau_{\text {ate }}$ and, to a lesser degree, $\tau_{\text {att }}$, even under ignorability.

Table 21.1 contains the estimates and the standard errors. For the first two estimation methods, we use the heteroskedasticity-robust standard error from the OLS regressions; by assumption, $\tau_{\text {att }}=\tau_{\text {ate }}$. For regression adjustment using separate regression functions, as well as propensity score weighting, we use 1,000 bootstrap replications in Stata 10. It is slightly more difficult to use the analytical formulas derived in equation (21.31) and expression (21.47).

Not surprisingly, using the experimental data the estimates are consistent across estimation method. Because of random assignment, $\tau_{\text {att }}=\tau_{\text {ate }}$, so it is also expected that the estimates of these parameters would be similar. The job training program

\begin{table}[h]
\begin{center}
\captionsetup{labelformat=empty}
\caption{Table 21.1}
\begin{tabular}{|l|l|l|l|l|l|l|}
\hline
\multirow[b]{2}{*}{Estimation Method} & \multicolumn{2}{|l|}{JTRAIN2} & \multicolumn{2}{|l|}{JTRAIN3 (Full Sample)} & \multicolumn{2}{|l|}{JTRAIN3 (Reduced Sample)} \\
\hline
 & $\hat{\tau}_{\text {ate }}$ & $\hat{\tau}_{\text {att }}$ & $\hat{\tau}_{\text {ate }}$ & $\hat{\tau}_{\text {att }}$ & $\hat{\tau}_{\text {ate }}$ & $\hat{\tau}_{\text {att }}$ \\
\hline
Difference in means & 1.794 (0.671) & 1.794 (0.671) & -15.205 (0.656) & -15.205 (0.656) & -5.005 (0.657) & -5.005 (0.657) \\
\hline
Pooled regression adjustment & 1.683 (0.658) & 1.683 (0.658) & 0.860 (0.767) & 0.860 (0.767) & 2.059 (0.801) & 2.059 (0.801) \\
\hline
Separate regression adjustment & 1.633 (0.642) & 1.774 (0.661) & -8.910 (3.721) & . 931 (0.794) & -2.340 (1.480) & 2.323 (0.817) \\
\hline
Propensity score weighting & 1.627 (0.637) & 1.798 (0.660) & 11.029 (40.809) & 1.627 (0.835) & 7.049 (17.295) & 1.649 (0.986) \\
\hline
Sample size & 445 & 445 & 2,675 & 2,675 & 1,162 & 1,162 \\
\hline
\end{tabular}
\end{center}
\end{table}

is estimated to increase real labor market earnings in 1978 by between $\$ 1,627$ and $\$ 1,798$. This is a huge effect considering that the average value of $r e 78$ for the untreated sample is $\$ 4,554$.

The pattern of estimates is very different when the nonexperimental data are used. It is not too surprising that the simple comparison-of-means estimate, which provides no control for self-selection into training, is negative and large. But the other estimates of $\tau_{\text {ate }}$ also appear unreliable. The estimate of $\tau_{\text {ate }}$ using propensity score weighting is suspiciously large (though very imprecise). In fact, of the 2,675 observations, the logit model for the propensity score perfectly predicts 158 of the train $=0$ outcomes. In effect, we are dividing by zero (but in practice it is a very small number). Before one uses propensity score weighting, one should study the distribution of the propensity score. As discussed earlier, one might reduce the sample to observations with, say, $\hat{p}\left(\mathbf{x}_{i}\right)$ between 0.10 and 0.90 , or 0.05 and 0.95 .

Controlling for the covariates, either through regression adjustment or propensity score weighting, provides sensible estimates of $\tau_{a t t}$, but the estimates are imprecise. The propensity score weighting is sensible for $\tau_{\text {att }}$ because small probabilities do not affect $\hat{\tau}_{\text {att }}$ (see equation (21.40)). The largest estimated propensity score in the entire sample is about 0.936 . It is clear by looking at simple summary statistics of the propensity score that lack of balance is a serious problem: in the treated subsample, the average propensity score is about 0.631 ; in the control subsample, it is about 0.027 .

To address the lack of overlap somewhat crudely, we also use a subset of the data in JTRAIN3.RAW. We use only observations where the average of re74 and re75 is less than or equal to $\$ 15,000$. This choice is essentially arbitrary but provides a mechanism for dropping men who likely have little chance of being part of such a program. The estimates are given in the last two columns of Table 21.1. Restricting\\
the sample produces more sensible estimates, but those for $\tau_{\text {ate }}$ are still unstable. Again, the three estimates of $\tau_{\text {att }}$ that account for covariates are fairly stable but not especially precise. A more careful analysis of these data would use more flexible functional forms and cause one to think more carefully about how one might restrict the sample.

\subsection*{21.3.4 Combining Regression Adjustment and Propensity Score Weighting}
In the previous two subsections, we described methods for estimating ATEs based on two strategies: the first is based on estimating $\mu_{g}(x)=\mathrm{E}\left(y_{g} \mid \mathbf{x}\right)$ for $g=0,1$ and averaging the differences in fitted values, as in equation (21.27), and the second is based on propensity score weighting, as in equation (21.39). For each approach we have discussed estimators that achieve the asymptotic efficiency bound. If we have large sample sizes relative to the dimension of $\mathbf{x}_{i}$, we might think nonparametric estimators of the conditional means or propensity score are sufficiently accurate to invoke the asymptotic efficiency results.

In other cases we might choose flexible parametric models because full nonparametric estimation is difficult. As shown earlier, one reason for viewing estimators of conditional means or propensity scores as flexible parametric models is that it simplifies standard error calculations for treatment effect estimates. But if we use standard error calculations that rely on parametric models, we should admit the possibility that those parametric models are misspecified. As it turns out, we can combine regression adjustment and propensity score methods to achieve some robustness to misspecification of the parametric models. The resulting estimator of $\tau_{\text {ate }}$ is said to be doubly robust because it only requires either the conditional mean model or the propensity score model to be correctly specified, not both.

The idea behind the doubly robust estimators is developed in Robins and Rotnitzky (1995), Robins, Rotnitzky, and Zhao (1995), and van der Laan and Robins (2003). Wooldridge (2007) provides a simple proof that double robustness holds for certain combinations of conditional mean specifications and estimation methods. To describe the approach, let $m_{0}\left(\cdot, \boldsymbol{\delta}_{0}\right)$ and $m_{1}\left(\cdot, \boldsymbol{\delta}_{1}\right)$ be parametric functions for $\mathrm{E}\left(y_{g} \mid \mathbf{x}\right), g=0,1$, and let $p(\cdot, \gamma)$ be a parametric model for the propensity score. In the first step we estimate $\boldsymbol{\gamma}$ by Bernoulli maximum likelihood and obtain the estimated propensity scores as $p\left(\mathbf{x}_{i}, \hat{\gamma}\right)$ (probably logit or probit). In the second step, we use regression or a quasi-likelihood method, where we weight by the inverse probability. For example, if we use linear functional forms for the conditional mean, to estimate $\boldsymbol{\delta}_{1}=\left(\alpha_{1}, \boldsymbol{\beta}_{1}^{\prime}\right)^{\prime}$ we would use the IPW linear least squares problem


\begin{equation*}
\min _{\alpha_{1}, \boldsymbol{\beta}_{1}} \sum_{i=1}^{N} w_{i}\left(y_{i}-\alpha_{1}-\mathbf{x}_{i} \boldsymbol{\beta}_{1}\right)^{2} / p\left(\mathbf{x}_{i}, \hat{\boldsymbol{\gamma}}\right) ; \tag{21.53}
\end{equation*}


for $\boldsymbol{\delta}_{0}$, we weight by $1 /\left[1-\hat{p}\left(\mathbf{x}_{i}\right)\right]$ and use the $w_{i}=0$ sample. Then, we estimate $\tau_{\text {ate }}$ as the average of the difference in predicted values,\\
$\hat{\tau}_{\text {ate,pswreg }}=N^{-1} \sum_{i=1}^{N}\left[\left(\hat{\alpha}_{1}+\mathbf{x}_{i} \hat{\boldsymbol{\beta}}_{1}\right)-\left(\hat{\alpha}_{0}+\mathbf{x}_{i} \hat{\boldsymbol{\beta}}_{0}\right)\right]$.

This is the same formula as linear regression adjustment, but we are using different estimates of $\alpha_{g}, \boldsymbol{\beta}_{g}, g=0,1$.

To carefully describe the double robustness result, denote the probability limits of $\hat{\boldsymbol{\gamma}}, \hat{\boldsymbol{\delta}}_{0}$, and $\hat{\boldsymbol{\delta}}_{1}$ by $\boldsymbol{\gamma}^{*}, \boldsymbol{\delta}_{0}^{*}$, and $\boldsymbol{\delta}_{1}^{*}$, respectively. Now, if the conditional mean functions are truly linear and ignorability holds conditional on $\mathbf{x}_{i}$, then weighted least squares estimators using any function of $\mathbf{x}_{i}$ consistently estimate $\boldsymbol{\delta}_{0}^{*}$ on the control sample and $\boldsymbol{\delta}_{1}^{*}$ on the treated sample. We covered the general case for missing data in Section 19.8 (see also Wooldridge, 2007). Therefore, even if $p(\mathbf{x}, \gamma)$ is arbitrarily misspecified, weighting by functions of $p\left(\mathbf{x}_{i}, \gamma^{*}\right)$ does not cause inconsistency for estimating the parameters of the correctly specified conditional mean. (Replacing $\gamma^{*}$ with $\hat{\gamma}$ does not change this claim because $\hat{\gamma}$ converges in probability to $\boldsymbol{\gamma}^{*}$.) Now if $\mathrm{E}\left(y_{g} \mid \mathbf{x}\right)= \alpha_{g}^{*}+\mathbf{x} \boldsymbol{\beta}_{g}^{*}, g=0,1$, then $\tau_{\text {ate }}=\mathrm{E}\left[\left(\alpha_{1}^{*}+\mathbf{x} \boldsymbol{\beta}_{1}^{*}\right)-\left(\alpha_{0}^{*}+\mathbf{x} \boldsymbol{\beta}_{0}^{*}\right)\right]$, which is the usual identification result for $\tau_{\text {ate }}$ for regression adjustment.

The other half of the double robustness result is more subtle. Now suppose that $p(\mathbf{x}, \gamma)$ is correctly specified for $\mathrm{P}(w=1 \mid \mathbf{x})$ but allow for the possibility that the conditional means are not linear. As we discussed in Section 19.8, the IPW estimator (under ignorability) recovers the solution to the unweighted minimization problem in the population. In the linear regression case, that means $\boldsymbol{\delta}_{g}^{*}$ minimizes $\mathrm{E}\left[\left(y_{g}-\alpha_{g}-\mathbf{x} \boldsymbol{\beta}_{g}\right)^{2}\right]$ for $g=0,1$. In other words, the $\boldsymbol{\delta}_{g}^{*}$ are the parameters in the linear projection $\mathrm{L}\left(y_{g} \mid 1, \mathbf{x}\right)$; because we include a constant in this projection, the unconditional mean of $y_{g}$ is the mean of the linear projection, $\mathrm{E}\left(y_{g}\right)=\mathrm{E}\left[\left(\alpha_{g}^{*}+\mathbf{x} \boldsymbol{\beta}_{g}^{*}\right)\right]$. Therefore, $\tau_{\text {ate }}=\mathrm{E}\left[\left(\alpha_{1}^{*}+\mathbf{x} \boldsymbol{\beta}_{1}^{*}\right)-\left(\alpha_{0}^{*}+\mathbf{x} \boldsymbol{\beta}_{0}^{*}\right)\right]$, just as before, except now we do not need the linear functions to be correctly specified for the conditional means. In other words, linear regression adjustment can still produce a consistent estimator of $\tau_{\text {ate }}$ when the conditional means are misspecified, but we must use IPW estimation with a correctly specified model for the propensity score. The argument is essentially unchanged when we replace $\mathbf{x}$ with any functions $\mathbf{h}(\mathbf{x})$.

As illustrated for the linear regression case, the key to the second part of the double robustness property-that is, when the conditional means are misspecified-is that we can still recover the counterfactual unconditional means, $\mathrm{E}\left(y_{g}\right)$, as $\mathrm{E}\left(y_{g}\right)= \mathrm{E}\left[m_{g}\left(\mathbf{x}, \boldsymbol{\delta}_{g}^{*}\right)\right]$. Thus, to extend the double robustness result to nonlinear conditional mean models, we need to find combinations of conditional mean functions and\\
objective functions with this property, where, again, the $\boldsymbol{\delta}_{g}^{*}$ solve the population optimization problem. This is a special property, but a couple of very useful cases are known. They turn out to be quasi-maximum likelihood estimators in particular linear exponential family distributions with particular conditional mean functions. We already saw the linear case with the least squares (normal log likelihood) objective function. If $y_{g}$ is a binary or fractional response, the mean function that delivers double robustness, along with the Bernoulli quasi-log likelihood, is the logistic function


\begin{equation*}
m_{g}\left(\mathbf{x}, \boldsymbol{\delta}_{g}\right)=\Lambda\left[\alpha_{g}+\mathbf{h}(\mathbf{x}) \boldsymbol{\beta}_{g}\right] \tag{21.55}
\end{equation*}


where $\mathbf{h}(\mathbf{x})$ can be any function of $\mathbf{x}$. In other words, for the treated group, we solve the IPW QMLE problem

$$
\min _{\alpha_{1}, \boldsymbol{\beta}_{1}} \sum_{i=1}^{N} w_{i}\left\{\left(1-y_{i}\right) \log \left[1-\Lambda\left(\alpha_{1}+\mathbf{h}_{i} \boldsymbol{\beta}_{1}\right)\right]+y_{i} \log \left[\Lambda\left(\alpha_{1}+\mathbf{h}_{i} \boldsymbol{\beta}_{1}\right)\right]\right\} / p\left(\mathbf{x}_{i}, \hat{\boldsymbol{\gamma}}\right),
$$

where $\mathbf{h}_{i} \equiv \mathbf{h}\left(\mathbf{x}_{i}\right)$. Once we have used IPW estimation in each case, the ATE is estimated as before:

$$
\hat{\tau}_{\text {ate,pswreg }}=N^{-1} \sum_{i=1}^{N}\left[\Lambda\left(\hat{\alpha}_{1}+\mathbf{h}_{i} \hat{\boldsymbol{\beta}}_{1}\right)-\Lambda\left(\hat{\alpha}_{0}+\mathbf{h}_{i} \hat{\boldsymbol{\beta}}_{0}\right)\right] .
$$

If $\mathrm{E}\left(y_{g} \mid \mathbf{x}\right)=\Lambda\left[\alpha_{g}^{*}+\mathbf{h}(\mathbf{x}) \boldsymbol{\beta}_{g}^{*}\right], g=0,1$ or $\mathrm{P}(w=1 \mid \mathbf{x})=p\left(\mathbf{x}, \gamma^{*}\right)$, then $\hat{\tau}_{\text {ate, pswreg }} \xrightarrow{p} \tau_{\text {ate }}$. We might very well use a logit model for $\mathrm{P}(w=1 \mid \mathbf{x})$, but that is not necessary. Also, notice that if we replace $\Lambda(\cdot)$ with, say, the standard normal cdf, $\Phi(\cdot)$, in equation (21.55), we lose the double robustness property.

If $y_{g}$ is a nonnegative response-it could be continuous, discrete, or have both features-an exponential mean function coupled with the Poisson QLL delivers double robustness. With $\mathbf{h}_{i}=\mathbf{h}\left(\mathbf{x}_{i}\right)$ functions of $\mathbf{x}_{i}$, we solve an IPW Poisson estimation,

$$
\min _{\alpha_{1}, \boldsymbol{\beta}_{1}} \sum_{i=1}^{N} w_{i}\left[y_{i}\left(\alpha_{1}+\mathbf{h}_{i} \boldsymbol{\beta}_{1}\right)-\exp \left(\alpha_{1}+\mathbf{h}_{i} \boldsymbol{\beta}_{1}\right)\right] / p\left(\mathbf{x}_{i}, \hat{\gamma}\right),
$$

for the treated sample, and similarly for the control sample. The average treatment effect then has the same form as equation (21.37).

Why do the Bernoulli and Poisson QMLEs with, respectively, the logistic and exponential mean functions yield doubly robust estimators of $\tau_{\text {ate }}$ ? Again, the first half of double robustness is straightforward. The Bernoulli and Poisson QMLEs are fully robust for estimating the parameters of a correctly specified mean regardless of\\
the nature of $y_{i}$ (and with any mean function); weighting by a strictly positive function of $\mathbf{x}_{i}$, either $1 / p\left(\mathbf{x}_{i}, \gamma^{*}\right)$ or $1 /\left[1-p\left(\mathbf{x}_{i}, \gamma^{*}\right)\right]$, does not change the robustness property when assignment is ignorable. This reasoning gives the first half of the double robustness. For the second half, we need the specific conditional mean functions for the corresponding QLL. When $p(\mathbf{x}, \boldsymbol{\gamma})$ is correctly specified, so that $\mathrm{P}(w=1 \mid \mathbf{x})= p\left(\mathbf{x}, \gamma^{*}\right)$, the IPW estimator consistently estimates the solution to the unweighted population problem-as always. The key is that for the two combinations of mean/ QLLs just described, the probability limits satisfy $\mathrm{E}\left(y_{g}\right)=\mathrm{E}\left[m_{g}\left(\mathbf{x}, \boldsymbol{\delta}_{g}^{*}\right)\right]$. These are easy to establish by studying the first-order conditions for $\boldsymbol{\delta}_{g}^{*}$. For example, in the Poisson case, $\boldsymbol{\delta}_{g}^{*}$ satisfies (for a random draw $i$ )\\
$\mathrm{E}\left[\left(1, \mathbf{h}_{i}\right)^{\prime} y_{i g}\right]=\mathrm{E}\left[\left(1, \mathbf{h}_{i}\right)^{\prime} \exp \left(\alpha_{g}^{*}+\mathbf{h}_{i} \boldsymbol{\beta}_{g}^{*}\right)\right]$,\\
and the first of these equations is simply $\mathrm{E}\left(y_{i g}\right)=\mathrm{E}\left[\exp \left(\alpha_{g}^{*}+\mathbf{h}_{i} \boldsymbol{\beta}_{g}^{*}\right)\right]$. A similar calculation establishes $\mathrm{E}\left(y_{g}\right)=\mathrm{E}\left[\Lambda\left(\alpha_{g}^{*}+\mathbf{h}(\mathbf{x}) \boldsymbol{\beta}_{g}^{*}\right)\right]$ in the logistic/Bernoulli case. (The sample analogue of these population conditions also holds: the sample average of $y_{i}$ is equal to the average of the fitted values.)

For the previous three cases discussed, the mean functions for the normal (squared residual), Bernoulli, and Poisson quasi-log likelihoods correspond to the "canonical link" functions in the language of generalized linear models. This correspondence is not a coincidence. The mean function associated with a canonical link always has the property $\mathrm{E}\left(y_{g}\right)=\mathrm{E}\left[m_{g}\left(\mathbf{x}, \boldsymbol{\delta}_{g}^{*}\right)\right]$ provided a constant is included in the index function-as would always be the case in treatment effect applications. It is important to remember that these three conditional mean/QLL combinations can be used for a variety of response variables. As a practical matter, we should ensure that the chosen mean function is logically consistent with the nature of $y_{g}$. If $y_{g}$ has unbounded support and takes on positive and negative values, a linear model seems natural. If $y_{g}$ is binary or fractional (with possible mass at zero or one), the logit approach seems natural. And if $y_{g} \geq 0$ and unbounded, the exponential function (combined with the Poisson QLL) seems natural. Remember, the Poisson QMLE can be applied to any kind of $y_{g}$.

Because the estimates of $\tau_{\text {ate }}$ take the form of equation (21.30), we can use equation (21.31) - with proper formulas for $\hat{\mathbf{V}}_{0}$ and $\hat{\mathbf{V}}_{1}$-to compute an asymptotic standard error for $\hat{\tau}_{\text {ate,pswreg }}$. If the conditional means are correctly specified, the usual robust variance matrix is valid for the parameters using the weighted quasi-MLEs. If the conditional means are misspecified but the propensity score is correctly specified, a better (and smaller) estimate of the asymptotic variances of the $\hat{\boldsymbol{\delta}}_{g}$ is obtained by netting out the gradient of the propensity score log likelihood from the weighted,\\
selected score for the QMLE, just as in Section 19.8; see also Wooldridge (2007). We might just use the variance matrix that ignores estimation of $\gamma^{*}$ because, at worst, it is conservative. However, there is no harm in adjusting for the MLE estimation of $\gamma^{*}$ because, if the means are correctly specified, the adjustment will be minimal in large samples. Of course, bootstrapping the two-step procedure and the formula for $\hat{\tau}_{\text {ate,pswreg }}$ is easy and provides asymptotically correct inference.

\subsection*{21.3.5 Matching Methods}
The motivation for matching estimators is similar to regression adjustment. In particular, for each $i$, we impute values for the counterfactuals, $y_{i 0}$ and $y_{i 1}$. Matching estimators use the observed outcomes when possible. In other words, if we let $\hat{y}_{i 0}$ and $\hat{y}_{i 1}$ denote the imputed values, $\hat{y}_{i 0}=y_{i}$ when $w_{i}=0$ and $\hat{y}_{i 1}=y_{i}$ when $w_{i}=1$. Generally, matching estimators take the forms


\begin{equation*}
\hat{\tau}_{\text {ate, match }}=N^{-1} \sum_{i=1}^{N}\left(\hat{y}_{i 1}-\hat{y}_{i 0}\right) \tag{21.56}
\end{equation*}


$\hat{\tau}_{a t t, \text { match }}=N_{1}^{-1} \sum_{i=1}^{N} w_{i}\left(y_{i}-\hat{y}_{i 0}\right)$,\\
where the latter formula uses the fact that $y_{i 1}=y_{i}$ for the treated subsample. (In other words, we never need to impute $y_{i 1}$ for the treated subsample.)

The key question in matching is how to impute $y_{i 0}$ for the treated units (for both $\tau_{\text {ate }}$ and $\tau_{\text {att }}$ ) and how to impute $y_{i 1}$ for the control units (for $\tau_{\text {ate }}$ ). Abadie and Imbens (2006) consider several approaches, each of which involves finding one or more matches based on the covariate values. In the simplest case, a single match is found for each observation. For concreteness, suppose $i$ is a treated observation $\left(w_{i}=1\right)$. Then $\hat{y}_{i 1}=y_{i}$ and $\hat{y}_{i 0}=y_{h(i)}$ for $h(i)$ such that $w_{h(i)}=0$ and unit $h(i)$ is "closest" to unit $i$ in the sense that $\mathbf{x}_{h(i)}$ is "closest" to $\mathbf{x}_{i}$ based on a chosen metric (or distance). In other words, for the treated unit $i$ we find the "most similar" untreated observation, and use its response as our best estimate of $y_{i 0}$. Similarly, if $w_{i}=0, \hat{y}_{i 0}=y_{i}$ and $\hat{y}_{i 1}=y_{h(i)}$ where now $w_{h(i)}=1$. If we choose matches based on the full set of covariates, we call this method matching on the covariates. If we settle on a single match for each unit and we have the list of covariates, the only issue is in choosing the distance measure. A common metric is the Mahalanobis distance, which for observations $h$ and $i$ is (the square root of) $\left(\mathbf{x}_{h}-\mathbf{x}_{i}\right)^{\prime} \hat{\Sigma}_{\mathbf{x}}^{-1}\left(\mathbf{x}_{h}-\mathbf{x}_{i}\right)$, where $\hat{\Sigma}_{\mathbf{x}}$ is the sample $K \times K$ variance-covariance matrix of the covariates. Some packages use a diagonal version as the default, which gives the weighted average $\sum_{j=1}^{K}\left(x_{h j}-x_{i j}\right)^{2} / \hat{\sigma}_{j}^{2}$.

Rather than using a single "nearest neighbor," we can impute the missing values using an average of $M$ nearest neighbors. If $w_{i}=1$ then


\begin{equation*}
\hat{y}_{i 0}=M^{-1} \sum_{h \in \aleph_{M}(i)} y_{h} \tag{21.58}
\end{equation*}


where $\aleph_{M}(i)$ contains the $M$ untreated nearest matches to observation $i$, again based on the covariates. In particular, for all $h \in \aleph_{M}(i), w_{h}=0$. (With ties in the distances, there can be more than $M$ elements in $\aleph_{M}(i)$, and then $M$ is replaced with the number of elements in $\aleph_{M}(i)$.) Similarly, if $w_{i}=0$,


\begin{equation*}
\hat{y}_{i 1}=M^{-1} \sum_{h \in \mathfrak{I}_{M}(i)} y_{h} \tag{21.59}
\end{equation*}


where $\mathfrak{J}_{M}(i)$ contains the $M$ treated nearest matches to observation $i$.\\
Matching on the full set of covariates can be computationally intensive, but with modern computers the burden is manageable. The method produces a consistent estimator of $\tau_{\text {ate }}$ under the ignorability-in-mean Assumption ATE.1' along with overlap Assumption ATE.2. (Naturally, the weaker assumptions are sufficient for $\tau_{\text {att. }}$.) Matching can be motivated by the following thought experiment. Suppose that we draw a value $\mathbf{x}$ from the distribution of covariates in the population. Then, for the given covariate values, we randomly draw a control unit and a treated unit from the subpopulation and record the outcomes. The expected difference in the outcomes is\\
$\mathrm{E}(y \mid w=1, \mathbf{x})-\mathrm{E}(y \mid w=0, \mathbf{x})=\tau_{a t e}(\mathbf{x})$,\\
where the equality holds under Assumption ATE.1'. By iterated expectations, if we average the difference in outcomes across the distribution of $\mathbf{x}$, then we have $\tau_{\text {ate }}$. Matching is just the sample analogue of the thought experiment.

It is clear that lack of overlap will cause problems for matching estimators, just as with regression adjustment and propensity score weighting. Suppose $i$ is a control unit and $\mathbf{x}_{i}$ is very far from all the covariate values in the treated subsample. Then the match used to obtain $\hat{y}_{i 0}$ could be very poor, and averaging several poor matched values need not help. Fundamentally, if there are regions in the support $\mathscr{X}$ without both control and treated units, matching can produce poor results.

The large-sample properties of covariate matching have been obtained by Abadie and Imbens (2006); see also Imbens and Wooldridge (2009). Unless $K=1$ (matching on a single covariate), matching estimators are not $\sqrt{N}$-consistent, and the bias dominates the variance when $K>3$. See Abadie and Imbens (2006) for bias and variance calculations, and for a bootstrap procedure for conducting valid inference.

Because of Proposition 21.4, matching on the propensity score also produces consistent estimators of $\tau_{\text {ate }}$ and $\tau_{\text {att }}$, which is very convenient because $p(\mathbf{x})$ is a scalar in the unit interval. If we knew the propensity score, we could apply the results of Abadie and Imbens (2006); in particular, the estimator would be $\sqrt{N}$-consistent and variance calculations would be fairly straightforward, as would applying the bootstrap. Rosenbaum and Rubin (1983) first proposed propensity score matching using propensity scores obtained from a preliminary logit estimation. Unfortunately, using an estimated propensity score complicates the statistical properties of the matching estimator because matching is an inherently discontinuous operation. In particular, if standard matching methods are used-for example, $h(i)$ is chosen as the match for observation $i$ if $w_{h(i)}=1-w_{i}$ and $h(i)=\operatorname{argmin}_{h}\left|\hat{p}\left(\mathbf{x}_{h}\right)-\hat{p}\left(\mathbf{x}_{i}\right)\right|$-then bootstrapping is no longer justified. (Sometimes a different function of the propensity score, such as the log-odds ratio, is used in the matching, but that in itself does not fix the problems with bootstrapping.) Various methods of smoothing propensity score matching have been proposed; see, for example, Frölich (2004).

Matching methods can also be combined with regression adjustments. See Imbens (2004) and Imbens and Wooldridge (2009) for surveys.

Example 21.2 (Causal Effect of Job Training on Earnings): We now compute matching estimates of the causal effects from Example 21.1. We use a single-match and diagonal-weighting matrix with the inverse of the variances down the diagonal. The standard errors are computed by Stata 10 using the methods in Abadie and Imbens (2006). Table 21.2 contains the results.

Not suprisingly, on the experimental data set, the matching estimates are very similar to the regression-adjustment and propensity-score-weighting estimates obtained in Example 21.1. The estimates are somewhat less precise than the regression and propensity score estimates.

Unfortunately, like the other methods, matching does not work very well on the nonexperimental data-even when the data set is restricted such that (re74 + re75) /2 $\leq 15$. Further work would need to be done to obtain a sample with better overlap.

\begin{table}[h]
\begin{center}
\captionsetup{labelformat=empty}
\caption{Table 21.2}
\begin{tabular}{|l|l|l|l|l|l|l|}
\hline
\multirow[b]{2}{*}{Estimation Method} & \multicolumn{2}{|l|}{JTRAIN2} & \multicolumn{2}{|l|}{JTRAIN3 (Full Sample)} & \multicolumn{2}{|l|}{JTRAIN3 (Reduced Sample)} \\
\hline
 & $\hat{\tau}_{\text {ate }}$ & $\hat{\tau}_{\text {att }}$ & $\hat{\tau}_{\text {ate }}$ & $\hat{\tau}_{\text {att }}$ & $\hat{\tau}_{\text {ate }}$ & $\hat{\tau}_{\text {att }}$ \\
\hline
Covariate matching & 1.628 (0.773) & 1.824 (0.882) & -12.869 (3.815) & 0.155 (1.478) & -3.846 (2.495) & -0.232 (1.649) \\
\hline
Sample size & 445 & 445 & 2,675 & 2,675 & 1,162 & 1,162 \\
\hline
\end{tabular}
\end{center}
\end{table}

\subsection*{21.4 Instrumental Variables Methods}
We now turn to instrumental variables estimation of average treatment effects when we suspect failure of the ignorability-of-treatment assumption (ATE. 1 or ATE.1'). IV methods for estimating ATEs can be very effective if a good instrument for treatment is available. We need the instrument to predict treatment (after partialing out any controls). As we discussed in Section 5.3.1, the instrument should be redundant in a certain conditional expectation and unrelated to unobserved heterogeneity; we give precise assumptions in the following subsections.

Our primary focus in this section is on the average treatment effect defined in equation (21.1). Section 21.4.1 considers IV estimation, including methods where fitted values from estimation of a binary response model for treatment are used as instruments-and makes the case for preferring fitted values as instruments rather than regressors. Section 21.4.2 provides methods for estimating the ATE that can be used when the gain to treatment depends on unobservables (as well as observables). Two approaches are suggested; one adds a "correction function" and applies IV to the resulting equation. The other is based on a control function approach. In Section 21.4.3 we briefly discuss estimating the local average treatment effect by instrumental variables when we are not willing to make functional form or distributional assumptions.

\subsection*{21.4.1 Estimating the Average Treatment Effect Using IV}
In studying IV procedures, it is useful to write the observed outcome $y$ as\\
$y=\mu_{0}+\left(\mu_{1}-\mu_{0}\right) w+v_{0}+w\left(v_{1}-v_{0}\right)$,\\
where $\mu_{g}=E\left(y_{g}\right)$ and $v_{g}=y_{g}-\mu_{g}, g=0,1$. However, unlike in Section 21.3, we do not assume that $v_{0}$ and $v_{1}$ are mean independent of $w$, given $\mathbf{x}$. Instead, we assume the availability of instruments, which we collect in the vector $\mathbf{z}$. (Here we separate the extra instruments from the covariates, so that $\mathbf{x}$ and $\mathbf{z}$ do not overlap. In many cases $\mathbf{z}$ is a scalar, but the analysis is no easier in that case.)

If we assume that the stochastic parts of $y_{1}$ and $y_{0}$ are the same, that is, $v_{1}=v_{0}$, then the interaction term disappears (and $\tau_{\text {ate }}=\tau_{\text {att }}$ ). Without the interaction term we can use standard IV methods under weak assumptions.

ASSUMPTION ATEIV.1: (a) In equation (21.60), $v_{1}=v_{0}$; (b) $\mathrm{L}\left(v_{0} \mid \mathbf{x}, \mathbf{z}\right)=\mathrm{L}\left(v_{0} \mid \mathbf{x}\right)$; and (c) $\mathrm{L}(w \mid \mathbf{x}, \mathbf{z}) \neq \mathrm{L}(w \mid \mathbf{x})$.

All linear projections in this chapter contain unity, which we suppress for notational simplicity.

Under parts a and b of Assumption ATEIV.1, we can write\\
$y=\delta_{0}+\tau w+\mathbf{x} \boldsymbol{\beta}_{0}+u_{0}$,\\
where $\tau=\tau_{\text {ate }}$ and $u_{0} \equiv v_{0}-\mathrm{L}\left(v_{0} \mid \mathbf{x}, \mathbf{z}\right)$. By definition, $u_{0}$ has zero mean and is uncorrelated with $(\mathbf{x}, \mathbf{z})$, but $w$ and $u_{0}$ are generally correlated, which makes OLS estimation of equation (21.61) inconsistent. The redundancy of $\mathbf{z}$ in the linear projection $\mathrm{L}\left(v_{0} \mid \mathbf{x}, \mathbf{z}\right)$ means that $\mathbf{z}$ is appropriately excluded from equation (21.61); this is the part of identification that we cannot test (except indirectly using the overidentification test from Chapter 6). Part c means that $\mathbf{z}$ has predictive power in the linear projection of treatment on $(\mathbf{x}, \mathbf{z})$; this is the standard rank condition for identification from Chapter 5, and we can test it using a first-stage regression and heteroskedasticity-robust tests of exclusion restrictions. Under Assumption ATEIV.1, $\tau$ (and the other parameters in equation (21.61)), are identified, and they can be consistently estimated by 2SLS. Because the only endogenous explanatory variable in equation (21.61) is binary, equation (21.60) is called a dummy endogenous variable model (Heckman, 1978). As we discussed in Chapter 5, there are no special considerations in estimating equation (21.61) by 2SLS when the endogenous explanatory variable is binary.

Assumption ATEIV.1b holds if the instruments $\mathbf{z}$ are independent of ( $y_{0}, \mathbf{x}$ ). For example, suppose $z$ is a scalar determining eligibility in a job training program or some other social program. Actual participation, $w$, might be correlated with $v_{0}$, which could contain unobserved ability. If eligibility is randomly assigned, it is often reasonable to assume that $z$ is independent of $\left(y_{0}, \mathbf{x}\right)$. Eligibility would positively influence participation, and so Assumption ATEIV.1c should hold.

Random assignment of eligibility is no guarantee that eligibility is a valid instrument for participation. The outcome of $z$ could affect other behavior, which could feed back into $u_{0}$ in equation (21.61). For example, consider Angrist's (1990) draft lottery application, where draft lottery number is used as an instrument for enlisting. Lottery number clearly affected enlistment, so Assumption ATEIV.1c is satisfied. Assumption ATEIV.1b is also satisfied if men did not change behavior in unobserved ways that affect wage, based on their lottery number. One concern is that men with low lottery numbers may get more education as a way of avoiding service through a deferment. Including years of education in $\mathbf{x}$ effectively solves this problem. But what if men with high draft lottery numbers received more job training because employers did not fear losing them? If a measure of job training status cannot be included in $\mathbf{x}$, lottery number would generally be correlated with $u_{0}$. See Angrist, Imbens, and Rubin (1996) and Heckman (1997) for additional discussion.

As the previous discussion implies, the redundancy condition in Assumption ATEIV.1b allows the instruments $\mathbf{z}$ to be correlated with elements of $\mathbf{x}$. For example, in the population of high school graduates, if $w$ is a college degree indicator and the instrument $z$ is distance to the nearest college while attending high school, then $z$ is allowed to be correlated with other controls in the wage equation, such as geographic indicators.

Under $v_{1}=v_{0}$ and the key assumptions on the instruments, 2SLS on equation (21.61) is consistent and asymptotically normal. But if we make stronger assumptions, we can find a more efficient IV estimator.\\
assumption ATEIV.1': (a) In equation (21.60), $v_{1}=v_{0}$; (b) $\mathrm{E}\left(v_{0} \mid \mathbf{x}, \mathbf{z}\right)=\mathrm{L}\left(v_{0} \mid \mathbf{x}\right)$; (c) $\mathrm{P}(w=1 \mid \mathbf{x}, \mathbf{z}) \neq \mathrm{P}(w=1 \mid \mathbf{x})$ and $\mathrm{P}(w=1 \mid \mathbf{x}, \mathbf{z})=G(\mathbf{x}, \mathbf{z} ; \boldsymbol{\gamma})$ is a known parametric form (usually probit or logit); and (d) $\operatorname{Var}\left(v_{0} \mid \mathbf{x}, \mathbf{z}\right)=\sigma_{0}^{2}$.

Part b assumes that $\mathrm{E}\left(v_{0} \mid \mathbf{x}\right)$ is linear in $\mathbf{x}$, and so it is more restrictive than Assumption ATEIV.1b. It does not usually hold for discrete response variables $y$, although it may be a reasonable approximation in some cases. Under parts a and b , the error $u_{0}$ in equation (21.61) has a zero conditional mean:\\
$\mathrm{E}\left(u_{0} \mid \mathbf{x}, \mathbf{z}\right)=0$.

Part d implies that $\operatorname{Var}\left(u_{0} \mid \mathbf{x}, \mathbf{z}\right)$ is constant. From the results on efficient choice of instruments in Section 14.4.3, the optimal IV for $w$ is $\mathrm{E}(w \mid \mathbf{x}, \mathbf{z})=G(\mathbf{x}, \mathbf{z} ; \boldsymbol{\gamma})$. Therefore, we can use a two-step IV method:

Procedure 21.1 (Under Assumption ATEIV.1'): (a) Estimate the binary response model $\mathrm{P}(w=1 \mid \mathbf{x}, \mathbf{z})=G(\mathbf{x}, \mathbf{z} ; \boldsymbol{\gamma})$ by maximum likelihood. Obtain the fitted probabilities, $\hat{G}_{i}$. The leading case occurs when $\mathrm{P}(w=1 \mid \mathbf{x}, \mathbf{z})$ follows a probit model.\\
(b) Estimate equation (21.61) by IV using instruments 1, $\hat{G}_{i}$, and $\mathbf{x}_{i}$.

There are several nice features of this IV estimator. First, it can be shown that the conditions sufficient to ignore the estimation of $y$ in the first stage hold; see Section 6.1.2. Therefore, the usual 2SLS standard errors and test statistics are asymptotically valid. Second, under Assumption ATEIV.1', the IV estimator from step $b$ is asymptotically efficient in the class of estimators where the IVs are functions of $\left(\mathbf{x}_{i}, \mathbf{z}_{i}\right)$; see Problem 8.11. If Assumption ATEIV.1'd does not hold, all statistics should be made robust to heteroskedasticity, and we no longer have the efficient IV estimator.

Procedure 21.1 has an important robustness property. Because we are using $\hat{G}_{i}$ as an instrument for $w_{i}$, the model for $\mathrm{P}(w=1 \mid \mathbf{x}, \mathbf{z})$ does not have to be correctly specified. For example, if we specify a probit model for $\mathrm{P}(w=1 \mid \mathbf{x}, \mathbf{z})$, we do not need the probit\\
model to be correct. Generally, what we need is that the linear projection of $w$ onto $\left[\mathbf{x}, G\left(\mathbf{x}, \mathbf{z} ; \boldsymbol{\gamma}^{*}\right)\right]$ actually depends on $G\left(\mathbf{x}, \mathbf{z} ; \boldsymbol{\gamma}^{*}\right)$, where we use $\boldsymbol{\gamma}^{*}$ to denote the plim of the maximum likelihood estimator when the model is misspecified (see White, 1982a, and Section 13.11.1). These requirements are fairly weak when $\mathbf{z}$ is partially correlated with $w$.

Technically, $\tau$ and $\boldsymbol{\beta}$ are identified even if we do not have extra exogenous variables excluded from $\mathbf{x}$. But we can rarely justify the estimator in this case. For concreteness, suppose that $w$ given $\mathbf{x}$ follows a probit model (and we have no $\mathbf{z}$, or $\mathbf{z}$ does not appear in $\mathrm{P}(w=1 \mid \mathbf{x}, \mathbf{z}))$. Because $G(\mathbf{x}, \boldsymbol{\gamma}) \equiv \Phi\left(\gamma_{0}+\mathbf{x} \gamma_{1}\right)$ is a nonlinear function of $\mathbf{x}$, it is not perfectly correlated with $\mathbf{x}$, so it can be used as an IV for $w$. This situation is very similar to the one discussed in Section 19.6.1: while identification holds for all values of $\alpha$ and $\boldsymbol{\beta}$ if $\boldsymbol{\gamma}_{1} \neq \mathbf{0}$, we are achieving identification off of the nonlinearity of $\mathrm{P}(w=1 \mid \mathbf{x})$. Further, $\Phi\left(\gamma_{0}+\mathbf{x} \gamma_{1}\right)$ and $\mathbf{x}$ are typically highly correlated. As we discussed in Section 5.2.6, severe multicollinearity among the IVs can result in very imprecise IV estimators. In fact, if $\mathrm{P}(w=1 \mid \mathbf{x})$ followed a linear probability model, $\tau$ would not be identified. See Problem 21.5 for an illustration.

Example 21.3 (Estimating the Effects of Education on Fertility): We use the data in FERTIL2.RAW to estimate the effect of attaining at least seven years of education on fertility. The data are for women of childbearing age in Botswana. Seven years of education is, by far, the modal amount of positive education. (About 21 percent of women report zero years of education. For the subsample with positive education, about 33 percent report seven years of education.) Let $y=$ children, the number of living children, and let $w=e d u c 7$ be a binary indicator for at least seven years of education. The elements of $\mathbf{x}$ are age, age ${ }^{2}$, evermarr (ever married), urban (lives in an urban area), electric (has electricity), and $t v$ (has a television).

The OLS estimate of $\tau$ is -.394 (se $=.050$ ). We also use the variable frsthalf, a binary variable equal to one if the woman was born in the first half of the year, as an IV for educ7. It is easily shown that educ7 and frsthalf are significantly negatively related. The usual IV estimate is much larger in magnitude than the OLS estimate, but only marginally significant: $-1.131(\mathrm{se}=.619)$. The estimate from Procedure 21.1 is even bigger in magnitude, and very significant: $-1.975(\mathrm{se}=.332)$. The standard error that is robust to arbitrary heteroskedasticity is even smaller. Therefore, using the probit fitted values as an IV, rather than the usual linear projection, produces a more precise estimate (and one notably larger in magnitude).

The IV estimate of education effect seems very large. One possible problem is that, because children is a nonnegative integer that piles up at zero, the assumptions underlying Procedure 21.1-namely, Assumptions ATEIV.1'a and ATEIV.1'b-\\
might not be met. We could instead apply methods for exponential response functions in Section 18.5. Both the Terza (1998) and Mullahy (1997) approaches can be derived using a counterfactual framework.

In principle, it is important to recognize that Procedure 21.1 is not the same as using $\hat{G}$ as a regressor in place of $w$. That is, IV estimation of equation (21.61) is not the same as the OLS estimator from\\
$y_{i}$ on 1, $\hat{G}_{i}, \mathbf{x}_{i}$

Consistency of the OLS estimators from regression (21.63) relies on having the model for $\mathrm{P}(w=1 \mid \mathbf{x}, \mathbf{z})$ correctly specified. If the first three parts of Assumption ATEIV.1' hold, then\\
$\mathrm{E}(y \mid \mathbf{x}, \mathbf{z})=\delta_{0}+\tau G(\mathbf{x}, \mathbf{z} ; \boldsymbol{\gamma})+\mathbf{x} \boldsymbol{\beta}$,\\
and, from the results on generated regressors in Chapter 6, the estimators from regression (21.63) are generally consistent. But Procedure 21.1 is more robust because it does not require Assumption ATEIV.1'c for consistency.

A different way to see the robustness of the IV approach compared with the regression approach is to think about the underlying first-stage regression in the population, which is the linear projection of $w$ on $\left[1, G\left(\mathbf{x}, \mathbf{z} ; \gamma^{*}\right), \mathbf{x}\right] ;$ write this as $\eta_{0}+\eta_{1} G\left(\mathbf{x}, \mathbf{z} ; \gamma^{*}\right)+\mathbf{x} \boldsymbol{\eta}_{2}$. The IV estimator is consistent for any values of the etas provided $\eta_{1} \neq 0$ because we simply need an exogenous variable that moves around $w$ that is not perfectly correlated with $\mathbf{x}$. By contrast, consistency of the two-step regression estimator for $\tau$ requires $\eta_{1}=1$. If $\mathrm{P}(w=1 \mid \mathbf{x}, \mathbf{z})=G\left(\mathbf{x}, \mathbf{z} ; \gamma^{*}\right)$, then $\eta_{0}=0$, $\eta_{1}=1$, and $\boldsymbol{\eta}_{2}=\mathbf{0}$, and so the linear projection of $y$ on $\left[1, G\left(\mathbf{x}, \mathbf{z} ; \gamma^{*}\right), \mathbf{x}\right]$ is $\delta_{0}+ \tau G\left(\mathbf{x}, \mathbf{z} ; \gamma^{*}\right)+\mathbf{x} \boldsymbol{\beta}$; thus, regression (21.63) consistently estimates all parameters. But if $G(\mathbf{x}, \mathbf{z} ; \gamma)$ is misspecified $\eta_{1}$ can differ from unity (and $\boldsymbol{\eta}_{2} \neq \mathbf{0}$ ). Generally, the regression (21.63) consistently estimates $\delta_{0}+\tau \eta_{0}, \tau \eta_{1}$, and $\boldsymbol{\beta}+\tau \boldsymbol{\eta}_{2}$ as the intercept, coefficient on $\hat{G}_{i}$, and coefficients on $\mathbf{x}_{i}$, respectively.

Another problem with regression (21.63) is that the usual OLS standard errors and test statistics are not valid, for two reasons. First, if $\operatorname{Var}\left(u_{0} \mid \mathbf{x}, \mathbf{z}\right)$ is constant, $\operatorname{Var}(y \mid \mathbf{x}, \mathbf{z})$ cannot be constant because $\operatorname{Var}(w \mid \mathbf{x}, \mathbf{z})$ is not constant. By itself this is a minor nuisance because heteroskedasticity-robust standard errors and test statistics are easy to obtain. (However, it does call into question the efficiency of the estimator from regression (21.63).) A more serious problem is that the asymptotic variance of the estimator from regression (21.63) depends on the asymptotic variance of $\hat{\gamma}$ unless $\alpha=0$, and the heteroskedasticity-robust standard errors do not correct for this.

In summary, using fitted probabilities from a first-stage binary response model, such as probit or logit, as an instrument for $w$ is a nice way to exploit the binary nature of the endogenous explanatory variable. In addition, the asymptotic inference is always standard. Using $\hat{G}_{i}$ as an instrument does require the assumption that $\mathrm{E}\left(v_{0} \mid \mathbf{x}, \mathbf{z}\right)$ depends only on $\mathbf{x}$ and is linear in $\mathbf{x}$, which can be more restrictive than Assumption ATEIV.1'b.

Allowing for the interaction $w\left(v_{1}-v_{0}\right)$ in equation (21.60) is notably harder. In general, when $v_{1} \neq v_{0}$, the IV estimator (using $\mathbf{z}$ or $\hat{G}$ as IVs for $w$ ) does not consistently estimate $\tau_{\text {ate }}$ (or $\tau_{\text {att }}$ ). Nevertheless, it is useful to find assumptions under which IV estimation does consistently estimate $A T E$. This problem has been studied by Angrist (1991), Heckman (1997), and Wooldridge (1997b, 2003b), and we synthesize results from these papers.

Under the conditional mean redundancy assumptions\\
$\mathrm{E}\left(v_{0} \mid \mathbf{x}, \mathbf{z}\right)=\mathrm{E}\left(v_{0} \mid \mathbf{x}\right) \quad$ and $\quad \mathrm{E}\left(v_{1} \mid \mathbf{x}, \mathbf{z}\right)=\mathrm{E}\left(v_{1} \mid \mathbf{x}\right)$,\\
we can always write equation (21.60) as\\
$y=\mu_{0}+\tau w+g_{0}(\mathbf{x})+w\left[g_{1}(\mathbf{x})-g_{0}(\mathbf{x})\right]+e_{0}+w\left(e_{1}-e_{0}\right)$,\\
where $\tau=\tau_{\text {ate }}$ and\\
$v_{0}=g_{0}(\mathbf{x})+e_{0}, \quad \mathrm{E}\left(e_{0} \mid \mathbf{x}, \mathbf{z}\right)=0$,\\
$v_{1}=g_{1}(\mathbf{x})+e_{1}, \quad \mathrm{E}\left(e_{1} \mid \mathbf{x}, \mathbf{z}\right)=0$.

Given functional form assumptions for $g_{0}$ and $g_{1}$-which would typically be linear in parameters-we can estimate equation (21.65) by IV, where the error term is $e_{0}+w\left(e_{1}-e_{0}\right)$. For concreteness, suppose that\\
$g_{0}(\mathbf{x})=\eta_{0}+\mathbf{x} \boldsymbol{\beta}_{0}, \quad g_{1}(\mathbf{x})-g_{0}(\mathbf{x})=(\mathbf{x}-\boldsymbol{\psi}) \boldsymbol{\delta}$,\\
where $\boldsymbol{\psi}=\mathrm{E}(\mathbf{x})$. If we plug these equations into equation (21.65), we need instruments for $w$ and $w(\mathbf{x}-\boldsymbol{\psi})$ (note that $\mathbf{x}$ does not contain a constant here). If $q \equiv q(\mathbf{x}, \mathbf{z})$ is the instrument for $w$ (such as the response probability in Procedure 21.1), the natural instrument for $w \cdot \mathbf{x}$ is $q \cdot \mathbf{x}$. (And, if $q$ is the efficient IV for $w, q \cdot \mathbf{x}$ is the efficient instrument for $w \cdot \mathbf{x}$.) When will applying IV to


\begin{equation*}
y=\gamma+\tau w+\mathbf{x} \boldsymbol{\beta}_{0}+w(\mathbf{x}-\boldsymbol{\psi}) \boldsymbol{\delta}+e_{0}+w\left(e_{1}-e_{0}\right) \tag{21.69}
\end{equation*}


be consistent? If the last term disappears, and, in particular, if\\
$e_{1}=e_{0}$,\\
then the error $e_{0}$ has zero mean given $(\mathbf{x}, \mathbf{z})$; this result means that IV estimation of equation (21.69) produces consistent, asymptotically normal estimators.

ASSUMPTION ATEIV.2: With $y$ expressed as in equation (21.60), conditions (21.64), (21.65), and (21.70) hold. In addition, Assumption ATEIV.1'c holds.

We have the following extension of Procedure 21.1:\\
Procedure 21.2 (Under Assumption ATEIV.2): (a) Same as Procedure 21.1.\\
(b) Estimate the equation\\
$y_{i}=\gamma+\tau w_{i}+\mathbf{x}_{i} \boldsymbol{\beta}_{0}+\left[w_{i}\left(\mathbf{x}_{i}-\overline{\mathbf{x}}\right)\right] \boldsymbol{\delta}+$ error $_{i}$\\
by IV, using instruments $1, \hat{G}_{i}, \mathbf{x}_{i}$ and $\hat{G}_{i}\left(\mathbf{x}_{i}-\overline{\mathbf{x}}\right)$.\\
If we add Assumption ATEIV.1'd, Procedure 21.2 produces the efficient IV estimator (when we ignore estimation of $\mathrm{E}(\mathbf{x})$ ). As with Procedure 21.1, we do not actually need the binary response model to be correctly specified for identification. As an alternative, we can use $\mathbf{z}_{i}$ and interactions between $\mathbf{z}_{i}$ and $\mathbf{x}_{i}$ as instruments, which generally results in testable overidentifying restrictions.

Technically, the fact that $\overline{\mathbf{x}}$ is an estimator of $\mathrm{E}(\mathbf{x})$ should be accounted for in computing the standard errors of the IV estimators. But, as shown in Problem 6.10, the adjustments for estimating $\mathrm{E}(\mathbf{x})$ often will have a trivial effect on the standard errors; in practice, we can just use the usual or heteroskedasticity-robust standard errors. Alternatively, we can apply the bootstrap.

Example 21.4 (An IV Approach to Evaluating Job Training): To evaluate the effects of a job training program on subsequent wages, suppose that $\mathbf{x}$ includes education, experience, and the square of experience. If $z$ indicates eligibility in the program, we would estimate the equation

$$
\begin{aligned}
\log (\text { wage })= & \mu_{0}+\tau \text { jobtrain }+\beta_{01} \text { educ }+\beta_{02} \text { exper }+\beta_{03} \text { exper }^{2} \\
& +\delta_{1} \text { jobtrain } \cdot(\text { educ }-\overline{\text { educ }})+\delta_{2} \text { jobtrain } \cdot(\text { exper }-\overline{\text { exper }}) \\
& +\delta_{3} \text { jobtrain } \cdot\left(\text { exper }^{2}-\overline{\text { exper }^{2}}\right)+\text { error }
\end{aligned}
$$

by IV, using instruments $1, z$, educ, exper, exper ${ }^{2}$, and interactions of $z$ with all demeaned covariates. Notice that for the last interaction, we subtract off the average of exper ${ }^{2}$. Alternatively, we could use in place of $z$ the fitted values from a probit of jobtrain on $(\mathbf{x}, z)$.

Procedure 21.2 is easy to carry out, but its consistency generally hinges on condition (21.70), not to mention the functional form assumptions in equation (21.68). We\\
can relax condition (21.70) to


\begin{equation*}
\mathrm{E}\left[w\left(e_{1}-e_{0}\right) \mid \mathbf{x}, \mathbf{z}\right]=\mathrm{E}\left[w\left(e_{1}-e_{0}\right)\right] \tag{21.72}
\end{equation*}


We do not need $w\left(e_{1}-e_{0}\right)$ to have zero mean, as a nonzero mean only affects the intercept. It is important to see that correlation between $w$ and ( $e_{1}-e_{0}$ ) does not invalidate the IV estimator of $\tau$ from Procedure 21.2. However, we must assume that the covariance conditional on ( $\mathbf{x}, \mathbf{z}$ ) is constant. Even if this assumption is not exactly true, it might be approximately true.

It is easy to see why, along with conditions (21.64) and (21.68), condition (21.72) implies consistency of the IV estimator. We can write equation (21.69) as\\
$y=\xi+\tau w+\mathbf{x} \boldsymbol{\beta}_{0}+w(\mathbf{x}-\boldsymbol{\psi}) \boldsymbol{\delta}+e_{0}+r$,\\
where $r=w\left(e_{1}-e_{0}\right)-\mathrm{E}\left[w\left(e_{1}-e_{0}\right)\right]$ and $\xi=\gamma+\mathrm{E}\left[w\left(e_{1}-e_{0}\right)\right]$. Under condition (21.72), $\mathrm{E}(r \mid \mathbf{x}, \mathbf{z})=0$, and so the composite error $e_{0}+r$ has zero mean conditional on $(\mathbf{x}, \mathbf{z})$. Therefore, any function of $(\mathbf{x}, \mathbf{z})$ can be used as instruments in equation (21.73). Under the following modification of Assumption ATEIV.2, Procedure 21.2 is still consistent:\\
assumption ATEIV.2': With $y$ expressed as in equation (21.60), conditions (21.64), (21.68), and (21.72) hold. In addition, Assumption ATEIV.1'c holds.

Even if Assumption ATEIV.1'd holds in addition to Assumption ATEIV.1'c, the IV estimator is generally not efficient because $\operatorname{Var}(r \mid \mathbf{x}, \mathbf{z})$ would typically be heteroskedastic.

Angrist (1991) provided primitive conditions for assumption (21.72) in the case where $\mathbf{z}$ is independent of $\left(y_{0}, y_{1}, \mathbf{x}\right)$. Then, the covariates can be dropped entirely from the analysis (leading to IV estimation of the simple regression equation $y= \xi+\tau w+$ error $)$. We can extend those conditions here to allow $\mathbf{z}$ and $\mathbf{x}$ to be correlated. Assume that


\begin{equation*}
\mathrm{E}\left(w \mid \mathbf{x}, \mathbf{z}, e_{1}-e_{0}\right)=h(\mathbf{x}, \mathbf{z})+k\left(e_{1}-e_{0}\right) \tag{21.74}
\end{equation*}


for some functions $h(\cdot)$ and $k(\cdot)$ and that\\
$e_{1}-e_{0}$ is independent of $(\mathbf{x}, \mathbf{z})$.

Under these two assumptions,


\begin{align*}
\mathrm{E}\left[w\left(e_{1}-e_{0}\right) \mid \mathbf{x}, \mathbf{z}\right] & =h(\mathbf{x}, \mathbf{z}) \mathrm{E}\left(e_{1}-e_{0} \mid \mathbf{x}, \mathbf{z}\right)+\mathrm{E}\left[\left(e_{1}-e_{0}\right) k\left(e_{1}-e_{0}\right) \mid \mathbf{x}, \mathbf{z}\right] \\
& =h(\mathbf{x}, \mathbf{z}) \cdot 0+\mathrm{E}\left[\left(e_{1}-e_{0}\right) k\left(e_{1}-e_{0}\right)\right] \\
& =\mathrm{E}\left[\left(e_{1}-e_{0}\right) k\left(e_{1}-e_{0}\right)\right], \tag{21.76}
\end{align*}


which is just an unconditional moment in the distribution of $e_{1}-e_{0}$. We have used the fact that $\mathrm{E}\left(e_{1}-e_{0} \mid \mathbf{x}, \mathbf{z}\right)=0$ and that any function of $e_{1}-e_{0}$ is independent of $(\mathbf{x}, \mathbf{z})$ under assumption (21.75). If we assume that $k(\cdot)$ is the identity function (as in Wooldridge, 1997b), then equation (21.76) is $\operatorname{Var}\left(e_{1}-e_{0}\right)$.

Assumption (21.75) is reasonable for continuously distributed responses, but it would not generally be reasonable when $y$ is a discrete response or corner solution outcome. Further, even if assumption (21.75) holds, assumption (21.74) is violated when $w$ given $\mathbf{x}, \mathbf{z}$, and ( $e_{1}-e_{0}$ ) follows a standard binary response model. For example, a probit model would have\\
$\mathrm{P}\left(w=1 \mid \mathbf{x}, \mathbf{z}, e_{1}-e_{0}\right)=\Phi\left[\pi_{0}+\mathbf{x} \boldsymbol{\pi}_{1}+\mathbf{z} \boldsymbol{\pi}_{2}+\rho\left(e_{1}-e_{0}\right)\right]$,\\
which is not separable in ( $\mathbf{x}, \mathbf{z}$ ) and ( $e_{1}-e_{0}$ ). Nevertheless, assumption (21.74) might be a reasonable approximation in some cases. Without covariates, Angrist (1991) presents simulation evidence that suggests the simple IV estimator does quite well for estimating the ATE even when assumption (21.74) is violated.

\subsection*{21.4.2 Correction and Control Function Approaches}
Rather than assuming that the presence of $w\left(e_{1}-e_{0}\right)$ in the error term does not cause inconsistency for IV estimators, we can use functional form and distributional assumptions to directly account for this term. The first approach we consider, proposed by Wooldridge (2008), involves adding a correction function, which is a function of the exogenous variables ( $\mathbf{x}, \mathbf{z}$ ), to equation (21.73), and then applying instrumental variables to account for the endogeneity of $w$ and $w(\mathbf{x}-\boldsymbol{\psi})$. To derive the correction function, we add to assumptions (21.75) and (21.77) a normality assumption. Let $c \equiv e_{1}-e_{0}$ and assume


\begin{equation*}
c \sim \operatorname{Normal}\left(0, \omega^{2}\right) \tag{21.78}
\end{equation*}


Under assumptions (21.75), (21.77), and (21.78) we can derive an estimating equation to show that $\tau_{\text {ate }}$ is usually identified.

To derive an estimating equation, note that conditions (21.75), (21.77), and (21.78) imply that\\
$\mathrm{P}(w=1 \mid \mathbf{x}, \mathbf{z})=\Phi\left(\theta_{0}+\mathbf{x} \boldsymbol{\theta}_{1}+\mathbf{z} \boldsymbol{\theta}_{2}\right)$,\\
where each theta is the corresponding pi multiplied by $\left[1+\rho^{2} \omega^{2}\right]^{-1 / 2}$. If we let $a$ denote the latent error underlying equation (21.79) (with a standard normal distribution), then conditions (21.75), (21.77), and (21.78) imply that ( $a, c$ ) has a zero-mean bivariate normal distribution that is independent of $(\mathbf{x}, \mathbf{z})$. Therefore, $\mathrm{E}(c \mid a, \mathbf{x}, \mathbf{z})= \mathrm{E}(c \mid a)=\xi a$ for some parameter $\xi$, and\\
$\mathrm{E}(w c \mid \mathbf{x}, \mathbf{z})=\mathrm{E}[w \mathrm{E}(c \mid a, \mathbf{x}, \mathbf{z}) \mid \mathbf{x}, \mathbf{z}]=\xi \mathrm{E}(w a \mid \mathbf{x}, \mathbf{z})$.\\
Using the fact that $a \sim \operatorname{Normal}(0,1)$ and is independent of $(\mathbf{x}, \mathbf{z})$, we have


\begin{align*}
\mathrm{E}(w a \mid \mathbf{x}, \mathbf{z}) & =\int_{-\infty}^{\infty} 1\left[\theta_{0}+\mathbf{x} \boldsymbol{\theta}_{1}+\mathbf{z} \boldsymbol{\theta}_{2}+a \geq 0\right] a \phi(a) \mathrm{d} a \\
& =\phi\left(-\left\{\theta_{0}+\mathbf{x} \boldsymbol{\theta}_{1}+\mathbf{z} \boldsymbol{\theta}_{2}\right\}\right)=\phi\left(\theta_{0}+\mathbf{x} \boldsymbol{\theta}_{1}+\mathbf{z} \boldsymbol{\theta}_{2}\right), \tag{21.80}
\end{align*}


where $\phi(\cdot)$ is the standard normal density. Therefore, we can now write\\
$y=\gamma+\tau w+\mathbf{x} \boldsymbol{\beta}+w(\mathbf{x}-\boldsymbol{\psi}) \boldsymbol{\delta}+\xi \phi\left(\theta_{0}+\mathbf{x} \boldsymbol{\theta}_{1}+\mathbf{z} \boldsymbol{\theta}_{2}\right)+e_{0}+r$,\\
where $r=w c-\mathrm{E}(w c \mid \mathbf{x}, \mathbf{z})$. The composite error in equation (21.81) has zero mean conditional on ( $\mathbf{x}, \mathbf{z}$ ), and so we can estimate the parameters using IV methods. One catch is the nonlinear function $\phi\left(\theta_{0}+\mathbf{x} \boldsymbol{\theta}_{1}+\mathbf{z} \boldsymbol{\theta}_{2}\right)$. We could use nonlinear two-stage least squares, as described in Chapter 14. But a two-step approach is easier. First, we gather together the assumptions:\\
assumption ATEIV.3: With $y$ written as in equation (21.70), maintain assumptions (21.64), (21.68), (21.75), (21.77), (with $\boldsymbol{\pi}_{2} \neq \mathbf{0}$ ), and (21.78).

Procedure 21.3 (Under Assumption ATEIV.3): (a) Estimate $\theta_{0}, \boldsymbol{\theta}_{1}$, and $\boldsymbol{\theta}_{2}$ from a probit of $w$ on ( $1, \mathbf{x}, \mathbf{z}$ ). Form the predicted probabilities, $\hat{\Phi}_{i}$, along with $\hat{\phi}_{i}= \phi\left(\hat{\theta}_{0}+\mathbf{x}_{i} \hat{\boldsymbol{\theta}}_{1}+\mathbf{z}_{i} \hat{\boldsymbol{\theta}}_{2}\right), i=1,2, \ldots, N$.\\
(b) Estimate the equation


\begin{equation*}
y_{i}=\gamma+\tau w_{i}+\mathbf{x}_{i} \boldsymbol{\beta}_{0}+w_{i}\left(\mathbf{x}_{i}-\overline{\mathbf{x}}\right) \boldsymbol{\delta}+\xi \hat{\phi}_{i}+\operatorname{error}_{i} \tag{21.82}
\end{equation*}


by IV, using instruments $\left[1, \hat{\Phi}_{i}, \mathbf{x}_{i}, \hat{\Phi}_{i}\left(\mathbf{x}_{i}-\overline{\mathbf{x}}\right), \hat{\phi}_{i}\right]$.\\
Wooldridge (2008) calls the extra term $\hat{\phi}_{i} \equiv \phi\left(\hat{\theta}_{0}+\mathbf{x}_{i} \hat{\boldsymbol{\theta}}_{1}+\mathbf{z}_{i} \hat{\boldsymbol{\theta}}_{2}\right)$ a "correction function" to distinguish it from the more common "control function," which we turn to shortly. Unlike with the control function approach, adding $\phi\left(\theta_{0}+\mathbf{x}_{i} \boldsymbol{\theta}_{1}+\mathbf{z}_{i} \boldsymbol{\theta}_{2}\right)$ does not render $w$ or $w(\mathbf{x}-\boldsymbol{\psi})$ exogenous in equation (21.81). Rather, it ensures that the composite error, $e_{0}+r$, is mean independent of $(\mathbf{x}, \mathbf{z})$ when $w\left(e_{1}-e_{0}\right)$ is present and not assumed independent of ( $\mathbf{x}, \mathbf{z}$ ). Conveniently, IV estimation of equation (21.82) leads to a simple test of $H_{0}: \xi=0$. Under the null, the coefficient on the generated regressor, $\hat{\phi}_{i}$ is zero. Further, that the IVs are estimated in a first stage does not affect the $\sqrt{N}$-asymptotic distribution of the IV estimators; see Section 6.1.3. Therefore, if we ignore the estimation error in $\overline{\mathbf{x}}$, we can use a standard hetero-skedasticity-robust $t$ statistic on $\hat{\phi}_{i}$ to test whether the correction function is needed. (Note that $r=w c-\mathrm{E}(w c \mid \mathbf{x}, \mathbf{z})$ can be homoskedastic under the null but there is\\
never any harm in making the test robust to heteroskedasticity.) Notice that the only place we used normality of ( $c, a$ ) is in deriving the correction function. Because $\xi=0$ under the null, this normality assumption is not needed for the test. (Remember, our use of probit-fitted values as instruments for $w$ does not hinge on the probit model for $w$ being true; it motivates the choice of IVs, but the probit model can be arbitrarily misspecified.)

If we allow the possibility that $\xi \neq 0$, all standard errors need to be adjusted for the two-step estimation. One can use the delta method or express the two-step estimation as a generalized method of moments problem as in Chapter 14. (Either approach can ignore or account for sampling error in $\overline{\mathbf{x}}$.) Because the two-step estimation is computationally simple, the bootstrap is attractive as a way to account for all sampling error in $\hat{\tau}$.

Even if $\xi \neq 0$, adding the correction function $\hat{\phi}_{i}$ to the equation need not have much effect on the estimate of $\tau$. To see why, consider the version of equation (21.81) when we have no covariates $\mathbf{x}$ and with a scalar IV, $z$ :


\begin{equation*}
y=\gamma+\tau w+\xi \phi\left(\theta_{0}+\theta_{1} z\right)+u, \quad \mathrm{E}(u \mid z)=0 \tag{21.83}
\end{equation*}


This equation holds, for example, if the instrument $z$ is independent of $\left(v_{0}, v_{1}\right)$. The simple IV estimator of $\tau$ is obtained by omitting $\phi\left(\theta_{0}+\theta_{1} z\right)$. If we use $z$ as an IV for $w$, the simple IV estimator is consistent provided $z$ and $\phi\left(\theta_{0}+\theta_{1} z\right)$ are uncorrelated. (Remember, having an omitted variable that is uncorrelated with the IV does not cause inconsistency of the IV estimator.) Even though $\phi\left(\theta_{0}+\theta_{1} z\right)$ is a function of $z$, these two variables might have small correlation because $z$ is monotonic while $\phi\left(\theta_{0}+\theta_{1} z\right)$ is symmetric about $-\left(\theta_{0} / \theta_{1}\right)$. This discussion shows that condition (21.72) is not necessary for IV to consistently estimate the ATE: It could be that while $\mathrm{E}\left[w\left(e_{1}-e_{0}\right) \mid \mathbf{x}, \mathbf{z}\right]$ is not constant, it is roughly uncorrelated with $\mathbf{x}$ (or the functions of $\mathbf{x}$ ) that appear in equation (21.73), as well as with the functions of $\mathbf{z}$ used as instruments.

Equation (21.83) illustrates another important point: If $\xi \neq 0$ and the single instrument $z$ is binary, $\tau$ is not identified. Lack of identification occurs because $\phi\left(\theta_{0}+\theta_{1} z\right)$ takes on only two values, which means it is perfectly linearly related to $z$. So long as $z$ takes on more than two values, $\tau$ is generally identified, although the identification is due to the fact that $\phi(\cdot)$ is a different nonlinear function than $\Phi(\cdot)$. With $\mathbf{x}$ in the model $\hat{\phi}_{i}$ and $\hat{\Phi}_{i}$ might be collinear, resulting in imprecise IV estimates.

Because $r$ in equation (21.81) is heteroskedastic, the instruments below equation (21.82) are not optimal, and so we might simply use $\mathbf{z}_{i}$ along with interactions of $\mathbf{z}_{i}$ with $\left(\mathbf{x}_{i}-\overline{\mathbf{x}}\right)$ and $\hat{\phi}_{i}$ as IVs. If $\mathbf{z}_{i}$ has dimension greater than one, then we can test the overidentifying restrictions as a partial test of instrument selection and the normality\\[0pt]
assumptions. Of course, we could use the results of Chapter 14 to characterize and estimate the optimal instruments, but this approach is fairly involved [see, for example, Newey and McFadden (1994)].

We can use a similar set of assumptions to derive a control function (CF) approach for estimating $\tau=\tau_{\text {ate }}$. Recall that the CF approach involves finding $\mathrm{E}(y \mid w, \mathbf{x}, \mathbf{z})$ and then using regression methods. (Or, a maximum likelihood approach is often available under a stronger set of assumptions.) Typically, the CF approach is derived in the context of the endogenous switching regression model, but it is easily seen that the treatment effect model with heterogeneous treatment, where the treatment is correlated with unobservables even after conditioning on observables, fits that bill when the treatment is defined to be the binary switching variable. In particular, equation (21.69) results by writing $y=(1-w) y_{0}+w y_{1}$ and imposing the linear, additive structures on $y_{0}$ and $y_{1}$ :

$$
\begin{aligned}
y & =(1-w)\left(\alpha_{0}+\mathbf{x} \boldsymbol{\beta}_{0}+e_{0}\right)+w\left(\alpha_{1}+\mathbf{x} \boldsymbol{\beta}_{1}+e_{1}\right) \\
& =\alpha_{0}+\mathbf{x} \boldsymbol{\beta}_{0}+\left(\alpha_{1}-\alpha_{0}\right) w+w \mathbf{x}\left(\boldsymbol{\beta}_{1}-\boldsymbol{\beta}_{0}\right)+e_{0}+w\left(e_{1}-e_{0}\right) \\
& \equiv \gamma+\mathbf{x} \boldsymbol{\beta}_{0}+\tau w+w(\mathbf{x}-\boldsymbol{\psi}) \boldsymbol{\delta}+e_{0}+w c,
\end{aligned}
$$

where $\tau=\tau_{\text {ate }}=\mathrm{E}\left(y_{1}-y_{0}\right)$ and $\boldsymbol{\psi}=\mathrm{E}(\mathbf{x})$. This particular way of expressing $y$ in terms of the treatment, covariates, and unobservables emphasizes that we are primarily interested in $\tau$, although $\tau_{\text {ate }}(\mathbf{x})=\tau+(\mathbf{x}-\boldsymbol{\psi}) \boldsymbol{\delta}$ is of interest for studying how the average treatment effect changes as a function of observables.

We can derive a control function method under the following assumption.\\
ASSUMPTION ATEIV.4: With $y$ written as in equation (21.70), maintain assumptions (21.64) and (21.68). Furthermore, the treatment can be written as $w=1\left[\theta_{0}+\mathbf{x} \boldsymbol{\theta}_{1}+\right. \left.\mathbf{z} \boldsymbol{\theta}_{2}+a \geq 0\right]$, where ( $a, e_{0}, e_{1}$ ) is independent of ( $\mathbf{x}, \mathbf{z}$ ) with a trivariate normal distribution; in particular, $a \sim \operatorname{Normal}(0,1)$.

Under Assumption ATEIV.4, we can use calculations very similar to those used in Section 19.6.1 to obtain $\mathrm{E}(y \mid w, \mathbf{x}, \mathbf{z})$. In particular,


\begin{align*}
\mathrm{E}(y \mid w, \mathbf{x}, \mathbf{z})= & \gamma+\alpha w+\mathbf{x} \boldsymbol{\beta}_{0}+w(\mathbf{x}-\boldsymbol{\psi}) \boldsymbol{\delta}+\rho_{1} w[\phi(\mathbf{q} \boldsymbol{\theta}) / \Phi(\mathbf{q} \boldsymbol{\theta})] \\
& +\rho_{2}(1-w)\{\phi(\mathbf{q} \boldsymbol{\theta}) /[1-\Phi(\mathbf{q} \boldsymbol{\theta})]\} \tag{21.84}
\end{align*}


where $\mathbf{q} \boldsymbol{\theta} \equiv \theta_{0}+\mathbf{x} \boldsymbol{\theta}_{1}+\mathbf{z} \boldsymbol{\theta}_{2}$ and $\rho_{1}$ and $\rho_{2}$ are additional parameters. Heckman (1978) used this expectation to obtain two-step estimators of the switching regression model. (See Vella and Verbeek (1999) for a recent discussion of the switching regression model in the context of treatment effects.) Not surprisingly, equation (21.84) suggests\\
a simple two-step procedure, where the first step is identical to that in Procedure 21.3:

Procedure 21.4 (Under Assumption ATEIV.4): (a) Estimate $\theta_{0}, \boldsymbol{\theta}_{1}$, and $\boldsymbol{\theta}_{2}$ from a probit of $w$ on $(1, \mathbf{x}, \mathbf{z})$. Form the predicted probabilities, $\hat{\Phi}_{i}$, along with $\hat{\phi}_{i}=\phi\left(\hat{\theta}_{0}+\right. \left.\mathbf{x}_{i} \hat{\boldsymbol{\theta}}_{1}+\mathbf{z}_{i} \hat{\boldsymbol{\theta}}_{2}\right), i=1,2, \ldots, N$.\\
(b) Run the OLS regression\\
$y_{i}$ on $1, w_{i}, \mathbf{x}_{i}, w_{i}\left(\mathbf{x}_{i}-\overline{\mathbf{x}}\right), w_{i}\left(\hat{\phi}_{i} / \hat{\Phi}_{i}\right),\left(1-w_{i}\right)\left[\hat{\phi}_{i} /\left(1-\hat{\Phi}_{i}\right)\right]$\\
using all of the observations. The coefficient on $w_{i}$ is a consistent estimator of $\alpha$, the Ate.

When we restrict attention to the $w_{i}=1$ subsample, thereby dropping $w_{i}$ and $w_{i}\left(\mathbf{x}_{i}-\overline{\mathbf{x}}\right)$, we obtain the sample selection correction from Section 19.6.1. (The treatment $w_{i}$ becomes the sample selection indicator.) But the goal of sample selection corrections is very different from estimating an average treatment effect. For the sample selection problem, the goal is to estimate $\boldsymbol{\beta}_{0}$, which indexes $\mathrm{E}(y \mid \mathbf{x})$ in the population. By contrast, in estimating an ATE we are interested in the causal effect that $w$ has on $y$.

It makes sense to check for joint significance of the last two regressors in regression (21.85) as a test of endogeneity of $w$. Because the coefficients $\rho_{1}$ and $\rho_{2}$ are zero under $\mathrm{H}_{0}$, we can use the results from Chapter 6 to justify the usual Wald test (perhaps made robust to heteroskedasticity). If these terms are jointly insignificant at a sufficiently high level, we can justify the usual OLS regression without unobserved heterogeneity. If we reject $\mathrm{H}_{0}$, we must deal with the generated regressors problem in obtaining a valid standard error for $\hat{\alpha}$.

Technically, Procedure 21.3 is more robust than Procedure 21.4 because the former does not require a trivariate normality assumption. Linear conditional expectations, along with the assumption that $w$ given $(\mathbf{x}, \mathbf{z})$ follows a probit, suffice. In addition, Procedure 21.3 allows us to separate the issues of endogeneity of $w$ and nonconstant treatment effect.

Practically, the extra assumption in Procedure 21.4 is that $e_{0}$ is independent of $(\mathbf{x}, \mathbf{z})$ with a normal distribution. We may be willing to make this assumption, especially if the estimates from Procedure 21.3 are too imprecise to be useful. The efficiency issue is a difficult one because of the two-step estimation involved, but, intuitively, Procedure 21.4 is likely to be more efficient because it is based on $\mathrm{E}(y \mid w, \mathbf{x}, \mathbf{z})$. Procedure 21.3 involves replacing the unobserved composite error with its expectation conditional only on $(\mathbf{x}, \mathbf{z})$. In at least one case, Procedure 21.4 gives results when

Procedure 21.3 cannot: when $\mathbf{x}$ is not in the equation and there is a single binary instrument.

So far we have focused on estimating $\tau_{\text {ate }}$. Under a variant of Assumption ATEIV.2, we can consistently estimate $\tau_{\text {att }}$ by IV. As before, we express $y$ as in equation (21.60). First, we show how to consistently estimate $\tau_{a t t}(\mathbf{x})$, which can be written as\\
$\tau_{a t t}(\mathbf{x})=\mathrm{E}\left(y_{1}-y_{0} \mid \mathbf{x}, w=1\right)=\left(\mu_{1}-\mu_{0}\right)+\mathrm{E}\left(v_{1}-v_{0} \mid \mathbf{x}, w=1\right)$.\\
The following assumption identifies $\tau_{\text {att }}(\mathbf{x})$ :\\
ASSUMPTION ATTIV.1: (a) With $y$ expressed as in equation (21.60), the first part of assumption (21.64) holds, that is, $\mathrm{E}\left(v_{0} \mid \mathbf{x}, \mathbf{z}\right)=\mathrm{E}\left(v_{0} \mid \mathbf{x}\right)$; (b) $\mathrm{E}\left(v_{1}-v_{0} \mid \mathbf{x}, \mathbf{z}, w=1\right)= \mathrm{E}\left(v_{1}-v_{0} \mid \mathbf{x}, w=1\right)$; and (c) Assumption ATEIV.1'c holds.

We discussed part a of this assumption earlier, as it also appears in Assumption ATEIV.2'. It can be violated if agents change their behavior based on $\mathbf{z}$. Part $b$ deserves some discussion. Recall that $v_{1}-v_{0}$ is the person-specific gain from participation or treatment. Assumption ATTIV. 1 requires that for those in the treatment group, the gain is not predictable given $\mathbf{z}$, once $\mathbf{x}$ is controlled for. Heckman (1997) discusses Angrist's (1990) draft lottery example, where $z$ (a scalar) is draft lottery number. Men who had a large $z$ were virtually certain to escape the draft. But some men with large draft numbers chose to serve anyway. Even with good controls in $\mathbf{x}$, it seems plausible that, for those who chose to serve, a higher $z$ is associated with a higher gain to military service. In other words, for those who chose to serve, $v_{1}-v_{0}$ and $z$ are positively correlated, even after controlling for $\mathbf{x}$. This argument directly applies to estimation of $\tau_{\text {att }}$; the effect on estimation of $\tau_{\text {ate }}$ is less clear.

Assumption ATTIV.1b is plausible when $z$ is a binary indicator for eligibility in a program, which is randomly determined and does not induce changes in behavior other than whether or not to participate.

To see how Assumption ATTIV. 1 identifies $\tau_{\text {att }}(\mathbf{x})$, rewrite equation (21.60) as


\begin{align*}
y= & \mu_{0}+g_{0}(\mathbf{x})+w\left[\left(\mu_{1}-\mu_{0}\right)+\mathrm{E}\left(v_{1}-v_{0} \mid \mathbf{x}, w=1\right)\right] \\
& +w\left[\left(v_{1}-v_{0}\right)-\mathrm{E}\left(v_{1}-v_{0} \mid \mathbf{x}, w=1\right)\right]+e_{0} \\
= & \mu_{0}+g_{0}(\mathbf{x})+w \cdot \tau_{a t t}(\mathbf{x})+a+e_{0} \tag{21.86}
\end{align*}


where $a \equiv w\left[\left(v_{1}-v_{0}\right)-\mathrm{E}\left(v_{1}-v_{0} \mid \mathbf{x}, w=1\right)\right]$ and $e_{0}$ is defined in equation (21.66). Under Assumption ATTIV.1a, $\mathrm{E}\left(e_{0} \mid \mathbf{x}, \mathbf{z}\right)=0$. The hard part is dealing with the term $a$. When $w=0, a=0$. Therefore, to show that $\mathrm{E}(a \mid \mathbf{x}, \mathbf{z})=0$, it suffices to show\\
that $\mathrm{E}(a \mid \mathbf{x}, \mathbf{z}, w=1)=0$. (Remember, $\mathrm{E}(a \mid \mathbf{x}, \mathbf{z})=\mathrm{P}(w=0) \cdot \mathrm{E}(a \mid \mathbf{x}, \mathbf{z}, w=0)+ \mathrm{P}(w=1) \cdot \mathrm{E}(a \mid \mathbf{x}, \mathbf{z}, w=1)$.) But this result follows under Assumption ATTIV.1b:\\
$\mathrm{E}(a \mid \mathbf{x}, \mathbf{z}, w=1)=\mathrm{E}\left(v_{1}-v_{0} \mid \mathbf{x}, \mathbf{z}, w=1\right)-\mathrm{E}\left(v_{1}-v_{0} \mid \mathbf{x}, w=1\right)=0$.\\
Now, letting $r \equiv a+e_{0}$ and assuming that $g_{0}(\mathbf{x})=\eta_{0}+\mathbf{h}(\mathbf{x}) \boldsymbol{\beta}_{0}$ and $\tau_{\text {att }}(\mathbf{x})=\alpha+ \mathbf{f}(\mathbf{x}) \boldsymbol{\delta}$ for some row vector of functions $\mathbf{h}(\mathbf{x})$ and $\mathbf{f}(\mathbf{x})$, we can write\\
$y=\gamma_{0}+\mathbf{h}_{0}(\mathbf{x}) \boldsymbol{\beta}_{0}+\alpha w+[w \cdot \mathbf{f}(\mathbf{x})] \boldsymbol{\delta}+r, \quad \mathrm{E}(r \mid \mathbf{x}, \mathbf{z})=0$.\\
All the parameters of this equation can be consistently estimated by IV, using any functions of $(\mathbf{x}, \mathbf{z})$ as IVs. (These would include include $1, \mathbf{h}_{0}(\mathbf{x}), G(\mathbf{x}, \mathbf{z} ; \hat{\gamma})$-the fitted treatment probabilities-and $G(\mathbf{x}, \mathbf{z} ; \hat{\gamma}) \cdot \mathbf{f}(\mathbf{x})$.) The average treatment effect on the treated for any $\mathbf{x}$ is estimated as $\hat{\alpha}+\mathbf{f}(\mathbf{x}) \hat{\boldsymbol{\delta}}$. Averaging over the observations with $w_{i}=1$ gives a consistent estimator of $\tau_{\text {att }}$.

\subsection*{21.4.3 Estimating the Local Average Treatment Effect by IV}
We now discuss estimation of an evaluation parameter introduced by Imbens and Angrist (1994), the local average treatment effect (LATE), in the simplest possible setting. This requires a slightly more complicated notation. (More general cases require even more complicated notation, as in AIR.) As before, we let $w$ be the observed treatment indicator (taking on zero or one), and let the counterfactual outcomes be $y_{1}$ with treatment and $y_{0}$ without treatment. The observed outcome $y$ can be written as in equation (21.3).

To define $\tau_{\text {late }}$, we need to have an instrumental variable, $z$. In the simplest case $z$ is a binary variable, and we focus attention on that case here. For each unit $i$ in a random draw from the population, $z_{i}$ is zero or one. Associated with the two possible outcomes on $z$ are counterfactual treatments, $w_{0}$ and $w_{1}$. These are the treatment statuses we would observe if $z=0$ and $z=1$, respectively. For each unit, we observe only one of these. For example, $z$ can denote whether a person is eligible for a particular program, while $w$ denotes actual participation in the program.

Write the observed treatment status as\\
$w=(1-z) w_{0}+z w_{1}=w_{0}+z\left(w_{1}-w_{0}\right)$.

When we plug this equation into $y=y_{0}+w\left(y_{1}-y_{0}\right)$ we get\\
$y=y_{0}+w_{0}\left(y_{1}-y_{0}\right)+z\left(w_{1}-w_{0}\right)\left(y_{1}-y_{0}\right)$.\\
A key assumption is\\
$z$ is independent of $\left(y_{0}, y_{1}, w_{0}, w_{1}\right)$.

Under assumption (21.88), all expectations involving functions of $\left(y_{0}, y_{1}, w_{0}, w_{1}\right)$, conditional on $z$, do not depend on $z$. Therefore,\\
$\mathrm{E}(y \mid z=1)=\mathrm{E}\left(y_{0}\right)+\mathrm{E}\left[w_{0}\left(y_{1}-y_{0}\right)\right]+\mathrm{E}\left[\left(w_{1}-w_{0}\right)\left(y_{1}-y_{0}\right)\right]$,\\
and\\
$\mathrm{E}(y \mid z=0)=\mathrm{E}\left(y_{0}\right)+\mathrm{E}\left[w_{0}\left(y_{1}-y_{0}\right)\right]$.\\
Subtracting the second equation from the first gives\\
$\mathrm{E}(y \mid z=1)-\mathrm{E}(y \mid z=0)=\mathrm{E}\left[\left(w_{1}-w_{0}\right)\left(y_{1}-y_{0}\right)\right]$,\\
which can be written (see equation (2.49)) as

$$
\begin{aligned}
1 \cdot \mathrm{E}\left(y_{1}\right. & \left.-y_{0} \mid w_{1}-w_{0}=1\right) \mathrm{P}\left(w_{1}-w_{0}=1\right) \\
& +(-1) \mathrm{E}\left(y_{1}-y_{0} \mid w_{1}-w_{0}=-1\right) \mathrm{P}\left(w_{1}-w_{0}=-1\right) \\
& +0 \cdot \mathrm{E}\left(y_{1}-y_{0} \mid w_{1}-w_{0}=0\right) \mathrm{P}\left(w_{1}-w_{0}=0\right) \\
= & \mathrm{E}\left(y_{1}-y_{0} \mid w_{1}-w_{0}=1\right) \mathrm{P}\left(w_{1}-w_{0}=1\right) \\
& -\mathrm{E}\left(y_{1}-y_{0} \mid w_{1}-w_{0}=-1\right) \mathrm{P}\left(w_{1}-w_{0}=-1\right)
\end{aligned}
$$

To get further, we introduce another important assumption, called monotonicity by Imbens and Angrist:\\
$w_{1} \geq w_{0}$.

In other words, we are ruling out $w_{1}=0$ and $w_{0}=1$. This assumption has a simple interpretation when $z$ is a dummy variable representing eligibility for treatment: anyone in the population who would be in the treatment group in the absence of eligibility would be in the treatment group if eligible for treatment group. Units of the population who do not satisfy monotonicity are called defiers. In many applications, this assumption seems very reasonable. For example, if $z$ denotes randomly assigned eligibility in a job training program, assumption (21.90) simply requires that people who would participate without being eligible would also participate if eligible.

Under assumption (21.90), $\mathrm{P}\left(w_{1}-w_{0}=-1\right)=0$, so assumptions (21.88) and (21.90) imply\\
$\mathrm{E}(y \mid z=1)-\mathrm{E}(y \mid z=0)=\mathrm{E}\left(y_{1}-y_{0} \mid w_{1}-w_{0}=1\right) \mathrm{P}\left(w_{1}-w_{0}=1\right)$.

In this setup, Imbens and Angrist (1994) define $\tau_{\text {late }}$ to be\\
$\tau_{\text {late }}=\mathrm{E}\left(y_{1}-y_{0} \mid w_{1}-w_{0}=1\right)$.

Because $w_{1}-w_{0}=1$ is equivalent to $w_{1}=1, w_{0}=0, \tau_{\text {late }}$ has the following interpretation: it is the average treatment effect for those who would be induced to participate by changing $z$ from zero to one. There are two things about $\tau_{\text {late }}$ that make it different from the other treatment parameters. First, it depends on the instrument, $z$. If we use a different instrument, then $\tau_{\text {late }}$ generally changes. The parameters $\tau_{\text {ate }}$ and $\tau_{\text {att }}$ are defined without reference to an IV, but only with reference to a population. Second, because we cannot observe both $w_{1}$ and $w_{0}$, we cannot identify the subpopulation with $w_{1}-w_{0}=1$. By contrast, $\tau_{\text {ate }}$ averages over the entire population, while $\tau_{\text {att }}$ is the average for those who are actually treated.

Example 21.5 (LATE for Attending a Catholic High School): Suppose that $y$ is a standardized test score, $w$ is an indicator for attending a Catholic high school, and $z$ is an indicator for whether the student is Catholic. Then, generally, $\tau_{\text {late }}$ is the mean effect on test scores for those individuals who choose a Catholic high school because they are Catholic. Evans and Schwab (1995) use a high school graduation indicator for $y$, and they estimate a probit model with an endogenous binary explanatory variable, as described in Section 15.7.3. Under the probit assumptions, it is possible to estimate $\tau_{\text {ate }}$, whereas the simple IV estimator identifies $\tau_{\text {late }}$ under weaker assumptions.

Because $\mathrm{E}(y \mid z=1)$ and $\mathrm{E}(y \mid z=0)$ are easily estimated using a random sample, $\tau_{\text {late }}$ is identified if $\mathrm{P}\left(w_{1}-w_{0}=1\right)$ is estimable and nonzero. Importantly, from the monotonicity assumption, $w_{1}-w_{0}$ is a binary variable because $\mathrm{P}\left(w_{1}-w_{0}=-1\right)=0$. Therefore,

$$
\begin{aligned}
\mathrm{P}\left(w_{1}-w_{0}=1\right) & =\mathrm{E}\left(w_{1}-w_{0}\right)=\mathrm{E}\left(w_{1}\right)-\mathrm{E}\left(w_{0}\right)=\mathrm{E}(w \mid z=1)-\mathrm{E}(w \mid z=0) \\
& =\mathrm{P}(w=1 \mid z=1)-\mathrm{P}(w=1 \mid z=0)
\end{aligned}
$$

where the second-to-last equality follows from equations (21.87) and (21.88). Each conditional probability can be consistently estimated given a random sample on $(w, z)$. Therefore, the final assumption is\\
$\mathrm{P}(w=1 \mid z=1) \neq \mathrm{P}(w=1 \mid z=0)$.

To summarize, under assumptions (21.88), (21.90), and (21.93),\\
$\tau_{\text {late }}=[\mathrm{E}(y \mid z=1)-\mathrm{E}(y \mid z=0)] /[\mathrm{P}(w=1 \mid z=1)-\mathrm{P}(w=1 \mid z=0)]$.

Therefore, a consistent estimator is $\hat{\tau}_{\text {late }}=\left(\bar{y}_{1}-\bar{y}_{0}\right) /\left(\bar{w}_{1}-\bar{w}_{0}\right)$, where $\bar{y}_{1}$ is the sample average of $y_{i}$ over that part of the sample where $z_{i}=1$ and $\bar{y}_{0}$ is the sample average over $z_{i}=0$, and similarly for $\bar{w}_{1}$ and $\bar{w}_{0}$ (which are sample proportions).

From Problem 5.13b, we know that $\hat{\tau}_{\text {late }}$ is identical to the IV estimator of $\tau$ in the simple equation $y=\delta_{0}+\tau w+$ error , where $z$ is the IV for $w$.

Our conclusion is that, in the simple case of a binary instrument for the binary treatment, the usual IV estimator consistently estimates $\tau_{\text {late }}$ under weak assumptions. See Angrist, Imbens, and Rubin (1996) and the discussants' comments for much more, and Imbens and Wooldridge (2009) for a survey of extensions to the basic framework.

\subsection*{21.5 Regression Discontinuity Designs}
We now turn to estimating treatment effects in a special situation where certain institutional or logical structures act to determine, or at least affect, treatment assignment. Regression discontinuity (RD) designs have a long history. Important work in bringing the methods into the mainstream in economics includes van der Klaauw (2002) and Hahn, Todd, and van der Klaauw (2001) (or HTV (2001)). Imbens and Lemieux (2008) provide a nice overview and survey of the subject, and this section draws on this work, as well as Imbens and Wooldridge (2009).

Generally, RD designs exploit discontinuities in policy assignment. For example, there might be an age threshold at which one becomes eligible for pension plan vesting, or an income threshold at which one becomes eligible for financial aid. To exploit discontinuities that are often determined by rather ad hoc institutional structures, one assumes that units just on different sides of the discontinuity are essentially the same in unobservables that affect the relevant outcome. The treatment statuses of the two groups differ, say, because of the institutional setup, in which case differences in outcomes can be attributed to the different treatment statuses.

We consider the sharp design RD-where assignment follows a deterministic rule-and the fuzzy design, where the probability of being treated is discontinuous at a known point.

\subsection*{21.5.1 The Sharp Regression Discontinuity Design}
As before, let $y_{0}$ and $y_{1}$ denote the counterfactual outcomes without and with treatment. For a random draw $i$, these are denoted $y_{i 0}, y_{i 1}$. For now, assume there is a single covariate, $x_{i}$, determining treatment (sometimes called the forcing variable). In the sharp regression discontinuity (SRD) design case, treatment is determined as


\begin{equation*}
w_{i}=1\left[x_{i} \geq c\right] . \tag{21.95}
\end{equation*}


Along with the forcing variable $x_{i}$ and treatment status $w_{i}$, we observe $y_{i}= \left(1-w_{i}\right) y_{i 0}+w_{i} y_{i 1}$.

Define, as before, the counterfactual conditional means as $\mu_{g}(x)=\mathrm{E}\left(y_{g} \mid x\right)$, $g=0,1$. A critical assumption in RD designs is that these underlying mean functions are continuous. (Technically, they are continuous at $x=c$, but it is hard to imagine how we could ensure that they are without assuming continuity over the range of $x$.) Now, because $w$ is a deterministic function of $x$, ignorability necessarily holds. Stated in terms of conditional means,\\
$\mathrm{E}\left(y_{g} \mid x, w\right)=\mathrm{E}\left(y_{g} \mid x\right), \quad g=0,1$,\\
which is exactly Assumption ATE.1'. So, if ignorability holds, why can we not just apply previous methods, such as propensity score weighting? The key is that with a sharp RD design, the overlap assumption fails absolutely: by construction, $p(x)=0$ for all $x<c$ and $p(x)=1$ for $x \geq c$. Clearly we cannot use a method such as propensity score weighting.

Technically, we can use regression adjustment, but we would have to rely on extreme forms of extrapolation in using parametric models. For example, if we estimate $\tau_{\text {ate }}$ using general regression adjustment of the form (21.30)-where $m_{1}\left(x, \boldsymbol{\delta}_{1}\right)= \mathrm{E}(y \mid x, w=1)$ and $m_{0}\left(x, \boldsymbol{\delta}_{0}\right)=\mathrm{E}(y \mid x, w=0)$-then, say, for control observations we must compute $m_{1}\left(x_{i}, \hat{\boldsymbol{\delta}}_{1}\right)$ for $x_{i}<c$, even though no data points with $x_{i}<c$ were used in obtaining $\hat{\boldsymbol{\delta}}_{1}$. If we use local smoothing in a nonparametric setting, we could not convincingly estimate $m_{1}(x)$ for $x<c$ or $m_{0}(x)$ for $x \geq c$.

Rather than trying to estimate $\tau_{\text {ate }}$, which relies on strong functional form assumptions unless we just assume a constant treatment effect, the RD literature often settles for estimating the average treatment effect at the discontinuity point, defined as\\
$\tau_{c} \equiv \mathrm{E}\left(y_{1}-y_{0} \mid x=c\right)=\mu_{1}(c)-\mu_{0}(c)$.

This is the average treatment effect for those at the margin of receiving the treatment. By focusing on $\tau_{c}$ we are generally sacrificing external validity of the estimated treatment effect because in different settings the cutoff $c$ may not be particularly relevant.

A leading reason for focusing on $\tau_{c}$ is that it is generally identified without any assumptions other than that the $\mu_{g}(\cdot)$ functions are continuous at $x=c$. To see why, write\\
$y=(1-w) y_{0}+w y_{1}=1[x<c] y_{0}+1[x \geq c] y_{1}$,\\
and so\\
$m(x) \equiv \mathrm{E}(y \mid x)=1[x<c] \mu_{0}(x)+1[x \geq c] \mu_{1}(x)$.

Now, using continuity of $\mu_{0}(\cdot)$ and $\mu_{1}(\cdot)$ at $c$,\\
$m^{-}(c) \equiv \lim _{x \uparrow c} m(x)=\mu_{0}(c)$\\
$m^{+}(c) \equiv \lim _{x \downarrow c} m(x)=\mu_{1}(c)$.\\
It follows that\\
$\tau_{c}=m^{+}(c)-m^{-}(c)$.

Because we can generally estimate $\mathrm{E}(y \mid x, w=0)$ for all $x<c$ and $\mathrm{E}(y \mid x, w=1)$ for all $x \geq c$, we can estimate the limits of these functions as $x$ approaches $c$ (from the appropriate direction). As a technical point, we must estimate the two regression functions at the boundary value, $c$, but several strategies have proved useful.

One such strategy is local linear regression. Define $\mu_{0 c}=\mu_{0}(c), \mu_{1 c}=\mu_{1}(c)$ and write


\begin{equation*}
y_{0}=\mu_{0 c}+\beta_{0}(x-c)+u_{0} \tag{21.99}
\end{equation*}



\begin{equation*}
y_{1}=\mu_{1 c}+\beta_{1}(x-c)+u_{1} \tag{21.100}
\end{equation*}


so that\\
$y=\mu_{0 c}+\tau_{c} w+\beta_{0}(x-c)+\delta w \cdot(x-c)+r$,\\
where $r=u_{0}+w\left(u_{1}-u_{0}\right)$. The estimate of $\tau_{c}$ is just the jump in the linear function at $x=c$. We could use the entire data set to run the regression\\
$y_{i} \quad$ on $\quad 1, w_{i},\left(x_{i}-c\right), w_{i} \cdot\left(x_{i}-c\right)$\\
and obtain $\hat{\tau}_{c}$ as the coefficient on $w_{i}$, but this approach would be global estimation to estimate a localized average treatment effect, $\tau_{c}$. To make this a "local" procedure, choose a "small" value $h>0$ and only use the data satisfying $c-h<x_{i}<c+h$. Of course, this is equivalent to estimating two separate regressions: $y_{i}$ on $1,\left(x_{i}-c\right)$ for $c-h<x_{i}<c$ and $y_{i}$ on $1,\left(x_{i}-c\right)$ for $c \leq x_{i}<c+h$, and then $\hat{\tau}_{c}=\hat{\mu}_{1 c}-\hat{\mu}_{0 c}$ is the difference in the intercepts. For extra flexibility, we can use a quadratic or cubic in $\left(x_{i}-c\right)$; if a single regression is used, the polynomials should also be interacted with $w_{i}$.

If $h$ is viewed as a fixed value chosen ahead of time by the researcher, inference on $\hat{\tau}_{c}$ is standard: just use a heteroskedasticity-robust $t$ standard error. Imbens and Lemieux (2008) show that if $h$ shrinks to zero quickly enough, the usual inference is still valid. While one can experiment with different choices of $h$-trading off more bias when $h$ is large versus a smaller variance-one can use a data-based method.

Imbens and Kalyanaraman (2009) explicitly look at choosing $h$ to minimize


\begin{equation*}
\mathrm{E}\left\{\left[\hat{\mu}_{0}(c)-\mu_{0}(c)\right]^{2}+\left[\hat{\mu}_{1}(c)-\mu_{1}(c)\right]^{2}\right\} \tag{21.103}
\end{equation*}


a total mean squared error for the two regression functions at the jump point. The optimal bandwidth choice depends on second derivatives of the regression functions at $x=c$, the density of $x_{i}$ at $x=c$, the conditional variances, and the kernel used in local linear regression. See Imbens and Kalraynaram (2008) for details.

Adding regressors is no problem: if the regressors are $\mathbf{r}_{i}$, just run the regression $y_{i}$ on $1, w_{i},\left(x_{i}-c\right), w_{i} \cdot\left(x_{i}-c\right), \mathbf{r}_{i}$, again only using data $c-h<x_{i}<c+h$. As discussed in Imbens and Lemieux (2008), using extra regressors is likely to have more of an impact when $h$ is large, and it might help reduce the bias arising from the deterioration of the linear approximation. Another reason for adding $\mathbf{r}_{i}$ is that doing so can reduce the error variance, possibly improving the precision of $\hat{\tau}_{c}$.

For response variables with limited range, we can use local versions of other estimation methods. For example, suppose the $y_{g}$ are count variables. Then we might use the observations with $c-h<x_{i}<c$ to estimate a Poisson regression $\mathrm{E}(y \mid x, w=0)=\exp \left(\alpha_{0}+\beta_{0} x\right)$ and use $c \leq x_{i}<c+h$ to estimate a Poisson regression $\mathrm{E}(y \mid x, w=1)=\exp \left(\alpha_{1}+\beta_{1} x\right)$. Of course, we could include more flexible functions of $x$, too. If the exponential regression functions are correctly specified for $x$ "near" $c$, we can estimate $\tau_{c}$ as


\begin{equation*}
\hat{\tau}_{c}=\exp \left(\hat{\alpha}_{1}+\hat{\beta}_{1} c\right)-\exp \left(\hat{\alpha}_{0}+\hat{\beta}_{0} c\right) \tag{21.104}
\end{equation*}


\subsection*{21.5.2 The Fuzzy Regression Discontinuity Design}
In the fuzzy regression discontinuity (FRD) design, the probability of treatment changes discontinuously at $x=c$. Define the propensity score as a function of the scalar $x$ as


\begin{equation*}
\mathrm{P}(w=1 \mid x) \equiv F(x) \tag{21.105}
\end{equation*}


As in the SRD case, we still assume that the counterfactual conditional mean functions $\mu_{0}(\cdot)$ and $\mu_{1}(\cdot)$ are continuous at $c$. The key assumption for the FRD design is that $F(\cdot)$ is discontinuous at $c$, so that there is a discrete jump in the probability of treatment at the cutoff.

There are various ways to identify $\tau_{c}$. An assumption that leads to a fairly straightforward analysis is that the individual-specific gain, $y_{1}-y_{0}$, is independent of $w$, conditional on $x$. Therefore, treatment $w$ is allowed to be correlated with $y_{0}$ (after conditioning on $x$ ) but not with the unobserved gain from treatment. Compared with identifying other average treatment effects, estimating $\tau_{\text {att }}$ requires (some version of)\\
ignorability of $w$ with respect to $y_{0}$ while $\tau_{\text {ate }}$ requires that $w$ is ignorable with respect to $\left(y_{0}, y_{1}\right)$ (which clearly implies the other two assumptions).

To see how $\tau_{c}$ is identified, again write $y=y_{0}+w\left(y_{1}-y_{0}\right)$ and use conditional independence between $w$ and $\left(y_{1}-y_{0}\right)$ :

$$
\begin{aligned}
\mathrm{E}(y \mid x) & =\mathrm{E}\left(y_{0} \mid x\right)+\mathrm{E}(w \mid x) \mathrm{E}\left(y_{1}-y_{0} \mid x\right) \\
& =\mu_{0}(x)+\mathrm{E}(w \mid x) \cdot \tau(x)
\end{aligned}
$$

As in the SRD case, take limits from the right and left and use continuity of $\mu_{0}(\cdot)$ and $\tau(\cdot)$ at $c$ :

$$
\begin{aligned}
& m^{+}(c)=\mu_{0}(c)+F^{+}(c) \tau_{c} \\
& m^{-}(c)=\mu_{0}(c)+F^{-}(c) \tau_{c}
\end{aligned}
$$

It follows that, if $F^{+}(c) \neq F^{-}(c)$, then


\begin{equation*}
\tau_{c}=\frac{\left[m^{+}(c)-m^{-}(c)\right]}{\left[F^{+}(c)-F^{-}(c)\right]} \tag{21.106}
\end{equation*}


Because the mean and propensity score are generally identified in a neighborhood of $c, \tau_{c}$ is generally identified.

Given consistent estimators of the four quantities in equation (21.106), we have


\begin{equation*}
\hat{\tau}_{c}=\frac{\left[\hat{m}^{+}(c)-\hat{m}^{-}(c)\right]}{\left[\hat{F}^{+}(c)-\hat{F}^{-}(c)\right]} \tag{21.107}
\end{equation*}


Imbens and Lemieux (2008) suggest estimating $m^{+}(c), m^{-}(c), F^{+}(c)$, and $F^{-}(c)$ all by local linear regression. Namely, $\hat{m}^{+}(c)=\hat{\alpha}_{1 c}, \hat{m}^{-}(c)=\hat{\alpha}_{0 c}, \hat{F}^{+}(c)=\hat{\theta}_{1 c}$, and $\hat{F}^{-}(c)=\hat{\theta}_{0 c}$ are the intercepts from four local linear regressions. For example, $\hat{\alpha}_{1 c}$ is from $y_{i}$ on $1,\left(x_{i}-c\right), c \leq x_{i}<c+h$ and $\hat{\theta}_{1 c}$ is from $w_{i}$ on $1,\left(x_{i}-c\right), c \leq x_{i}<c+h$.

As a computational device, and also for simple inference, HTV (2001) show that equation (21.107) is numerically identical to the local IV estimator of $\tau_{c}$ in the equation


\begin{equation*}
y_{i}=\alpha_{0 c}+\tau_{c} w_{i}+\beta_{0}\left(x_{i}-c\right)+\delta 1\left[x_{i} \geq c\right] \cdot\left(x_{i}-c\right)+e_{i} \tag{21.108}
\end{equation*}


where $z_{i} \equiv 1\left[x_{i} \geq c\right]$ is the IV for $w_{i}$. One uses the data such that $c-h<x_{i}<c+h$. If $h$ is fixed or is decreasing at a rate described in Imbens and Lemieux (2008), one can use the usual heteroskedasticity-robust IV standard error for $\hat{\tau}_{c}$.

Rather than using a local linear model for the probability of treatment, we could, say, use a local logit model-for example,\\
$\mathrm{P}(w=1 \mid x)=\Lambda\left(\eta_{c 0}+\psi_{0}(x-c)\right), \quad x<c$\\
$\mathrm{P}(w=1 \mid x)=\Lambda\left(\eta_{c 1}+\psi_{1}(x-c)\right), \quad x \geq c$\\
and then use\\
$\hat{F}^{+}(c)-\hat{F}^{-}(c)=\Lambda\left(\hat{\eta}_{c 1}\right)-\Lambda\left(\hat{\eta}_{c 0}\right)$,\\
where $\left(\hat{\eta}_{c 0}, \hat{\psi}_{0}\right)$ are from a logit of $w_{i}$ on $1,\left(x_{i}-c\right)$ using $c-h<x_{i}<c$, and similarly for $\left(\hat{\eta}_{c 1}, \hat{\psi}_{1}\right)$.

In the FRD case, we have to choose two bandwidths (even assuming that we use symmetric bandwidths in estimating the regression functions on either side of the jump). We could use the same bandwidth for $\mathrm{P}(w=1 \mid x)$ and $\mathrm{E}(y \mid x)$ or choose them separately using, say, Imbens and Kalyanaraman (2009).

\subsection*{21.5.3 Unconfoundedness versus the Fuzzy Regression Discontinuity}
Unlike in the SRD case, it is possible that overlap can hold for the FRD (although it might be weak in practice). Therefore, we can compare regression adjustment to estimators that exploit the FRD.

It is useful to return to the linear formulation:\\
$y=\mu_{0 c}+\tau_{c} w+\beta_{0}(x-c)+\delta w \cdot(x-c)+u_{0}+w\left(u_{1}-u_{0}\right)$.

Under the ignorability assumption $\mathrm{D}\left(y_{0}, y_{1} \mid w, x\right)=\mathrm{D}\left(y_{0}, y_{1} \mid x\right)$, the composite error has zero mean conditional on ( $w, x$ ), and so OLS (or local regression) consistently estimates $\tau_{c}$. In fact, if we believe unconfoundedness and the linear functional form, we can use all the data and average across $x_{i}$ to estimate $\tau_{\text {ate }}$.

If we only assume $\mathrm{D}\left(y_{1}-y_{0} \mid w, x\right)=\mathrm{D}\left(y_{1}-y_{0} \mid x\right)$-but allow $u_{0}$ to be correlated with $w$-then the OLS estimator is inconsistent. Recall that the estimate of $\tau_{c}$ can be written as\\
$\tilde{\tau}_{c}=\tilde{m}^{+}(c)-\tilde{m}^{-}(c)$,\\
where $\tilde{m}_{1}(x)$ is estimated using the $w_{i}=1$ observations and $\tilde{m}_{0}(x)$ is estimated using the $w_{i}=0$ observations. In other words, the discontinuity in $\mathrm{P}(w=1 \mid x)$ at $x=c$ is essentially ignored. By contrast, the estimator in equation (21.107) is consistent under the weaker version of ignorability, and it directly exploits the jumps in the mean responses and treatment probabilities at $x=c$.

A further benefit of the estimator in equation (21.107) is that it is consistent for ATE for compliers at $x=c$ without unconfoundedness, provided we add a monotonicity assumption. See Imbens and Lemieux (2008) for a detailed treatment.

\subsection*{21.6 Further Issues}
We now discuss some special considerations and extensions ofthe basic methods previously discussed.

\subsection*{21.6.1 Special Considerations for Responses with Discreteness or Limited Range}
We have seen that, when ignorability of selection holds, several methods of estimating average treatment effects are available that depend in no essential way on the nature of $y$. Propensity score weighting and matching can be applied directly whether the response is continuous or discrete or has both features. Methods that rely on regression adjustment (parametric or even series estimation)-either by itself or in conjunction with other methods-work better when the chosen conditional mean functions are good approximations to $\mathrm{E}(y \mid w=1, \mathbf{x})$ and $\mathrm{E}(y \mid w=0, \mathbf{x})$. We already discussed how binary and fractional responses can be modeled using logistic and related functional forms (probably combined with the Bernoulli log likelihood) and how exponential regression functions can be used when $y$ is nonnegative (possibly combined with the Poisson log likelihood). But if $y$ has both discrete and continuous characteristics, we might want to try other models for the conditional expectations. For example, suppose $y_{0}$ and $y_{1}$ are corner solutions. Under Assumption ATE.1, $\mathbf{D}\left(y_{g} \mid w, \mathbf{x}\right)=\mathrm{D}\left(y_{g} \mid \mathbf{x}\right), g=0,1$. Therefore, we can estimate models for $\mathrm{D}(y \mid w=0, \mathbf{x})$ and $\mathrm{D}(y \mid w=1, \mathbf{x})$ that account for the corner nature of $y$. This could be the standard Type I Tobit model if the corner is at zero, or a two-limit Tobit model as in Section 17.7 if there are two corners. But we might also try a hurdle model separately for $w=0$ and $w=1$. Remember, the idea is to eventually obtain good estimates of the conditional means, and a Tobit or hurdle model might do that better than, say, a linear model, even though the Tobit or hurdle model is itself misspecified. Once we have those estimates of the mean functions, we can use equation (21.30), as always.

When we relaxed ignorability and allowed treatment to be correlated with unobservables even after conditioning on $\mathbf{x}$, the IV, correction function, and control function methods discussed in Section 21.4 relied heavily on linear (in parameters and unobservables) functional forms. Even though the linear functional forms can, under certain assumptions, consistently estimate a local average treatment effect, we may want to use nonlinear functions and try to better approximate $\tau_{\text {ate }}$ or $\tau_{\text {att }}$. If the counterfactual responses are binary, we might specify


\begin{align*}
& y_{0}=1\left[\alpha_{0}+\mathbf{x} \boldsymbol{\beta}_{0}+e_{0} \geq 0\right]  \tag{21.110}\\
& y_{1}=1\left[\alpha_{1}+\mathbf{x} \boldsymbol{\beta}_{1}+e_{1} \geq 0\right], \tag{21.111}
\end{align*}


where ( $e_{0}, e_{1}$ ) is independent of ( $\mathbf{x}, \mathbf{z}$ ) and each has a standard normal distribution. If we add, as in the linear case,\\
$w=1\left[\theta_{0}+\mathbf{x} \boldsymbol{\theta}_{1}+\mathbf{z} \boldsymbol{\theta}_{2}+a \geq 0\right]$,\\
where ( $e_{0}, e_{1}, a$ ) is independent of ( $\mathbf{x}, \mathbf{z}$ ) and trivariate normal, then we can estimate $\left(\alpha_{0}, \boldsymbol{\beta}_{0}^{\prime}\right)^{\prime}$ and $\left(\alpha_{1}, \boldsymbol{\beta}_{1}^{\prime}\right)^{\prime}$ by "selection" probits on the $w_{i}=0$ and $w_{i}=1$ subsamples, respectively; see Section 19.6.3. Then $\hat{\tau}_{\text {ate }}(\mathbf{x})=\Phi\left(\hat{\alpha}_{1}+\mathbf{x} \hat{\boldsymbol{\beta}}_{1}\right)-\Phi\left(\hat{\alpha}_{0}+\mathbf{x} \hat{\boldsymbol{\beta}}_{0}\right)$ and\\
$\hat{\tau}_{\text {ate }}=N^{-1} \sum_{i=1}^{N}\left[\Phi\left(\hat{\alpha}_{1}+\mathbf{x}_{i} \hat{\boldsymbol{\beta}}_{1}\right)-\Phi\left(\hat{\alpha}_{0}+\mathbf{x}_{i} \hat{\boldsymbol{\beta}}_{0}\right)\right]$.

If we impose $\boldsymbol{\beta}_{1}=\boldsymbol{\beta}_{0}$ and $e_{1}=e_{0}$, we obtain the bivariate probit model that we discussed in Section 15.7.3. Naturally, it is better to use the more flexible model unless evidence suggests it is not needed.

If the $y_{g}$ are corner solutions, we might use the specification $y_{g}=\max \left(0, \alpha_{g}+\right. \mathbf{x} \boldsymbol{\beta}_{g}+e_{g}$ ) and make the standard assumptions so that $y_{g}$ follows a Tobit. In fact, we could still assume that $\left(e_{0}, e_{1}, a\right)$ is independent of $(\mathbf{x}, \mathbf{z})$ and trivariate normal, but now where $\sigma_{g}^{2}=\operatorname{Var}\left(e_{g}\right)$ are variance parameters to be estimated. Although Tobit models with "endogenous switching" are not common, there is no reason they cannot be useful for estimating ATEs with corner solution responses. Statistically, estimation of the parameters for the control and treatment groups is the same as estimating a Tobit model with endogenous sample selection. (The case with $\boldsymbol{\beta}_{1}=\boldsymbol{\beta}_{0}$ and $e_{1}=e_{0}$ is covered in Problem 17.6.) The estimate of $\tau_{\text {ate }}$ takes the usual form\\
$\hat{\tau}_{\text {ate }}=N^{-1} \sum_{i=1}^{N}\left[m\left(\hat{\alpha}_{1}+\mathbf{x}_{i} \hat{\boldsymbol{\beta}}_{1}, \hat{\sigma}_{1}^{2}\right)-m\left(\hat{\alpha}_{0}+\mathbf{x}_{i} \hat{\boldsymbol{\beta}}_{0}, \hat{\sigma}_{0}^{2}\right)\right]$,\\
where $m(\cdot, \cdot)$ is the conditional mean function for a Tobit model. Hurdle models could be used, too, but estimation becomes even more complicated.

For exponential response functions we can use the "selection correction" described in Section 19.6.4 on the control and treated samples, and then, as usual, construct an average of the difference in estimated counterfactual means. These can be applied under the assumption $\mathrm{E}\left(y_{g} \mid a_{g}, w, \mathbf{x}, \mathbf{z}\right)=\exp \left(a_{g}+\mathbf{x} \boldsymbol{\beta}_{g}\right)$ with $a_{g}$ independent of $(\mathbf{x}, \mathbf{z})$, $g=0,1$ and suitable normality assumptions (at a minimum in the probit model for w). A similar approach can be used for fractional responses when just the counterfactural conditional means are assumed to be correctly specified.

\subsection*{21.6.2 Multivalued Treatments}
So far, we have considered the case with a single treatment level, which we indicate as $w_{i}=1$ with $w_{i}=0$ indicating a control unit. In some cases, program participation\\
can take place at different levels-say, part-time or full-time-or there could be different options for treatment, such as different job training programs. Both cases require an extension of the previous framework and estimation method.

Now suppose that the treatment variable can take on $G+1$ different values, which we label $\{0,1,2, \ldots, G\}$. Typically, zero indicates the control group and $1, \ldots, G$ different levels or options for treatment. Thus, $w_{i}$ takes a value in $\{0,1,2, \ldots, G\}$. Now, of course, there are $G+1$ counterfactual outcomes, which we denote, for a random draw $i,\left\{y_{i g}: g=0,1, \ldots, G\right\}$. The observed response, $y_{i}$, can be expressed as\\
$y_{i}=1\left[w_{i}=0\right] y_{i 0}+1\left[w_{i}=1\right] y_{i 1}+\cdots+1\left[w_{i}=G\right] y_{i G}$.

As usual, we have available a set of covariates, $\mathbf{x}_{i}$. Define $\mu_{g}=\mathrm{E}\left(y_{i g}\right)$ as the population means of the counterfactuals. A sufficient ignorability for identifying the means is the conditional mean independence assumption


\begin{equation*}
\mathrm{E}\left(y_{i g} \mid w_{i}, \mathbf{x}_{i}\right)=\mathrm{E}\left(y_{i g} \mid \mathbf{x}_{i}\right), \quad g=0,1, \ldots, G . \tag{21.116}
\end{equation*}


Under assumption (21.116) it follows easily that


\begin{align*}
\mathrm{E}\left(y_{i} \mid w_{i}, \mathbf{x}_{i}\right)= & 1\left[w_{i}=0\right] \mathrm{E}\left(y_{i 0} \mid \mathbf{x}_{i}\right)+1\left[w_{i}=1\right] \mathrm{E}\left(y_{i 1} \mid \mathbf{x}_{i}\right) \\
& +\cdots+1\left[w_{i}=G\right] \mathrm{E}\left(y_{i G} \mid \mathbf{x}_{i}\right), \tag{21.117}
\end{align*}


which immediately shows that the mean function $\mathrm{E}\left(y_{g} \mid \mathbf{x}\right)$ is identified because


\begin{equation*}
\mathrm{E}\left(y_{g} \mid \mathbf{x}\right)=\mathrm{E}(y \mid w=g, \mathbf{x}) \tag{21.118}
\end{equation*}


We can estimate $\mathrm{E}(y \mid w=g, \mathbf{x})$ for each $g$, given a random sample, by restricting attention to units with $w_{i}=g$. In other words, regression adjustment in the multipletreatment case is an obvious extension of the case when $w_{i}$ is binary. The considerations hold here as in the binary case, including using particular functional forms that account for the nature of $y$ and using nonparametric regression.

Given conditional mean estimates $\left\{\hat{m}_{g}(\mathbf{x}): g=0,1, \ldots, G\right\}$, we can estimate the average treatment effect for treatment level $h$ relative to $g$, say $\tau_{g h}$, as


\begin{equation*}
\hat{\tau}_{g h, r e g}=N^{-1} \sum_{i=1}^{N}\left[\hat{m}_{h}\left(\mathbf{x}_{i}\right)-\hat{m}_{g}\left(\mathbf{x}_{i}\right)\right] \tag{21.119}
\end{equation*}


Or, if $\alpha_{g h}$ is the average treatment effect for those in either group $g$ or $h, \hat{\alpha}_{g h, \text { reg }}$ is obtained by averaging the differences $\hat{m}_{h}\left(\mathbf{x}_{i}\right)-\hat{m}_{g}\left(\mathbf{x}_{i}\right)$ across the subsample with $w_{i}=g$ or $w_{i}=h$. For example, perhaps $g$ is a lower level of participation in job\\
training than $h$, and we want to estimate the average treatment effect of going from part-time to full-time for those who were in one of those treated groups. We can define a measure more like the average treatment effect as $\eta_{g h}=\mathrm{E}\left(y_{h}-y_{g} \mid w=h\right)$, which would be particularly interesting when $g=0$ is the control group. Now we would simply average $\hat{m}_{h}\left(\mathbf{x}_{i}\right)-\hat{m}_{g}\left(\mathbf{x}_{i}\right)$ across the subsample $w_{i}=h$.

It is pretty clear that overlap is needed to estimate average treatment effects with multiple treatment levels. While we can get by with less for some cases, the most straightforward statement is based on the set of propensity scores (which Imbens, 2000, calls the generalized propensity score):\\
$\mathrm{P}\left(w_{g}=1 \mid \mathbf{x}\right) \equiv p_{g}(\mathbf{x})>0, \quad \mathbf{x} \in \mathscr{X}, \quad g=0, \ldots, G$.

Because the propensity scores must sum to unity, assumption (21.120) rules out the case of unit probabilities for any $g$.

Propensity score weighting can also be used under the previous ignorability and overlap assumptions. For example, using the same argument in Section 21.3.3, it is easily shown that\\
$\mathrm{E}\left(y_{g}\right)=\mathrm{E}\left\{\frac{1\left[w_{i}=g\right] y_{i}}{p_{g}\left(\mathbf{x}_{i}\right)}\right\}, \quad g=0,1, \ldots, G$,\\
and so consistent estimates of the counterfactual means take the form\\
$\hat{\mu}_{g, p s}=N^{-1} \sum_{i=1}^{N}\left\{\frac{1\left[w_{i}=g\right] y_{i}}{\hat{p}_{g}\left(\mathbf{x}_{i}\right)}\right\}$,\\
where $\hat{p}_{g}(\cdot)$ are the estimated propensity scores. Given these estimates, we can form differences such as $\hat{\tau}_{g h, p s}=\hat{\mu}_{h, p s}-\hat{\mu}_{g, p s}$. Estimating the other kinds of treament effects mentioned previously requires more care, but is fairly straightforward. See, for example Imbens (2000) and Lechner (2001).

To implement IPW estimation, we have to estimate the propensity scores. A common parametric approach would be to use a multinomial logit (MNL) model with flexible functions in $\mathbf{x}_{i}$. Because we are just looking for good estimates of the probabilities, the MNL model will often work well, but other approaches-such as nested logit models-can be used. If the treatment categories are obviously ordered-say, $w_{i}$ is number of years in college for the population of high school graduates-one might try an ordered model, such as ordered logit.

In addition to regression adjustment and propensity score weighting, one can, of course, use matching-either on the covariates or propensity scores. See, for example, Lechner (2001).

\subsection*{21.6.3 Multiple Treatments}
The case of multiple treatments is more difficult to handle, and in some cases relies on additional functional form restrictions. Of course, if $\mathbf{w}_{i}$ is an $M$-vector of treatment variables and the $\mathbf{x}_{i}$ are the controls, we can always work with models for $\mathrm{E}\left(y_{i} \mid \mathbf{w}_{i}, \mathbf{x}_{i}\right)$ directly, but then we are circumventing the counterfactual framework.

A fairly general framework-which actually encompasses the framework of the previous section-is to use a random coefficient setting. Let $\mathbf{w}_{i}$ be $1 \times M$ and $\mathbf{c}_{i}$ an $M \times 1$ vector of unobserved heterogeneity. Assume that\\
$y_{i}=\mathbf{w}_{i} \mathbf{c}_{i}=w_{i 1} c_{i 1}+\cdots+w_{i M} c_{i M}$.

By choosing the $c_{i m}$ as counterfactual outcomes and the $w_{i m}$ as treatment indicators for each treatment level (including the control), we can put the multivalued treatment framework into the current setting. But we can also allow for truly different treatments. For example, suppose we have two programs, A and B , where participation is allowed in one or both. Then we can define $\mathbf{w}_{i}$ to have four dummy variables indicating every possible treatment state: participation in neither, participation in A but not B , participation in B but not A , and participation in both. Or, some of the $w_{\text {im }}$ can be continuous treatments along with discrete treatments. Further, as described in Wooldridge (2004), the $w_{\text {im }}$ can be different functions of the same (nonbinary) treatment in order to make functional forms more flexible. If we want, we can set, say, $w_{i 1} \equiv 1$ so that there is an intercept in the equation; in other cases it is more convenient to have a full set of treatment indicators and exclude an intercept.

Given covariates $\mathbf{x}_{i}$, the goal is to estimate $\mathrm{E}\left(\mathbf{c}_{i} \mid \mathbf{x}_{i}=\mathbf{x}\right)$ and then, eventually, $\boldsymbol{\mu}_{\mathbf{c}}=\mathrm{E}\left(\mathbf{c}_{i}\right)$. When we use appropriate linear combinations, the latter corresponds to average treatment effects across the entire population while $\mathrm{E}\left(\mathbf{c}_{i} \mid \mathbf{x}_{i}=\mathbf{x}\right)$ generally corresponds to conditional average treatment effects.

If we assume ignorability, we can consistently estimate $\mathrm{E}\left(\mathbf{c}_{i} \mid \mathbf{x}_{i}=\mathbf{x}\right)$, and, assuming overlap, we can then average to estimate $\boldsymbol{\mu}_{\mathbf{c}}$. The key assumptions in the population are


\begin{equation*}
\mathrm{E}(y \mid \mathbf{w}, \mathbf{c}, \mathbf{x})=\mathrm{E}(y \mid \mathbf{w}, \mathbf{c}) \tag{21.124}
\end{equation*}


and\\
$\mathrm{E}\left(\mathbf{w}^{\prime} \mathbf{w} \mid \mathbf{c}, \mathbf{x}\right)=\mathrm{E}\left(\mathbf{w}^{\prime} \mathbf{w} \mid \mathbf{x}\right)$.

In many cases, we would obtain assumption (21.125) from $\mathrm{E}(\mathbf{w} \mid \mathbf{c}, \mathbf{x})=\mathrm{E}(\mathbf{w} \mid \mathbf{x})$ and $\operatorname{Var}(\mathbf{w} \mid \mathbf{c}, \mathbf{x})=\operatorname{Var}(\mathbf{w} \mid \mathbf{x})$-in other words, from ignorability conditional on the first two moments of the distribution of $\mathbf{w}$ given $(\mathbf{c}, \mathbf{x})$.

If we also assume that $\boldsymbol{\Lambda}(\mathbf{x}) \equiv \mathrm{E}\left(\mathbf{w}^{\prime} \mathbf{w} \mid \mathbf{x}\right)$ is nonsingular, we can easily find an expression for $\mathrm{E}(\mathbf{c} \mid \mathbf{x})$. Write $y=\mathbf{w c}+u$ where $\mathrm{E}(u \mid \mathbf{w}, \mathbf{c}, \mathbf{x})=0$ by assumption (21.124). Then $\mathbf{w}^{\prime} y=\mathbf{w}^{\prime} \mathbf{w c}+\mathbf{w}^{\prime} u$, and so\\
$\mathrm{E}\left(\mathbf{w}^{\prime} y \mid \mathbf{x}\right)=\mathrm{E}\left(\mathbf{w}^{\prime} \mathbf{w c} \mid \mathbf{x}\right)+\mathrm{E}\left(\mathbf{w}^{\prime} u \mid \mathbf{x}\right)=\mathrm{E}\left(\mathbf{w}^{\prime} \mathbf{w c} \mid \mathbf{x}\right)$.\\
Now, by iterated expectations, $\mathrm{E}\left(\mathbf{w}^{\prime} \mathbf{w c} \mid \mathbf{x}\right)=\mathrm{E}\left[\mathrm{E}\left(\mathbf{w}^{\prime} \mathbf{w c} \mid \mathbf{c}, \mathbf{x}\right) \mid \mathbf{x}\right]$, and\\
$\mathrm{E}\left[\mathrm{E}\left(\mathbf{w}^{\prime} \mathbf{w c} \mid \mathbf{c}, \mathbf{x}\right) \mid \mathbf{x}\right]=\mathrm{E}\left[\mathrm{E}\left(\mathbf{w}^{\prime} \mathbf{w} \mid \mathbf{c}, \mathbf{x}\right) \mathbf{c} \mid \mathbf{x}\right]=\mathrm{E}\left[\mathrm{E}\left(\mathbf{w}^{\prime} \mathbf{w} \mid \mathbf{x}\right) \mathbf{c} \mid \mathbf{x}\right]=\mathrm{E}\left(\mathbf{w}^{\prime} \mathbf{w} \mid \mathbf{x}\right) \mathrm{E}(\mathbf{c} \mid \mathbf{x})$, where the second equality follows from assumption (21.125). We have shown that\\
$\mathrm{E}\left(\mathbf{w}^{\prime} y \mid \mathbf{x}\right)=\mathrm{E}\left(\mathbf{w}^{\prime} \mathbf{w} \mid \mathbf{x}\right) \mathrm{E}(\mathbf{c} \mid \mathbf{x})$,\\
and so, assuming invertibility of $\mathrm{E}\left(\mathbf{w}^{\prime} \mathbf{w} \mid \mathbf{x}\right)$,\\
$\mathrm{E}(\mathbf{c} \mid \mathbf{x})=\left[\mathrm{E}\left(\mathbf{w}^{\prime} \mathbf{w} \mid \mathbf{x}\right)\right]^{-1} \mathrm{E}\left(\mathbf{w}^{\prime} y \mid \mathbf{x}\right) \equiv[\boldsymbol{\Lambda}(\mathbf{x})]^{-1} \mathrm{E}\left(\mathbf{w}^{\prime} y \mid \mathbf{x}\right)$.

Because we can collect a random sample on $\left(y_{i}, \mathbf{w}_{i}, \mathbf{x}_{i}\right)$, we can estimate all the conditional moments that appear in $\mathrm{E}(\mathbf{c} \mid \mathbf{x})$.

For estimating the unconditional mean $\boldsymbol{\mu}_{\mathbf{c}}$ it suffices to estimate $\boldsymbol{\Lambda}(\mathbf{x})$ because\\
$\boldsymbol{\mu}_{\mathbf{c}}=\mathrm{E}\left\{[\boldsymbol{\Lambda}(\mathbf{x})]^{-1} \mathrm{E}\left(\mathbf{w}^{\prime} y \mid \mathbf{x}\right)\right\}=\mathrm{E}\left\{[\boldsymbol{\Lambda}(\mathbf{x})]^{-1} \mathbf{w}^{\prime} y\right\}$,\\
which just uses the fact that $\mathbf{w}^{\prime} y=\mathrm{E}\left(\mathbf{w}^{\prime} y \mid \mathbf{x}\right)+\mathbf{r}$ with $\mathrm{E}(\mathbf{r} \mid \mathbf{x}) \equiv \mathbf{0}$.\\
Expression (21.127) takes on a simple, recognizable form when $\mathbf{w}$ is a vector of mutually exclusive, exhaustive binary indicators for different treament groups. Then $\mathbf{w}^{\prime} \mathbf{w}=\operatorname{diag}\left(w_{i 1}, \ldots, w_{i M}\right)$ and so $\boldsymbol{\Lambda}(\mathbf{x})$ is the $M \times M$ diagonal matrix with $\mathrm{E}\left(w_{m} \mid \mathbf{x}\right)= \mathrm{P}\left(w_{m}=1 \mid \mathbf{x}\right)$-that is, the propensity scores-down the diagonal. It is easily seen that the $m$ th element of $\left.\mathrm{E}\left\{[\boldsymbol{\Lambda}(\mathbf{x})]^{-1} \mathbf{w}^{\prime} y\right)\right\}$ is simply\\
$\mathrm{E}\left[w_{i m} y_{i} / p_{m}\left(\mathbf{x}_{i}\right)\right]$,\\
which is exactly the expression we derived in the previous subsection to motivate propensity score weighting in the multiple treatment case.

More generally, $\boldsymbol{\Lambda}(\mathbf{x})$ might not be a diagonal matrix, although it always will be if we choose $\mathbf{w}_{i}$ to saturate all possible treatment scenarios. In other words, if we have $Q$ possible programs where, for program $q$ there are $J_{q}$ treatment levels, then $\mathbf{w}_{i}$ has $\sum_{q=1}^{Q} J_{q}$ elements. (If certain combinations are impossible, they can just be excluded.) With continuous treatments, or where we, say, let each program have its own set of treatment effects and these do not interact with treatment effects from other programs, then we generally have to estimate conditional covariances along with conditional means (and variances). Assuming we have a consistent estimator $\hat{\boldsymbol{\Lambda}}(\mathbf{x})$ of $\boldsymbol{\Lambda}(\mathbf{x})$\\
for each $\mathbf{x}$, consistent estimators of $\boldsymbol{\mu}_{\mathbf{c}}$ have the form\\
$\hat{\boldsymbol{\mu}}_{\mathbf{c}}=N^{-1} \sum_{i=1}^{N}\left[\hat{\boldsymbol{\Lambda}}\left(\mathbf{x}_{i}\right)\right]^{-1} \mathbf{w}_{i}^{\prime} y_{i}$.

In many cases it is useful to explicitly separate the random slopes from the random intercept, in which case we can write in the population\\
$\mathrm{E}(y \mid \mathbf{w}, \mathbf{c})=\mathrm{E}(y \mid \mathbf{w}, \mathbf{c}, \mathbf{x})=a+\mathbf{w b}$,\\
where $\mathbf{w}$ is now a set of $K$ treatment variables such that $\boldsymbol{\Omega}(\mathbf{x}) \equiv \operatorname{Var}(\mathbf{w} \mid \mathbf{x})$ is nonsingular for all (relevant) values of $\mathbf{x}$. Then, as can be derived from equation (21.127) using partitioned inverse (see also Wooldridge (2004)),\\
$\mathrm{E}(\mathbf{b} \mid \mathbf{x})=[\operatorname{Var}(\mathbf{w} \mid \mathbf{x})]^{-1} \operatorname{Cov}(\mathbf{w}, y \mid \mathbf{x}) \equiv[\boldsymbol{\Omega}(\mathbf{x})]^{-1} \operatorname{Cov}(\mathbf{w}, y \mid \mathbf{x})$\\
and\\
$\boldsymbol{\mu}_{\mathbf{b}} \equiv \mathrm{E}(\mathbf{b})=\mathrm{E}\left\{[\boldsymbol{\Omega}(\mathbf{x})]^{-1}[\mathbf{w}-\boldsymbol{\psi}(\mathbf{x})] y\right\}$,\\
where $\boldsymbol{\psi}(\mathbf{x}) \equiv \mathrm{E}(\mathbf{w} \mid \mathbf{x})$ is the $1 \times K$ vector of conditional mean functions. In this formulation, we see directly that the conditional mean and conditional variancecovariance matrix of the treatment $\mathbf{w}$ are needed to estimate the average slopes, $\boldsymbol{\mu}_{\mathbf{b}}$. Because we assume a random sample on ( $\mathbf{w}, \mathbf{x}$ ), these moments are generally identified. In practice, we might use parametric models that account for the nature of the elements of $\mathbf{w}$. Naturally, given such estimators, we estimate $\boldsymbol{\mu}_{\mathbf{b}}$ as a sample average, $\hat{\boldsymbol{\mu}}_{\mathbf{b}}=N^{-1} \sum_{i=1}^{N}\left[\hat{\boldsymbol{\Omega}}\left(\mathbf{x}_{i}\right)\right]^{-1}\left[\mathbf{w}_{i}-\hat{\boldsymbol{\psi}}\left(\mathbf{x}_{i}\right)\right] y_{i}$.

When treatment is not ignorable for a set of covariates $\mathbf{x}$, we need to obtain instrumental variables. Linear IV estimation is relatively straightforward provided unobserved heterogeneity is additive, as in Section 21.4.1. But allowing for general heterogeneous treatment effects without ignorability is difficult. Control function approaches would rely on\\
$\mathrm{E}(y \mid \mathbf{w}, \mathbf{x}, \mathbf{z})=\mathbf{w} \mathrm{E}(\mathbf{c} \mid \mathbf{w}, \mathbf{x}, \mathbf{z})$,\\
where $\mathbf{z}$ is the vector of instruments, and the latter can be difficult to obtain when $\mathbf{w}$ is a vector. Wooldridge (2008) considered correction function approaches under the assumptions


\begin{equation*}
\mathrm{E}(y \mid \mathbf{w}, \mathbf{c}, \mathbf{x}, \mathbf{z})=\mathrm{E}(y \mid \mathbf{w}, \mathbf{c}) \tag{21.129}
\end{equation*}


(which is not much of an assumption because can include lots of factors) and\\
$\mathrm{E}(\mathbf{c} \mid \mathbf{x}, \mathbf{z})=\mathrm{E}(\mathbf{c} \mid \mathbf{x})=\boldsymbol{\mu}_{\mathbf{c}}+\left(\mathbf{x}-\boldsymbol{\mu}_{\mathbf{x}}\right) \boldsymbol{\Gamma}$,\\
which is a standard exclusion restriction on the instruments, $\mathbf{z}$, along with linearity assumption. If we write $c_{m}=\mu_{c m}+\left(\mathbf{x}-\boldsymbol{\mu}_{\mathbf{x}}\right) \boldsymbol{\gamma}_{m}+v_{m}$, then we have

$$
\begin{aligned}
y_{i}= & \mu_{c 1} w_{i 1}+\mu_{c 2} w_{i 2}+\cdots+\mu_{c M} w_{i m}+w_{i 1}\left(\mathbf{x}_{i}-\boldsymbol{\mu}_{\mathbf{x}}\right) \gamma_{1}+\cdots+w_{i M}\left(\mathbf{x}_{i}-\boldsymbol{\mu}_{\mathbf{x}}\right) \gamma_{M} \\
& +w_{i 1} v_{i 1}+\cdots+w_{i M} v_{i M}
\end{aligned}
$$

The correction function approach now entails finding $\mathrm{E}\left(w_{\text {im }} u_{\text {im }} \mid \mathbf{x}_{i}, \mathbf{z}_{i}\right)$ and inserting these for the $w_{\text {im }} v_{\text {im }}$, and then applying instrumental variables using functions of $\left(\mathbf{x}_{i}, \mathbf{z}_{i}\right)$. More specifically, suppose $w_{i m}=f_{m}\left(\mathbf{x}_{i}, \mathbf{z}_{i}, u_{i m}, \boldsymbol{\alpha}_{m}\right)$ and $\mathrm{E}\left(v_{i m} \mid u_{i m}, \mathbf{x}_{i}, \mathbf{z}_{i}\right)= \rho_{m} u_{i m}$. Then, with distributional assumptions on $u_{i m}$, we can estimate $\boldsymbol{\alpha}_{m}$ and parameters in the distribution of $u_{i m}$ to obtain $\mathrm{E}\left(w_{i m} u_{i m} \mid \mathbf{x}_{i}, \mathbf{z}_{i}\right)=h_{m}\left(\mathbf{x}_{i}, \mathbf{z}_{i}, \boldsymbol{\theta}_{m}\right)$, where the $\boldsymbol{\theta}_{m}$ are assumed to be identified and $h_{m}(\cdot)$ is known. (Section 21.4.2 considered the case where $w_{i}$ is binary and follows a probit; the correction function was the standard normal pdf evaluated at a linear index.) Now we have the estimating equation


\begin{align*}
y_{i}= & \mu_{c 1} w_{i 1}+\mu_{c 2} w_{i 2}+\cdots+\mu_{c M} w_{i m}+w_{i 1}\left(\mathbf{x}_{i}-\boldsymbol{\mu}_{\mathbf{x}}\right) \gamma_{1}+\cdots+w_{i M}\left(\mathbf{x}_{i}-\boldsymbol{\mu}_{\mathbf{x}}\right) \gamma_{M} \\
& +\rho_{1} h_{1}\left(\mathbf{x}_{i}, \mathbf{z}_{i}, \boldsymbol{\theta}_{1}\right)+\cdots+\rho_{M} h_{M}\left(\mathbf{x}_{i}, \mathbf{z}_{i}, \boldsymbol{\theta}_{M}\right)+r_{i} \tag{21.131}
\end{align*}


$\mathrm{E}\left(r_{i} \mid \mathbf{x}_{i}, \mathbf{z}_{i}\right)=0$.\\
After replacing $\boldsymbol{\mu}_{\mathbf{x}}$ with the sample average $\overline{\mathbf{x}}$, plugging in the $\hat{\boldsymbol{\theta}}_{m}$, and specifying the instruments, we can estimate the equation by IV. Natural choices for the instruments are $\hat{\mathrm{E}}\left(w_{i m} \mid \mathbf{x}_{i}, \mathbf{z}_{i}\right)$ and $\hat{\mathrm{E}}\left(w_{i m} \mid \mathbf{x}_{i}, \mathbf{z}_{i}\right) \cdot\left(\mathbf{x}_{i}-\overline{\mathbf{x}}\right)$, where the $\hat{\mathrm{E}}\left(w_{i m} \mid \mathbf{x}_{i}, \mathbf{z}_{i}\right)$ are obtained from the specified distribution $\mathrm{D}\left(w_{i m} \mid \mathbf{x}_{i}, \mathbf{z}_{i}\right)$. In the simple binary case where $w_{i}$ follows a probit, we used the probit fitted values, $\hat{\Phi}_{i}$.

A test for whether all interactions between the treatment variables and heterogeneity have zero coefficients is a standard, heteroskedasticity Wald test of $H_{0}: \rho_{1}=\cdots =\rho_{M}=0$ after IV estimation. If we want to allow some of the $\rho_{m}$ to be nonzero, the variance matrix of all second-step estimators needs to be adjusted for the estimation of the $\boldsymbol{\theta}_{m}$, using either the delta method or bootstrapping. Notice that the coefficients on the $w_{\text {im }}$ are the estimated counterfactual means, and then we can use these to construct various average treatment effects. We also have direct estimates of how the treatment effects vary with $\mathbf{x}$ (because we have the $\hat{\gamma}_{m}$ ).

The correction function method is easy to apply when the marginal models for the treatment indicators are easy to obtain. For example, it is much easier if, say, $w_{i m}$ is assumed to follow a probit model than if a collection of indicators follows a multivariate probit. The method applies also when some treatments have discrete and continuous characteristics, such as hours spent in a job training program. If $w_{\text {im }}$ follows a Tobit model, then $h_{m}\left(\mathbf{x}_{i}, \mathbf{z}_{i}, \boldsymbol{\theta}_{m}\right)$ is tractable; see Wooldridge (2008).

\subsection*{21.6.4 Panel Data}
Using individual panel data with general patterns of treatments is complicated by the different kinds of treatment effects-static and dynamic-that one might like to consider. In this subsection, we consider the case where, in each time period $t$, unit $i$ is part of a treatment or control group. Thus, $w_{i t}$ is a binary variable where $w_{i t}=1$ means treatment in time $t$. Let $\mathbf{w}_{i}=\left(w_{i 1}, \ldots, w_{i T}\right)$ denote the entire history of treatment indicators. If we focuse on the ATEs of treatment in a particular time period, it is fairly straightforward to state ignorability assumptions conditional on unobserved heterogeneity. Let $y_{i t}(0)$ and $y_{i t}(1)$ denote the counterfactual outcomes in the untreated and treated states, respectively. Then ignorability conditional on unobserved heterogeneity $\mathbf{c}_{i}$ and covariates $\mathbf{x}_{i}$ is


\begin{equation*}
\mathrm{E}\left[y_{i t}(g) \mid \mathbf{w}_{i}, \mathbf{c}_{i}, \mathbf{x}_{i}\right]=\mathrm{E}\left[y_{i t}(g) \mid \mathbf{c}_{i}, \mathbf{x}_{i}\right], \quad g=0,1 . \tag{21.132}
\end{equation*}


Notice that this is a strict exogeneity assumption on treatment assignment, conditional on $\left(\mathbf{c}_{i}, \mathbf{x}_{i}\right)$. Because $y_{i t}=y_{i t}(0)+w_{i t} \mathrm{E}\left[y_{i t}(1)-y_{i t}(0)\right]$, it follows that


\begin{align*}
\mathrm{E}\left(y_{i t} \mid \mathbf{w}_{i}, \mathbf{c}_{i}, \mathbf{x}_{i}\right) & =\mathrm{E}\left(y_{i t} \mid w_{i t}, \mathbf{c}_{i}, \mathbf{x}_{i}\right) \\
& =\mathrm{E}\left[y_{i t}(0) \mid \mathbf{c}_{i}, \mathbf{x}_{i}\right]+w_{i t} \mathrm{E}\left[y_{i t}(1)-y_{i t}(0) \mid \mathbf{c}_{i}, \mathbf{x}_{i}\right] \tag{21.133}
\end{align*}


If the treatment effect $y_{i t}(1)-y_{i t}(0)$ is constant for each $t$, say $\tau_{t}$, and $\mathrm{E}\left[y_{i t}(0) \mid \mathbf{c}_{i}, \mathbf{x}_{i}\right]$ is linear and additive with scalar heterogeneity, and only covariates at time $t$ appear, then\\
$\mathrm{E}\left(y_{i t} \mid \mathbf{w}_{i}, \mathbf{c}_{i}, \mathbf{x}_{i}\right)=c_{i 0}+\alpha_{0 t}+\mathbf{x}_{i t} \boldsymbol{\beta}_{0 t}+\tau_{t} w_{i t}, \quad t=1, \ldots, T$,\\
and this leads to a standard fixed-effects or first-differencing analysis (especially if we assume time homogeneity of $\boldsymbol{\beta}_{0 t}$ and $\tau_{t}$ ).

If we allow $\mathrm{E}\left[y_{i t}(1)-y_{i t}(0) \mid \mathbf{c}_{i}, \mathbf{x}_{i}\right]$ to depend on heterogeneity and covariates, we get an estimating equation where $w_{i t}$ interacts with $\mathbf{x}_{i t}$ as well as with heterogeneity, say $a_{i}$. So, assuming time homogeneity except in the intercepts $\alpha_{0 t}$ and the means of the covariates, the estimating equation looks like


\begin{equation*}
\mathrm{E}\left(y_{i t} \mid \mathbf{w}_{i}, \mathbf{c}_{i}, \mathbf{x}_{i}\right)=c_{i 0}+\alpha_{0 t}+\mathbf{x}_{i t} \boldsymbol{\beta}_{0}+a_{i} w_{i t}+w_{i t}\left(\mathbf{x}_{i t}-\boldsymbol{\psi}_{t}\right), \quad t=1, \ldots, T, \tag{21.134}
\end{equation*}


where $\boldsymbol{\psi}_{t}=\mathrm{E}\left(\mathbf{x}_{i t}\right)$ (which should be replaced by the sample average for each $t$ in estimation). The average treatment effect is $\tau=\mathrm{E}\left(a_{i}\right)$. We discussed how to estimate such models in Section 11.7, and we also noted that there are conditions where setting $a_{i}=\tau$ in estimation and using the usual FE estimator consistently estimates $\tau$.

If we are willing to assume strictly exogenous treatment, conditional on heterogeneity and covariates, then allowing lagged effects of treatment is straightforward in a\\
regression-type framework. As a practical matter, one just adds lags of treatment indicators to the standard models described previously. Then, one can estimate how long program participation effects last. Explicitly using a counterfactual setting is possible but notationally complicated. For example, suppose we are interested only in contemporaneous effects and effects lagged a year. Then, for each $t \geq 2$, we have counterfactual outcomes $y_{i t}\left(g_{t-1}, g_{t}\right)$ where $g_{t-1}$ and $g_{t}$ can be zero or one, corresponding to treatment in the same time period or one time period earlier. In other words, at time $t$ there are four counterfactuals for unit $i$ : $y_{i t}(0,0), y_{i t}(1,0), y_{i t}(0,1)$, and $y_{i t}(1,1)$. We can write the observed outcome as

$$
\begin{aligned}
y_{i t}= & \left(1-w_{i, t-1}\right)\left(1-w_{i t}\right) y_{i t}(0,0)+w_{i, t-1}\left(1-w_{i t}\right) y_{i t}(1,0) \\
& +\left(1-w_{i, t-1}\right) w_{i t} y_{i t}(0,1)+w_{i, t-1} w_{i t} y_{i t}(1,1) \\
= & y_{i t}(0,0)+w_{i, t-1}\left[y_{i t}(1,0)-y_{i t}(0,0)\right]+w_{i t}\left[y_{i t}(0,1)-y_{i t}(0,0)\right] \\
& +w_{i, t-1} w_{i t}\left[y_{i t}(1,1)-y_{i t}(1,0)-y_{i t}(0,1)+y_{i t}(0,0)\right]
\end{aligned}
$$

If we modify assumption (21.132) to\\
$\mathrm{E}\left[y_{i t}\left(g_{t-1}, g_{t}\right) \mid \mathbf{w}_{i}, \mathbf{c}_{i}, \mathbf{x}_{i}\right]=\mathrm{E}\left[y_{i t}\left(g_{t-1}, g_{t}\right) \mid \mathbf{c}_{i}, \mathbf{x}_{i}\right], \quad g_{t-1}, g_{t}=0,1$,\\
then unobserved-effects models fall out naturally. The standard additive model with covariates dated at time $t$ (and possibly earlier time periods) emerges if only $\mathrm{E}\left[y_{i t}(0,0) \mid \mathbf{c}_{i}, \mathbf{x}_{i}\right]$ depends on heterogeneity and covariates, but adding interactions with heterogeneity, and especially with observed covariates (contemporaneous or lagged) is straightforward. If we treat all means as constant over time- $\mu_{g h}= \mathrm{E}\left[y_{i t}(g, h)\right]$-then it is straightforward to estimate various treatment effects. For example, $\mu_{11}-\mu_{00}$ is the average treatment effect from participating in neither period versus participating in both periods.

We can also adapt an approach due to Altonji and Matzkin (2005); see also Wooldridge (2005a). Under assumption (21.135),

$$
\begin{aligned}
\mathrm{E}\left(y_{i t} \mid \mathbf{w}_{i}, \mathbf{c}_{i}, \mathbf{x}_{i}\right)= & \mathrm{E}\left(y_{i t} \mid w_{i, t-1}, w_{i t}, \mathbf{c}_{i}, \mathbf{x}_{i}\right)=\left(1-w_{i, t-1}\right)\left(1-w_{i t}\right) \mathrm{E}\left[y_{i t}(0,0) \mid \mathbf{c}_{i}, \mathbf{x}_{i}\right] \\
& +w_{i, t-1}\left(1-w_{i t}\right) \mathrm{E}\left[y_{i t}(1,0) \mid \mathbf{c}_{i}, \mathbf{x}_{i}\right] \\
& +\left(1-w_{i, t-1}\right) w_{i t} \mathrm{E}\left[y_{i t}(0,1) \mid \mathbf{c}_{i}, \mathbf{x}_{i}\right]+w_{i, t-1} w_{i t} \mathrm{E}\left[y_{i t}(1,1) \mid \mathbf{c}_{i}, \mathbf{x}_{i}\right]
\end{aligned}
$$

Unless we make special functional form assumptions of the kind just described, this equation is not directly useful because it depends on the unobserved heterogeneity, $\mathbf{c}_{i}$. Suppose, however, that we assume that the distribution of $\mathbf{c}_{i}$ given ( $\mathbf{w}_{i}, \mathbf{x}_{i}$ ) depends on a relatively simple function of the history of treatments, such as the fraction of\\
treated periods, $\bar{w}_{i}$. (There are more sophisticated possibilities, such as using fractions within various subperiods.) Formally, if $\mathrm{D}\left(\mathbf{c}_{i} \mid \mathbf{w}_{i}, \mathbf{x}_{i}\right)=\mathrm{D}\left(\mathbf{c}_{i} \mid \bar{w}_{i}, \mathbf{x}_{i}\right)$, then, by iterated expectations and assumption (21.135),


\begin{align*}
\mathrm{E}\left(y_{i t} \mid \mathbf{w}_{i}, \mathbf{x}_{i}\right)= & \mathrm{E}\left(y_{i t} \mid w_{i, t-1}, w_{i t}, \bar{w}_{i}, \mathbf{x}_{i}\right)=\left(1-w_{i, t-1}\right)\left(1-w_{i t}\right) \mathrm{E}\left[y_{i t}(0,0) \mid \bar{w}_{i}, \mathbf{x}_{i}\right] \\
& +w_{i, t-1}\left(1-w_{i t}\right) \mathrm{E}\left[y_{i t}(1,0) \mid \bar{w}_{i}, \mathbf{x}_{i}\right] \\
& +\left(1-w_{i, t-1}\right) w_{i t} \mathrm{E}\left[y_{i t}(0,1) \mid \bar{w}_{i}, \mathbf{x}_{i}\right] \\
& +w_{i, t-1} w_{i t} \mathrm{E}\left[y_{i t}(1,1) \mid \bar{w}_{i}, \mathbf{x}_{i}\right] \tag{21.136}
\end{align*}


Actually, we could have derived this equation directly if we had just assumed that treatment at $t-1$ and $t$ is unconfounded conditional on $\left(\bar{w}_{i}, \mathbf{x}_{i}\right)$, but the current derivation makes a link to approaches with unobserved heterogeneity.

It follows now from equation (21.136) that we can estimate $\mathrm{E}\left[y_{i t}\left(g_{t-1}, g_{t}\right) \mid \bar{w}_{i}, \mathbf{x}_{i}\right]$ by estimating $\mathrm{E}\left(y_{i t} \mid w_{i, t-1}=g_{t-1}, w_{i t}=g_{t}, \bar{w}_{i}, \mathbf{x}_{i}\right) \equiv m_{t}\left(g_{t-1}, g_{t}, \bar{w}_{i}, \mathbf{x}_{i}\right)$. That is, for each of the four treatment group combinations, we estimate regressions of $y_{i t}$ on $\left(\bar{w}_{i}, \mathbf{x}_{i}\right)$, or use quasi-MLE. (If $\mathbf{x}_{i}$ is truly a sequence of covariates that change over time, we might restrict the way $\mathbf{x}_{i}$ appears, such as $\left(\mathbf{x}_{i t}, \mathbf{x}_{i, t-1}, \overline{\mathbf{x}}_{i}\right)$, but this restriction is only intended to conserve on the dimension of the estimation problem.) Given the $\hat{m}_{t}\left(g_{t-1}, g_{t}, \bar{w}_{i}, \mathbf{x}_{i}\right)$, we have


\begin{equation*}
\hat{\mathrm{E}}\left[y_{i t}\left(g_{t-1}, g_{t}\right)\right]=N^{-1} \sum_{i=1}^{N} \hat{m}_{t}\left(g_{t-1}, g_{t}, \bar{w}_{i}, \mathbf{x}_{i}\right), \tag{21.137}
\end{equation*}


that is, we average out the control variables ( $\bar{w}_{i}, \mathbf{x}_{i}$ ). This approach can be made quite flexible and allows estimation of average treatment effects that compare any of the four groups to any other group. And, of course, nothing about the method requires us to focus on only current and one lag of treatment status. The same general approach applies to histories $\mathbf{g}^{t}=\left(g_{1}^{t}, g_{2}^{t}, \ldots, g_{t}^{t}\right)$ where each $g_{r}^{t}$ is zero or one.

The previous approaches assume ignorability of treatment conditional on heterogeneity, or a sufficient statistic for the entire treatment history (such as the time average), and a sequence of observed covariates. It may be more realistic to assume ignorability conditional on past observed outcomes, treatment assignments, and covariates. For example, if workers are being assigned into job training, an administrator may be more likely to make assignments on the basis of past observed labor market outcomes and past assignment status. General frameworks can be found in Gill and Robins (2001), Lechner and Miquel (2005), and Abbring and Heckman (2007). Here I follow Lechner (2004) and use a dynamic ignorability\\
assumption. The setup is that for a total of $T$ time periods we have a balanced panel. For each $i$, we observe the sequence of binary treatment indicators, $\left\{w_{i t}: t=\right. 1, \ldots, T\}$. Let the $1 \times t$ vector $\mathbf{g}^{t}$ be defined as in the previous paragraph: a sequence of zeros and ones indicating a treatment regime through time $t$. We are interested in the counterfactual outcomes $y_{i t}\left(\mathbf{g}^{t}\right)$, the outcome under regime $\mathbf{g}^{t}$. (Incidentally, unlike in some approaches to treatment effect estimation with panel data, we do not have to consider counterfactual outcomes as a function of future treatments. Therefore, we need not put restrictions on the counterfactuals such as the "no anticipation" requirement; see, for example, Abbring and Heckman (2007).) Naturally, we can write the observed outcome, $y_{i t}$, as a function of the $y_{i t}\left(\mathbf{g}^{t}\right)$ for $\mathbf{g}^{t} \in \mathscr{G}^{t}$, the set of all valid treatment regimes. Further, at each time $t$ we observe a set of covariates, $\mathbf{x}_{i}^{t}$ such that $\mathbf{x}_{i}^{t-1} \subset \mathbf{x}_{i}^{t}, t=2, \ldots, T$. At a minimum, $\mathbf{x}_{i}^{t}$ would typically include all past observed outcomes, $\left\{y_{i, t-1}, \ldots, y_{i 1}\right\}$, but it can also include time-constant variables, say, $\mathbf{z}_{i}$, collected in the intial period. In some cases, $\mathbf{x}_{i}^{t}$ can contain variables dated at time $t$, but one must be sure that such variables are not themselves affected by treatment.

We state the dynamic ignorability assumption (for each $t$ ) as


\begin{align*}
& \mathrm{D}\left[w_{i r} \mid y_{i t}\left(\mathbf{g}^{t}\right), w_{i, r-1}, \ldots, w_{i 1}, \mathbf{x}_{i}^{r}\right] \\
& \quad=\mathrm{D}\left[w_{i r} \mid w_{i, r-1}, \ldots, w_{i 1}, \mathbf{x}_{i}^{r}\right], \quad \mathbf{g}^{t} \in \mathscr{G}^{t}, r \leq t . \tag{21.138}
\end{align*}


When $r=t$, this condition is\\
$\mathbf{D}\left[w_{i t} \mid y_{i t}\left(\mathbf{g}^{t}\right), w_{i, t-1}, \ldots, w_{i 1}, \mathbf{x}_{i}^{t}\right]=\mathbf{D}\left[w_{i t} \mid w_{i, t-1}, \ldots, w_{i 1}, \mathbf{x}_{i}^{t}\right], \quad \mathbf{g}^{t} \in \mathscr{G}^{t}$,\\
which says that, conditional on past assignments and outcomes (contained in $\mathbf{x}_{i}^{t}$ ), assignment is independent of the counterfactual outcomes. Lechner (2004) refers to this as the weak dynamic conditional independence assumption. Equation (21.138) requires that ignorability of assignment with respect to the counterfactual at time $t$ hold in periods before $t$.

With a suitable overlap assumption, which will become clear shortly, we can show that equation (21.138) is sufficient to identify the means $\mathrm{E}\left[y_{i t}\left(\mathbf{g}^{t}\right)\right]$, which we can compare across different $\mathbf{g}^{t}$ to obtain average treatment effects. To simplify the notation, consider the mean $\mathrm{E}\left[y_{i t}\left(\mathbf{1}^{t}\right)\right]$ where $\mathbf{1}^{t}=(1,1, \ldots, 1)$, a $t$-vector of ones. This corresponds to treatment in every time period. Define


\begin{equation*}
p_{i r}\left(\mathbf{x}_{i}^{r}\right) \equiv \mathrm{P}\left(w_{i r}=1 \mid w_{i, r-1}=1, \ldots, w_{i 1}=1, \mathbf{x}_{i}^{r}\right), \quad r \leq t \tag{21.139}
\end{equation*}


(which is specific to the particular treatment sequence we are considering). By assumption (21.138), this is also the probability conditional on $y_{i t}\left(\mathbf{1}^{t}\right)$. Now, we want\\
to show that a particular kind of inverse probability weighting identifies $\mathrm{E}\left[y_{i t}\left(\mathbf{1}^{t}\right)\right]$. (This is similar to the attrition problem in Section 19.9.3 except that here we assume that the same units are observed in each time period. Thus, though there is a "missing data" problem in that we do not observe $y_{i t}\left(\mathbf{1}^{t}\right)$, we assume that, at a given time period $t$, for each unit we observe the entire history of past outcomes and other covariates.) The selected, weighted outcome $y_{i t}$ for $\mathbf{g}^{t}=(1, \ldots, 1)$ is\\
$\frac{w_{i t} w_{i, t-1} \cdots w_{i 1} y_{i t}}{p_{i t}\left(\mathbf{x}_{i}^{t}\right) \cdots p_{i 1}\left(\mathbf{x}_{i}^{1}\right)}=\frac{w_{i t} w_{i, t-1} \cdots w_{i 1} y_{i t}\left(\mathbf{1}^{t}\right)}{p_{i t}\left(\mathbf{x}_{i}^{t}\right) \cdots p_{i 1}\left(\mathbf{x}_{i}^{1}\right)}$,\\
where the equality follows because $w_{i t} w_{i, t-1} \cdots w_{i 1}=1$ implies that $y_{i t}=y_{i t}\left(\mathbf{1}^{t}\right)$. Because we have divided by the propensity scores in equation (21.140), we are, of course, assuming they are nonzero. This is the overlap assumption for this particular counterfactual outcome. If, for example, it is impossible to receive treatment in every period, then $\mathrm{E}\left[y_{i t}\left(\mathbf{1}^{t}\right)\right]$ is generally unidentified. Essentially, we must restrict attention to averages of counterfactual outcomes that correspond to possible treatment sequences.

We now show that the expected value of the second term is $\mathrm{E}\left[y_{i t}\left(\mathbf{1}^{t}\right)\right]$. We do so sequentially by iterated expectations and assumption (21.138). First,\\
$\mathrm{E}\left[\frac{w_{i t} w_{i, t-1} \cdots w_{i 1} y_{i t}\left(\mathbf{1}^{t}\right)}{p_{i t}\left(\mathbf{x}_{i}^{t}\right) \cdots p_{i 1}\left(\mathbf{x}_{i}^{1}\right)}\right]=\mathrm{E}\left\{\mathrm{E}\left[\left.\frac{w_{i t} w_{i, t-1} \cdots w_{i 1} y_{i t}\left(\mathbf{1}^{t}\right)}{p_{i t}\left(\mathbf{x}_{i}^{t}\right) \cdots p_{i 1}\left(\mathbf{x}_{i}^{1}\right)} \right\rvert\, y_{i t}\left(\mathbf{1}^{t}\right), w_{i, t-1}, \ldots, w_{i 1}, \mathbf{x}_{i}^{t}\right]\right\}$\\
Now, because $\mathbf{x}_{i}^{t-1} \subset \mathbf{x}_{i}^{t}$, the denominator is a function of the conditioning variables, and the conditional expectation of the numerator is

$$
\begin{aligned}
& \mathrm{E}\left[w_{i t} \mid y_{i t}\left(\mathbf{1}^{t}\right), w_{i, t-1}, \ldots, w_{i 1}, \mathbf{x}_{i}^{t}\right] w_{i t} w_{i, t-1} \cdots w_{i 1} y_{i t}\left(\mathbf{1}^{t}\right) \\
& \quad=\mathrm{P}\left(w_{i t}=1 \mid w_{i, t-1}, \ldots, w_{i 1}, \mathbf{x}_{i}^{t}\right) w_{i t} w_{i, t-1} \cdots w_{i 1} y_{i t}\left(\mathbf{1}^{t}\right) \\
& \quad=\mathrm{P}\left(w_{i t}=1 \mid w_{i, t-1}=1, \ldots, w_{i 1}=1, \mathbf{x}_{i}^{t}\right) w_{i t} w_{i, t-1} \cdots w_{i 1} y_{i t}\left(\mathbf{1}^{t}\right) \\
& \quad \equiv p_{i t}\left(\mathbf{x}_{i}^{t}\right) w_{i t} w_{i, t-1} \cdots w_{i 1} y_{i t}\left(\mathbf{1}^{t}\right)
\end{aligned}
$$

where the first equality follows from assumption (21.138) and the second follows because the term is zero unless all $w_{i s}, s=1, \ldots, t-1$, are unity. Therefore, we have shown that

$$
\mathrm{E}\left[\left.\frac{w_{i t} w_{i, t-1} \cdots w_{i 1} y_{i t}\left(\mathbf{1}^{t}\right)}{p_{i t}\left(\mathbf{x}_{i}^{t}\right) \cdots p_{i 1}\left(\mathbf{x}_{i}^{1}\right)} \right\rvert\, y_{i t}\left(\mathbf{1}^{t}\right), w_{i, t-1}, \ldots, w_{i 1}, \mathbf{x}_{i}^{t}\right]=\frac{w_{i, t-1} \cdots w_{i 1} y_{i t}\left(\mathbf{1}^{t}\right)}{p_{i, t-1}\left(\mathbf{x}_{i}^{t-1}\right) \cdots p_{i 1}\left(\mathbf{x}_{i}^{1}\right)},
$$

and so


\begin{equation*}
\mathrm{E}\left[\frac{w_{i t} w_{i, t-1} \cdots w_{i 1} y_{i t}\left(\mathbf{1}^{t}\right)}{p_{i t}\left(\mathbf{x}_{i}^{t}\right) \cdots p_{i 1}\left(\mathbf{x}_{i}^{1}\right)}\right]=\mathrm{E}\left[\frac{w_{i, t-1} \cdots w_{i 1} y_{i t}\left(\mathbf{1}^{t}\right)}{p_{i, t-1}\left(\mathbf{x}_{i}^{t-1}\right) \cdots p_{i 1}\left(\mathbf{x}_{i}^{1}\right)}\right] \tag{21.141}
\end{equation*}


Now we can repeat the argument conditioning on $\left[y_{i t}\left(\mathbf{1}^{t}\right), w_{i, t-2}, \ldots, w_{i 1}, \mathbf{x}_{i}^{t-1}\right]$ and use assumption (21.138) for $r=t-1$. And so on, until we get the expectation to\\
$\mathrm{E}\left[\frac{w_{i 1} y_{i t}\left(\mathbf{1}^{t}\right)}{p_{i 1}\left(\mathbf{x}_{i}^{1}\right)}\right]=\mathrm{E}\left[y_{i t}\left(\mathbf{1}^{t}\right)\right]$,\\
where one final application of assumption (21.138) and iterated expectations gets us the desired equality.

Given that\\
$\mathrm{E}\left[\frac{w_{i t} w_{i, t-1} \cdots w_{i 1} y_{i t}}{p_{i t}\left(\mathbf{x}_{i}^{t}\right) \cdots p_{i 1}\left(\mathbf{x}_{i}^{1}\right)}\right]=\mathrm{E}\left[y_{i t}\left(\mathbf{1}^{t}\right)\right]$,\\
where all quantities in the left-hand-side expectation are either observed or estimable, we can consistently estimate $\mathrm{E}\left[y_{i t}\left(\mathbf{1}^{t}\right)\right]$ by\\
$\hat{\mathrm{E}}\left[y_{i t}\left(\mathbf{1}^{t}\right)\right]=N^{-1} \sum_{i=1}^{N} \frac{w_{i t} w_{i, t-1} \cdots w_{i 1} y_{i t}}{\hat{p}_{i t}\left(\mathbf{x}_{i}^{t}\right) \cdots \hat{p}_{i 1}\left(\mathbf{x}_{i}^{1}\right)}$,\\
where $\hat{p}_{i r}\left(\mathbf{x}_{i}^{r}\right)$ is an estimate of $\mathrm{P}\left(w_{i r}=1 \mid w_{i, r-1}=1, \ldots, w_{i 1}=1, \mathbf{x}_{i}^{r}\right)$. We can estimate these probabilities very generally. With lots of data, we might estimate a flexible binary response model for $w_{i r}$ on the subsample with $w_{i, r=1}=1, \ldots, w_{i 1}=1$ with $\mathbf{x}_{i}^{r}$ as the covariates. Or, we might estimate a model for $\mathrm{P}\left(w_{i r}=1 \mid w_{i, r-1}, \ldots, w_{i 1}, \mathbf{x}_{i}^{r}\right)$ using all $i$, and then insert $w_{i, r=1}=1, \ldots, w_{i 1}=1$.

The approach for any sequence of potential treatments $\mathbf{g}^{t}$ should be clear. We use the treatment indicators to select out the appropriate subsample, and then estimate $\mathrm{P}\left(w_{i r}=g_{r}^{t} \mid w_{i, r-1}=g_{r-1}^{t}, \ldots, w_{i 1}=g_{1}^{t}, \mathbf{x}_{i}^{r}\right)$ for $r=1, \ldots, t$. For example, with $\mathbf{g}^{t}=(0,0, \ldots, 0)$ we select out the subsample using $\left(1-w_{i t}\right) \cdots\left(1-w_{i 1}\right)$ and then weight by the inverse of the product of the conditional probabilities. By estimating several combinations of treatment patterns we can obtain distributed lag effects of the policy intervention on current and future outcomes.

The two approaches just described-ignorability conditional on unobserved heterogeneity and dynamic ignorability-are generally consistent under different assumptions. Nevertheless, it is sometimes (at least implicitly) claimed that dynamic ignorability is less restrictive. To see that this claim can seem to be true but is not, it is useful to work through a simple case.

Suppose that $T=2$ and there is no treatment in the first time period, so $w_{i 1} \equiv 0$. Consequently, the observed outcome in period one is $y_{i 1}$, and there is no need\\
to distinguish between it and counterfactual outcomes. In period 2 the counterfactual outcomes are $y_{i 2}(0)$ and $y_{i 2}(1)$, and the observed outcome is $y_{i 2}=y_{i 2}(0)+ w_{i 2}\left[y_{i 2}(1)-y_{i 2}(0)\right]$. Assume that the treatment effect is the constant $\tau$. If we assume ignorability conditional on heterogeneity $c_{i}$, then we can write $y_{i 1}=c_{i}+u_{i 1}, y_{i 2}= c_{i}+\eta+\tau w_{i 2}+u_{i 2}$, where $\eta$ is the difference in intercepts between the two time periods (assumed to be constant). These simple representations lead to a simple differenced equation:\\
$\Delta y_{i 2}=\eta+\tau w_{i 2}+\Delta u_{i 2}$,\\
where $\Delta y_{i 2}=y_{i 2}-y_{i 1}$ and we use the fact that $w_{i 1}=0$. If we estimate equation (21.144) by OLS, the OLS estimate is the first-difference (FD) estimate,\\
$\hat{\tau}_{F D}=\Delta \bar{y}_{\text {treat }}-\Delta \bar{y}_{\text {control }}$,\\
which is the difference in average changes over time between the treatment and control groups. This estimator is consistent because treatment is strictly exogenous conditional on the heterogeneity.

A commonly used alternative in the statistics literature is to add the first-period outcome as a control, that is, use the regression\\
$\Delta y_{i 2}$ on $1, w_{i 2}, y_{i 1}$.

If we add the assumption that the shocks $\left\{u_{i 1}, u_{i 2}\right\}$ are serially uncorrelatedspecifically, $\mathrm{E}\left(u_{i 1} u_{i 2} \mid w_{i 1}, c_{i}\right)=0$-then adding $y_{i 1}$ overcontrols and leads to inconsistency. In fact, it can be shown (for example, Angrist and Pischke (2009), Section 5.4) that\\
$\operatorname{plim}\left(\hat{\tau}_{L D V}\right)=\tau+\pi_{1}\left(\sigma_{u_{1}}^{2} / \sigma_{r_{2}}^{2}\right)$,\\
where $w_{i 2}=\pi_{0}+\pi_{1} y_{i 1}+r_{i 2}$ is a linear projection, and so $\pi_{1}=\operatorname{Cov}\left(c_{i}, w_{i 2}\right) / \left(\sigma_{c}^{2}+\sigma_{u_{1}}^{2}\right)$. We can easily sign the bias given the sign of $\pi_{1}$. For example, if $w_{i 2}$ indicates a job training program and less productive workers are more likely to participate ( $\pi_{1}<0$ ), then the regression that controls for $y_{i 1}$ underestimates the job training effect. If more productive workers participate, it overestimates the effect of job training. This simple example illustrates a general lesson in estimating treatment effects, whether one uses regression, propensity score weighting, or matching: it is not necessarily true that controlling for more covariates results in less bias. This example shows that adding covariates can actually induce bias when there was none.

Now suppose that ignorability of treatment holds conditional on $y_{i 1}$ (and the treatment effect is constant). Then we can write\\
$\Delta y_{i 2}=\gamma+\tau w_{i 2}+\lambda y_{i 1}+e_{i 2}$\\
$\mathrm{E}\left(e_{i 2}\right)=0, \quad \operatorname{Cov}\left(w_{i 2}, e_{i 2}\right)=\operatorname{Cov}\left(y_{i 1}, e_{i 2}\right)=0$.\\
Now, of course, controlling for $y_{i 1}$ consistently estimates $\tau$ because we are assuming that controlling for $y_{i 1}$ (in a linear way) leads to ignorable treatment. On the other hand, the FD estimator suffers from omitted variable "bias":\\
$\operatorname{plim}\left(\hat{\tau}_{F D}\right)=\tau+\lambda \frac{\operatorname{Cov}\left(w_{i 2}, y_{i 1}\right)}{\operatorname{Var}\left(w_{i 2}\right)}$.

Suppose there is regression to the mean, so that $\lambda<0$. If workers observed with low first-period earnings are more likely to participate, so that $\operatorname{Cov}\left(w_{i 2}, y_{i 1}\right)<0$, then $\operatorname{plim}\left(\hat{\tau}_{F D}\right)>\tau$, and so FD overestimates the effect. In fact, if the correlation between $w_{i 2}$ and $y_{i 1}$ is negative in the sample, $\hat{\tau}_{F D}>\hat{\tau}_{L D V}$ is an algebraic fact. But this does not allow us to determine which estimator has less inconsistency because the derivations rely on different ignorability assumptions.

\section*{Problems}
21.1. Consider the difference-in-means estimator, $\bar{d}=\bar{y}_{1}-\bar{y}_{0}$, where $\bar{y}_{g}$ is the sample average of the $y_{i}$ with $w_{i}=g, g=0,1$.\\
a. Show that, as an estimator of $\tau_{a t t}$, the bias in $\bar{y}_{1}-\bar{y}_{0}$ is $\mathrm{E}\left(y_{0} \mid w=1\right)- \mathrm{E}\left(y_{0} \mid w=0\right)$.\\
b. Let $y_{0}$ be the earnings someone would earn in the absence of job training, and let $w=1$ denote the job training indicator. Explain the meaning of $\mathrm{E}\left(y_{0} \mid w=1\right)< \mathrm{E}\left(y_{0} \mid w=0\right)$. Intuitively, does it make sense that $\mathrm{E}(\bar{d})<\tau_{a t t}$ ?\\
21.2. Explain how you would estimate $\tau_{\text {ate }, \mathscr{R}}=\mathrm{E}\left(y_{1}-y_{0} \mid \mathbf{x} \in \mathscr{R}\right)$ using propensity score weighting under Assumption ATE.1'.\\
21.3. Use the data in JTRAIN3.RAW to answer these questions. The response in this case is the binary variable unem 78.\\
a. Estimate $\tau_{\text {ate }}$ from the simple regression of unem $78_{i}$ on 1, train $_{i}$. Does the training program have the anticipated effect on being unemployed?\\
b. Add as controls the variables age, educ, black, hisp, married, re74, re75, unem75, and unem74. How does the estimate of $\tau_{\text {ate }}$ compare with the simple regression estimate from part a? What is the (heteroskedasticity-robust) $95 \%$ confidence interval?\\
c. Using the same controls as in part b , estimate separate regressions for the treatment and control groups. Now what is $\hat{\tau}_{\text {ate }}$ ? Does it differ much from the estimate in part b ? What is $\hat{\tau}_{\text {att }}$ ?\\
d. Now use regression adjustment with controls age, educ, black, hisp, married, re74, and re75, but use only the subsample of men who were unemployed in 1974 or 1975. Again use two separate regression functions. What are $\hat{\tau}_{\text {ate }}$ and $\hat{\tau}_{\text {att }}$ ? How do these compare with the estimates from part c?\\
e. Estimate the propensity score using logit and the explanatory variables in part b. How many outcomes are perfectly predicted? What does this result mean for IPW estimation?\\
f. Now estimate the propensity score using only observations with avgre $\leq 15$. Use the $\hat{p}\left(\mathbf{x}_{i}\right)$ in IPW estimation to obtain $\hat{\tau}_{\text {ate }}$ and $\hat{\tau}_{\text {att }}$. Be sure to only use observations with avgre $\leq 15$. Compare these with the corresponding regression adjustment estimates using the avgre $\leq 15$ observations.\\
21.4. Carefully derive equation (21.80).\\
21.5. Use the data in JTRAIN2.RAW for this question.\\
a. As in Example 21.2, run a probit of train on $1, \mathbf{x}$, where $\mathbf{x}$ contains the covariates from Example 21.2. Obtain the probit fitted values, say $\hat{\Phi}_{i}$.\\
b. Estimate the equation $r e 78_{i}=\gamma_{0}+\tau \operatorname{train}_{i}+\mathbf{x}_{i} \gamma+u_{i}$ by IV, using instruments $\left(1, \hat{\Phi}_{i}, \mathbf{x}_{i}\right)$. Comment on the estimate of $\tau$ and its standard error.\\
c. Regress $\hat{\Phi}_{i}$ on $\mathbf{x}_{i}$ to obtain the $R$-squared. What do you make of this result?\\
d. Does the nonlinearity of the probit model for train allow us to estimate $\tau$ when we do not have an additional instrument? Explain.\\
21.6. In Procedure 21.2, explain why it is better to estimate equation (21.71) by IV rather than to run the OLS regression $y_{i}$ on $1, \hat{G}_{i}, \mathbf{x}_{i}, \hat{G}_{i}\left(\mathbf{x}_{i}-\overline{\mathbf{x}}\right), i=1, \ldots, N$.\\
21.7. As a special case of the setup in Section 21.6.3 we can write $y_{i}=a_{i}+b_{i} w_{i}$ where $w_{i}$ is a scalar. If we define $\beta=\mathrm{E}\left(b_{i}\right)$, then, under assumptions (21.124) and (21.125),\\
$\beta=\mathrm{E}\left\{\frac{\left[w_{i}-\psi\left(\mathbf{x}_{i}\right)\right] y_{i}}{\omega\left(\mathbf{x}_{i}\right)}\right\}$,\\
where $\psi\left(\mathbf{x}_{i}\right)=\mathrm{E}\left(w_{i} \mid \mathbf{x}_{i}\right)$ and $\omega\left(\mathbf{x}_{i}\right)=\operatorname{Var}\left(w_{i} \mid \mathbf{x}_{i}\right)$.\\
a. Suppose that the treatment variable $w_{i}$ is a nonnegative count variable (such as visits to one's family physician during a year) and you think $\mathrm{E}(w \mid \mathbf{x})=\exp \left(\gamma_{0}+\mathbf{x} \boldsymbol{\gamma}\right)$\\
and $\operatorname{Var}(w \mid \mathbf{x})=\exp \left(\delta_{0}+\mathbf{x} \boldsymbol{\delta}\right)$. Propose a method for estimating the mean and variance parameters. [Hint: For the latter, define appropriate errors, say $r=w-\mathrm{E}(w \mid \mathbf{x})$, and note that $\mathrm{E}\left(r^{2} \mid \mathbf{x}\right)=\operatorname{Var}(w \mid \mathbf{x})$.]\\
b. Given the estimators from part a, suggest a consistent, $\sqrt{N}$-asymptotically normal estimator of $\beta$. How would you conduct inference on $\beta$ ?\\
c. Apply the method from parts a and b to the data in JTRAIN2.RAW, using as $w$ the variable mostrn (months spent in job training) and $y=r e 78$. Use the same variables in $\mathbf{x}$ as Example 21.1.\\
d. Now suppose $w_{i}$ is a proportion (such as proportion of lectures attended) and you think $\mathrm{E}(w \mid \mathbf{x})=\Lambda\left(\gamma_{0}+\mathbf{x} \boldsymbol{\gamma}\right)$ and $\operatorname{Var}(w \mid \mathbf{x})=\delta_{0}+\delta_{1} \mathrm{E}(w \mid \mathbf{x})+\delta_{2}[\mathrm{E}(w \mid \mathbf{x})]^{2}$, where $\Lambda(\cdot)$ is the logistic function. Now how would you estimate the mean and variance parameters? What practical problem might arise with the estimated variances?\\
e. Apply the approach from part d to the data set ATTEND.RAW with $w=$ atndrte, the proportion (not percentage) of lectures attended. For the response use $y=$ stndful, for the elements of $\mathbf{x}$ use cubics in priGPA and $A C T$, and use the frosh and soph binary indicators. What is the estimate of $\beta$ ? How does it compare to the multiple regression estimate of stndfnl on atndrte and the given controls?\\
21.8. In the IV setup of Section 21.6.3, suppose that $b=\beta$, and therefore we can write\\
$y=a+\beta w+e, \quad \mathrm{E}(e \mid a, \mathbf{x}, \mathbf{z})=0$.\\
Assume that conditions (21.129) and (21.130) hold for $a$.\\
a. Suppose $w$ is a corner solution outcome, such as hours spent in a job training program. If $\mathbf{z}$ is used as IVs for $w$ in $y=\gamma_{0}+\beta w+\mathbf{x} \gamma+r$, what is the identification condition?\\
b. If $w$ given $(\mathbf{x}, \mathbf{z})$ follows a standard Tobit model, propose an IV estimator that uses the Tobit fitted values for $w$.\\
c. If $\operatorname{Var}(e \mid a, \mathbf{x}, \mathbf{z})=\sigma_{e}^{2}$ and $\operatorname{Var}(a \mid \mathbf{x}, \mathbf{z})=\sigma_{a}^{2}$, argue that the IV estimator from part b is asymptotically efficient.\\
d. What is an alternative to IV estimation that would use the Tobit fitted values for $w$ ? Which method do you prefer?\\
e. If $b \neq \beta$, but assumptions (21.129) and (21.130) hold, how would you estimate $\beta$ ?\\
21.9. a. Using the data in JTRAIN3.RAW, estimate the logit model for the propensity score in Example 21.1, using the same explanatory variables described there.

Plot histograms for the $\operatorname{train}_{i}=1$ and $\operatorname{train}_{i}=0$ sample separately. What do you conclude about the overlap assumption?\\
b. Now use the subsample with $(r e 74+r e 75) / 2 \leq 15$, and repeat the exercise from part a. How does the situation compare with the full data set?\\
21.10. In this problem you are to show that one can improve efficiency in estimating a constant ATE under random assignment by including regressors (provided those regressors are also independent of assignment). Let $w_{i}$ be the treatment, and assume that $w_{i}$ is independent of $\left(y_{i 0}, y_{i 1}, \mathbf{x}_{i}\right)$. Assume that $y_{i 1}-y_{i 0}=\tau$ for all $i$. Thus, write\\
$y_{i}=y_{i 0}+\tau w_{i} \equiv \mu_{0}+\tau w_{i}+v_{i 0}$.\\
a. Derive the asymptotic variance of the simple OLS estimator $\tilde{\tau}$ from the regression of $y_{i}$ on $1, w_{i}$ in terms of $\operatorname{Var}\left(v_{i 0}\right), \rho=\mathrm{P}\left(w_{i}=1\right)$, and $N$.\\
b. Write the linear projection of $y_{i 0}$ on $\left(1, \mathbf{x}_{i}\right)$ as\\
$y_{i 0}=\alpha_{0}+\mathbf{x}_{i} \boldsymbol{\beta}_{0}+u_{i 0}$\\
$\mathrm{E}\left(u_{i 0}\right)=0, \quad \mathrm{E}\left(\mathbf{x}_{i}^{\prime} u_{i 0}\right)=\mathbf{0}$.\\
Show that we can write\\
$y_{i}=\alpha_{0}+\tau w_{i}+\mathbf{x}_{i} \boldsymbol{\beta}_{0}+u_{i 0}$,\\
and explain why $w_{i}$ is uncorrelated with $\mathbf{x}_{i}$ and $u_{i 0}$.\\
c. Using part b , show that the asymptotic variance of $\hat{\tau}$ from the regression $y_{i}$ on 1 , $w_{i}, \mathbf{x}_{i}$ is $\operatorname{Var}\left(u_{i 0}\right) /[N \rho(1-\rho)]$.\\
d. Show that if $\boldsymbol{\beta}_{0} \neq \mathbf{0}, \operatorname{Var}\left(u_{i 0}\right)<\operatorname{Var}\left(v_{i 0}\right)$. Conclude that the asymptotic variance of $\hat{\tau}$ is strictly smaller than that of $\tilde{\tau}$ if $\mathbf{x}_{i}$ helps to predict $y_{i 0}$.\\
e. Suppose that $\mathrm{E}\left(y_{i 0} \mid \mathbf{x}\right) \neq \alpha_{0}+\mathbf{x}_{i} \boldsymbol{\beta}_{0}$-that is, the linear projection differs from the conditional mean. In this situation, why might you use the simple difference-in-means estimator under random assignment. (Hint: Think about small-sample versus largesample properties.)\\
21.11. Suppose that we allow full slope, as well as intercept, heterogeneity in a linear representation of two counterfactual outcomes,\\
$y_{i 0}=a_{i 0}+\mathbf{x}_{i} \mathbf{b}_{i 0}$\\
$y_{i 1}=a_{i 1}+\mathbf{x}_{i} \mathbf{b}_{i 1}$.

Assume that the vector $\left(\mathbf{x}_{i}, \mathbf{z}_{i}\right)$ is independent of $\left(a_{i 0}, \mathbf{b}_{i 0}, a_{i 1}, \mathbf{b}_{i 1}\right)$-which makes, as we will see, $\mathbf{z}_{i}$ instrumental variables candidates in a control function or correction function setting.\\
a. Define $\alpha_{g} \equiv \mathrm{E}\left(a_{i g}\right)$ and $\boldsymbol{\beta}_{g} \equiv \mathrm{E}\left(\mathbf{b}_{i g}\right), g=0,1$ and $\boldsymbol{\psi} \equiv \mathrm{E}\left(\mathbf{x}_{i}\right)$. Find $\mu_{g} \equiv \mathrm{E}\left(y_{i g}\right)$, $g=0,1$ in terms of these population parameters.\\
b. Let $\tau=\tau_{\text {ate }}=\mu_{1}-\mu_{0}$ be the ATE. Also, write $a_{i g}=\alpha_{g}+c_{i g}$ and $\mathbf{b}_{i g}=\boldsymbol{\beta}_{g}+\mathbf{f}_{i g}$. Show that\\
$y_{i}=\mu_{0}+\tau w_{i}+\left(\mathbf{x}_{i}-\boldsymbol{\psi}\right) \boldsymbol{\beta}_{0}+w_{i}\left(\mathbf{x}_{i}-\boldsymbol{\psi}\right) \boldsymbol{\delta}+c_{i 0}+\mathbf{x}_{i} \mathbf{f}_{i 0}+w_{i} e_{i}+w_{i} \mathbf{x}_{i} \mathbf{d}_{i}$,\\
where $\boldsymbol{\delta} \equiv \boldsymbol{\beta}_{1}-\boldsymbol{\beta}_{0}, e_{i} \equiv c_{i 1}-c_{i 0}$, and $\mathbf{d}_{i} \equiv \mathbf{f}_{i 1}-\mathbf{f}_{i 0}$.\\
c. Assume that\\
$w_{i}=1\left[\theta_{0}+\mathbf{x}_{i} \boldsymbol{\theta}_{1}+\mathbf{z}_{i} \boldsymbol{\theta}_{2}+v_{i} \geq 0\right] \equiv 1\left[\mathbf{q}_{i} \boldsymbol{\theta}+v_{i} \geq 0\right]$,\\
and assume that $\mathrm{E}\left(c_{i 0} \mid v_{i}, \mathbf{x}_{i}, \mathbf{z}_{i}\right)=\rho_{0} v_{i}, \mathrm{E}\left(\mathbf{f}_{i 0} \mid v_{i}, \mathbf{x}_{i}, \mathbf{z}_{i}\right)=\boldsymbol{\eta}_{0} v_{i}, \mathrm{E}\left(e_{i} \mid v_{i}, \mathbf{x}_{i}, \mathbf{z}_{i}\right)=\xi v_{i}$, and $\mathrm{E}\left(\mathbf{d}_{i} \mid v_{i}, \mathbf{x}_{i}, \mathbf{z}_{i}\right)=\boldsymbol{\theta} a_{i}$. Find $\mathrm{E}\left(y_{i} \mid v_{i}, \mathbf{x}_{i}, \mathbf{z}_{i}\right)$.\\
d. Assuming that $v_{i}$ is independent of $\left(\mathbf{x}_{i}, \mathbf{z}_{i}\right)$ and has a standard normal distribution, use part c to find $\mathrm{E}\left(y_{i} \mid w_{i}, \mathbf{x}_{i}, \mathbf{z}_{i}\right)$. (Hint: This should depend on the generalized residual $h\left(w_{i}, \mathbf{q}_{i} \boldsymbol{\theta}\right)=w_{i} \lambda\left(\mathbf{q}_{i} \boldsymbol{\theta}\right)-\left(1-w_{i}\right) \lambda\left(-\mathbf{q}_{i} \boldsymbol{\theta}\right)$, where $\lambda(\cdot)$ is the inverse Mills ratio.)\\
e. Propose a two-step control function approach to estimating $\tau$. How does this differ from the method derived in Section 21.4.2?\\
f. How would you obtain a valid standard error for $\hat{\tau}$ from part e?\\
g. How would you estimate $\tau_{\text {ate }}(\mathbf{x})=\mathrm{E}\left(y_{1}-y_{0} \mid \mathbf{x}\right)$ ?\\
21.12. Consider the same setup as Problem 21.11, under the same assumptions.\\
a. Of the four unobserved terms in equation (21.149), which two have a zero mean conditional on $\left(\mathbf{x}_{i}, \mathbf{z}_{i}\right)$ ?\\
b. Derive the correction functions for the other terms.\\
c. Propose a two-step correction function estimator of $\tau$ (and the other parameters).\\
d. How would you test whether the correction functions are needed? Be very precise.\\
21.13. For the sharp regression discontinuity design, what kind of local estimation method might you use if $y_{i}$ is a fractional response? Specifically describe the model and estimation method, and provide the formula for $\hat{\tau}_{c}$.\\
21.14. Consider a treatment effect setup with a binary treatment, $w_{i t}$, in each time period. We are interested in contemporaneous ATEs, $\tau_{t, \text { ate }}=\mathrm{E}\left[y_{i t}(1)-y_{i t}(0)\right] \equiv \mu_{t 1}-\mu_{t 0}$. Assume that for covariates $\mathbf{x}_{i t}, y_{i t}(g)=a_{i t g}+\mathbf{x}_{i t} \boldsymbol{\beta}_{g}, g=0,1$, where $a_{i t g}$ are\\
random variables and $\boldsymbol{\beta}_{g}$ is assumed (for simplicity) to be constant over time. Write $a_{i t g}=a_{t g}+c_{i t g}$, where $\alpha_{t g} \equiv \mathrm{E}\left(a_{i t g}\right)$. Define $\boldsymbol{\psi}_{t} \equiv \mathrm{E}\left(\mathbf{x}_{i t}\right)$.\\
a. Find $\mu_{t 0}, \mu_{t 1}$, and $\tau_{t, \text { ate }}$ in terms of the $\alpha_{t g}, \boldsymbol{\beta}_{g}$, and $\boldsymbol{\psi}_{t}$.\\
b. Let $y_{i t}=\left(1-w_{i t}\right) y_{i t}(0)+w_{i t} y_{i t}(1)$ and let $\tau_{t}=\tau_{t, a t e}$ for notational simplicity. Show that we can write\\
$y_{i t}=\mu_{t 0}+\tau_{t} w_{i t}+\left(\mathbf{x}_{i t}-\boldsymbol{\psi}_{t}\right) \boldsymbol{\beta}_{0}+w_{i t}\left(\mathbf{x}_{i t}-\boldsymbol{\psi}_{t}\right) \boldsymbol{\delta}+c_{i t 0}+w_{i t} e_{i t}$,\\
where $\boldsymbol{\delta} \equiv \boldsymbol{\beta}_{1}-\boldsymbol{\beta}_{0}$ and $e_{i t} \equiv c_{i t 1}-c_{i t 0}$.\\
c. We think in equation (21.150) that $w_{i t}$ is correlated with $c_{i t 0}$ and $e_{i t}$, possibly because these terms each contain time-constant heterogeneity and idiosyncratic unobservables that are related to treatment. Suppose we have a vector $\mathbf{z}_{i t}$ of potential instrumental variables. We allow the time average of the instruments and covariates to be correlated with $\left(c_{i t 0}, e_{i t}\right)$. In particular,\\
$c_{i t 0}=\left(\overline{\mathbf{x}}_{i}-\boldsymbol{\mu}_{\overline{\mathbf{x}}}\right) \boldsymbol{\xi}_{1}+\left(\overline{\mathbf{z}}_{i}-\boldsymbol{\mu}_{\overline{\mathbf{z}}}\right) \boldsymbol{\xi}_{2}+r_{i t 0}, \quad \mathrm{E}\left(r_{i t 0} \mid \mathbf{x}_{i}, \mathbf{z}_{i}\right)=0$\\
$e_{i t}=\left(\overline{\mathbf{x}}_{i}-\boldsymbol{\mu}_{\overline{\mathbf{x}}}\right) \boldsymbol{\eta}_{1}+\left(\overline{\mathbf{z}}_{i}-\boldsymbol{\mu}_{\overline{\mathbf{z}}}\right) \boldsymbol{\eta}_{2}+v_{i t}, \quad \mathrm{E}\left(v_{i t} \mid \mathbf{x}_{i}, \mathbf{z}_{i}\right)=0$,\\
where $\mathbf{x}_{i}$ and $\mathbf{z}_{i}$ are the entire histories of the covariates and instruments. (Removing the population means of $\overline{\mathbf{x}}_{i}$ and $\overline{\mathbf{z}}_{i}$ ensures $c_{i t 0}$ and $e_{i t}$ have zero means.) Rewrite equation (21.150) so that it includes the time averages.\\
d. Assume that $w_{i t}$ can be expressed as\\
$w_{i t}=1\left[\theta_{0}+\mathbf{x}_{i t} \boldsymbol{\theta}_{1}+\mathbf{z}_{i t} \boldsymbol{\theta}_{2}+\overline{\mathbf{x}}_{i} \boldsymbol{\theta}_{3}+\overline{\mathbf{z}}_{i} \boldsymbol{\theta}_{4}+q_{i t} \geq 0\right]$\\
$\mathrm{D}\left(q_{i t} \mid \mathbf{x}_{i}, \mathbf{z}_{i}\right)=\operatorname{Normal}(0,1)$.\\
Further, assume that $\mathrm{E}\left(c_{i t 0} \mid q_{i t}, \mathbf{x}_{i}, \mathbf{z}_{i}\right)=\alpha_{0} q_{i t}$ and $\mathrm{E}\left(e_{i t} \mid q_{i t}, \mathbf{x}_{i}, \mathbf{z}_{i}\right)=\rho q_{i t}$. Find $\mathrm{E}\left(y_{i t} \mid w_{i t}, \mathbf{x}_{i}, \mathbf{z}_{i}\right)$.\\
e. Use the expression from part d to propose a two-step control function approach to estimating the $\tau_{t}$.\\
21.15. Use the data in CATHETER.RAW to answer this question. The treatment variable is rhc, which is unity if a patient received a right-heart catheterization and zero if not. The response variable is death, equal to unity if the patient died within 180 days. The patients were all admitted to the intensive care unit, and so the death rate is very high: almost $65 \%$ of the patients died in the first six months. See Li , Racine, and Wooldridge (2008).\\
a. How many observations are in the sample? What fraction of patients received the RHC treatment?\\
b. Use the difference in means to estimate the average treatment effect of rhc. Does the treatment reduce the probability of death? Is the estimate practically large and statistically significant?\\
c. Now use regression adjustment, estimating logit models for death separately for $r h c=0$ and $r h c=1$. As covariates, use sex, race, income, cat1, cat2, ninsclas, and age. For all but the age variable include dummy variables for all categories (except, of course, a base category). What is the regression adjustment estimate of $\tau_{\text {ate }}$ ? Of $\tau_{\text {att }}$ ? Obtain standard errors using, say, 1,000 bootstrap replications.\\
d. Estimate a logit model for rhc using the same covariates as in part c. Compare the range and average of the estimated propensity scores for the treated and untreated samples. Plot a histogram in each case, and comment on overlap.\\
e. Estimate $\tau_{\text {ate }}$ and $\tau_{\text {att }}$ by propensity score weighting. Again use 1,000 bootstrap replications to obtain standard errors. How do the estimates compare with those from part d? Overall, does it appear RHC reduces the probability of death?\\
21.16. Use the data in REGDISC.RAW to answer this question. These are simulated data of a fuzzy regression discontinuity design with forcing variable $x$. The discontinuity is at $x=5$.\\
a. What fraction of the observations have $x_{i} \geq 5$ ? What fraction of units are in the treated group (that is, have $w_{i}=1$ )?\\
b. Estimate separate linear probability models for $x_{i}<5$ and $x_{i} \geq 5$, and obtain the fitted line. Do the same for logit models, and obtain the fitted logit functions. Graph these functions on the same graph. What is the estimated jump in the probability of treatment in each case?\\
c. Compute the estimate of $\tau_{c}$ with $c=5$ in equation (21.107) using linear regression for $y$ and $w$, and using all the data.\\
d. Now use the IV estimate described in equation (21.108), again using all the data. Is it the same as the estimate in part c? What is its standard error?\\
e. Now use a "local" version of the IV method, restricting attention to data with $x_{i}>3$ and $x_{i}<7$. What happens to the estimated $\tau_{c}$ and its standard error compared with that in part d ?

\section*{22 Duration Analysis}

\end{document}
